% Raskin《上同调类域论》中文试译 — 主文件（本目录下 `xelatex main` 三遍即可出 PDF）
\documentclass[11pt,a4paper]{article}
\usepackage[margin=2.6cm]{geometry}
\usepackage{xeCJK}
\setCJKmainfont[AutoFakeBold=2.5,AutoFakeSlant=0.2]{SimSun}
\setCJKsansfont{Microsoft YaHei}
\setCJKmonofont{SimSun}
\usepackage{amsmath,amssymb,amsthm,amscd,mathrsfs}
\usepackage{array,booktabs,longtable}
\usepackage{enumitem}
\usepackage{xcolor}
\usepackage[colorlinks=true,linkcolor=blue!50!black,urlcolor=blue!50!black,bookmarks=false]{hyperref}
\setlength{\parindent}{2em}
\setlength{\parskip}{0.35em}
\linespread{1.18}
\allowdisplaybreaks

\newcommand{\Z}{\mathbb{Z}}
\newcommand{\Q}{\mathbb{Q}}
\newcommand{\F}{\mathbb{F}}
\newcommand{\C}{\mathbb{C}}
\newcommand{\R}{\mathbb{R}}
\newcommand{\OK}{\mathcal{O}}
\newcommand{\Gal}{\operatorname{Gal}}
\newcommand{\Br}{\operatorname{Br}}
\newcommand{\Hh}{\mathrm{H}}
\newcommand{\HH}{\widehat{H}}
\newcommand{\coker}{\operatorname{coker}}
\newcommand{\im}{\operatorname{im}}
\newcommand{\Spec}{\operatorname{Spec}}
\newcommand{\Gm}{\mathbb{G}_m}
\providecommand{\check}[1]{\v{#1}}

\title{\bfseries 上同调类域论讲义\\[2pt]
\large 麻省理工学院 18.786（2016 春）\quad 主讲：Sam Raskin\quad 整理：Oron Propp}
\author{中文试译（自动流程 + 人工合稿）}
\date{底本：2017-08-21 版；试译：2026-11-25；校：2026-11-26}

\providecommand{\Hom}{\operatorname{Hom}}
\providecommand{\Aut}{\operatorname{Aut}}
\providecommand{\End}{\operatorname{End}}
\providecommand{\Spec}{\operatorname{Spec}}
\providecommand{\Ext}{\operatorname{Ext}}
\providecommand{\Tor}{\operatorname{Tor}}
\providecommand{\Coker}{\operatorname{Coker}}
\providecommand{\Frob}{\operatorname{Frob}}

\begin{document}
\maketitle
\tableofcontents
\newpage

\section*{译者说明}

这份 PDF 是 Sam Raskin 教授 2016 年春季在 MIT 讲授的《数论 II（18.786）》讲义的中文试译。底本为 Oron Propp 整理的英文讲义（2017 年 8 月 21 日版，共 23 讲，119 页）。译文由自动翻译流程生成、经人工合稿校对，用于检验「数学长文 + 检索插件」这条流水线；不保证与英文原稿逐字对应，引用请以英文底本为准。

几点约定：
\begin{itemize}[leftmargin=2em]
  \item 原讲义中的定义、定理、例、式均保留原有编号（如「定理 1.5」「式 (1.1)」），不重新编号；讲标题按原文顺序排列。
  \item 原稿是从 PDF 文本层提取的，公式在提取时普遍出现上下标错位、符号粘连。译文中的公式均按数学含义重排为 \LaTeX{}，并已逐条对照公开的原讲义 PDF（MIT OCW 版）核校；原文本身的排印错误，凡能确证者已在译文中订正，并在 \texttt{\% [原文疑误：…]} 注释中留痕（注释不进入 PDF），凡仅属抽取丢字、原文可通者则照原文排并加注说明。
  \item 术语以随书术语表为准（见文末）；人名首次出现时附英文原文，术语中的姓氏保留拉丁拼写（如 Artin 互反律、Tate 上同调）。
  \item 参考文献条目保留英文原貌，以便读者按原书索引查证。
\end{itemize}

\section*{前言}

这份讲义对应 2016 年春季学期由 Sam Raskin 在 MIT 讲授的《数论 II》（18.786）课程。课程原网页见 \url{http://math.mit.edu/~sraskin/cft/index.html}；除本讲义外，网页上还有带注释的参考书目与习题集，讲义中会经常引用。习题集是课程的重要部分，其中常常引入讲义正文没有覆盖的材料（不过索引里大多列了出来）。

我写这份讲义时，总体上尽量复现 Sam 讲课的内容，偶尔补全细节、或为便于自己理解而稍作重组。手写笔记在学期中就开始录入，约一半工作是在大二暑假有更多空闲时完成的。因此，讲义有些地方比较粗糙，但我猜越往后质量越好。当然，其中的错误都只属于我，与 Sam 无关；如有笔误、错误或其它意见，欢迎发邮件至 opropp@mit.edu。

课程从第 2–5 讲开始，用 Hilbert 符号分析类域论的二次情形，以便对类域论的思想作更具体的引入。第 6–13 讲发展此后处理类域论所需的上同调工具。接着，第 14–18 讲证明局部类域论的主定理；第 16 讲由 Marc Hoyois 在 Sam 缺席时代讲，谨此致谢。最后，第 19–23 讲证明整体类域论的主定理，这大概是全课最难的部分。课程以 Bjorn Poonen 教授的两讲收尾：第一讲关于二次型与 Hasse--Minkowski 定理，第二讲关于 Witt 消去与 Brauer--Manin 障碍；遗憾的是，我还没来得及把它们整理成文。

感谢 Sam 讲授这门课、耐心回答我关于材料的许多问题、鼓励我完成并上传这份讲义，以及对我早先草稿的评论。请注意，第 1 讲与第 16 讲很大程度上改编自 Sam 本人提供的讲义。

希望这份讲义对你有用，也希望你像我一样享受这门课。


% ==== LECTURES ====
\addcontentsline{toc}{section}{第 1 讲　引言}
% 第 1 讲 引言
% 源：src/lecture01.txt（Raskin, 18.786 讲义，2017-08-21 版）
% 说明：原文页眉「1. INTRODUCTION 3/5」为 PDF 页眉残留，已删。

\section*{第 1 讲\quad 引言}

本课程先研究类域论（class field theory，cft）的二次情形，重点在显式 cft。随后我们将发展一套处理 cft 的上同调方法。最后，若时间允许，我们还会讨论若干附加课题，例如（更深入的）显式 cft、用 Fontaine--Herr 方法处理局部类域论（local class field theory，lcft）、代数群，以及 Tate 对偶。

类域论在十九世纪至少从三条研究线索中浮现出来。第一条是根式可解性问题：$\mathbb{Q}$ 中哪些代数数可以用 $n$ 次根、和等等表示出来？Abel 与 Galois 证明了：对某个数域 $K$，不可约多项式 $f(x)\in K[x]$ 的根可用根式表示，当且仅当 $f$ 的分裂域的 Galois 群可解，也就是说，$f$ 的分裂域是一串阿贝尔扩张的迭代扩张，例如
\[
  \mathbb{Q}\overset{\mathbb{Z}/2\mathbb{Z}}{\subseteq}\mathbb{Q}(\zeta_3)
  \overset{\mathbb{Z}/3\mathbb{Z}}{\subseteq}\mathbb{Q}(\zeta_3,\sqrt[3]{2}),
\]
其中我们把每个子扩张的 Galois 群写在相应的包含关系上方。这条判据把「判定哪些代数数可以写成根式」的问题，归结为理解数域的阿贝尔（乃至循环）扩张。遗憾的是，这个问题尚未解决，尽管人们可以梦想前沿研究正在逐步逼近它。不过，$\mathbb{Q}$ 的阿贝尔扩张是已知的：

\textbf{定理 1.1（Kronecker--Weber）.} $\mathbb{Q}$ 的每个阿贝尔扩张都包含于某个 $\mathbb{Q}(\zeta_n)$ 之中，其中 $\zeta_n$ 是本原 $n$ 次单位根。

也就是说，若 $f\in\mathbb{Q}[x]$ 的分裂域的 Galois 群阿贝尔，则 $f$ 的全部（等价地，某一些）根都可写成某个 $n$ 对应的 $\zeta_n$ 的有理函数。简要回顾：$[\mathbb{Q}(\zeta_n):\mathbb{Q}]=\varphi(n)$（$\varphi$ 为 Euler 函数），且 $\operatorname{Gal}(\mathbb{Q}(\zeta_n)/\mathbb{Q})=(\mathbb{Z}/n\mathbb{Z})^{\times}$，其中 $m\in(\mathbb{Z}/n\mathbb{Z})^{\times}$ 的作用为 $\zeta_n\mapsto\zeta_n^m$。对 $\mathbb{Q}$ 而言，cft 本质上等价于 Kronecker--Weber 定理，但它对一般数域给出了额外（尽管不显式）的掌控。

第二条线索是寻找代数数之间的恒等式。我们将会看到，Gauss 说明了 $\mathbb{Q}$ 中不明显的恒等式具有非平凡的算术后果。例如如下恒等式
\[
  \sqrt{2}=\zeta_8+\zeta_8^{-1}=\zeta_8+\bar\zeta_8=\frac{1+i}{\sqrt{2}}+\frac{1-i}{\sqrt{2}},
\]
\[
  \sqrt{-3}=\zeta_3-\zeta_3^{-1}=\zeta_3-\bar\zeta_3=\frac{-1+\sqrt{-3}}{2}-\frac{-1-\sqrt{-3}}{2},
\]
都由 Kronecker--Weber 定理所预言（因为这些数对应的 Galois 群是阿贝尔群 $\mathbb{Z}/2\mathbb{Z}$）。这些算术后果表明，我们应当尝试更充分地理解此类恒等式。

最后，第三个领域是丢番图方程的可解性。下面是典型定理的一个例子：

\textbf{定理 1.2（Hasse 原理）.} 设 $K$ 为数域，且
\[
  q(x_1,\ldots,x_n)=\sum_i a_i x_i^2+\sum_{i\neq j}a_{ij}x_ix_j,
  \qquad a_i,a_{ij}\in K.
\]
则对任意 $y\in K$，方程 $q(x_1,\ldots,x_n)=y$ 有解，当且仅当它在 $\mathbb{R}$ 中以及在所有 $\mathbb{Q}_p$（$p$ 遍历素数）中都有解。

在 $\mathbb{R}$ 上检验有解是容易的；在 $\mathbb{Q}_p$ 上，问题化归为初等的同余性质；而事实证明这个问题完全可以算法化地解决。我们可以把这类问题改述为：$y\in\mathbb{Q}$ 是否为二次扩张 $\mathbb{Q}(\sqrt{d})/\mathbb{Q}$ 中的范数——至少对形式 $N(x+y\sqrt{d})=x^2-dy^2$（$x,y\in\mathbb{Q}$）是如此，而这本来就是上述情形中最难的情形。这个问题，以及「联系局部与整体」这一更宏大的想法，将会再次出现。

现在我们转向 cft 主定理的陈述，它们或许眼下还不那么令人振奋。设 $K$ 是局部域，于是或者 $K$ 是阿基米德的，此时 $K=\mathbb{R},\mathbb{C}$；或者 $K$ 是非阿基米德的，此时 $K$ 是 $\mathbb{Q}_p$ 的有限扩张，或 $K=\mathbb{F}_{p^n}((t))$（对某对 $p,n$）。设 $K^{\mathrm{sep}}$ 为其可分闭包。注意：若 $K$ 是任意域，而 $K^{\mathrm{sep}}/K_1/K$ 与 $K^{\mathrm{sep}}/K_2/K$ 是阿贝尔扩张（它们必然是 Galois 的），则复合域 $K_1\cdot K_2$ 也阿贝尔，这为下述定义提供了依据：

\textbf{定义 1.3.} 极大阿贝尔扩张 $K^{\mathrm{ab}}/K$ 是所有阿贝尔扩张 $K^{\mathrm{sep}}/K'/K$ 的复合域，而 $\operatorname{Gal}^{\mathrm{ab}}(K):=\operatorname{Gal}(K^{\mathrm{ab}}/K)$ 是绝对 Galois 群 $\operatorname{Gal}(K):=\operatorname{Gal}(K^{\mathrm{sep}}/K)$ 的交换化。

我们还回忆下述定义：

\textbf{定义 1.4.} 对群 $G$，逆向极限
\[
  \widehat{G}:=\varprojlim_{\substack{H\triangleleft G\\ [G:H]<\infty}} G/H
\]
（对有限指数正规子群的商取极限）称为 $G$ 的投射有限完备化。

现在可以陈述局部类域论的主定理了：

\textbf{定理 1.5（lcft 主定理）.} 对任意有限扩张 $K/\mathbb{Q}_p$，存在投射有限群的标准同构
\[
  \operatorname{Gal}^{\mathrm{ab}}(K)\cong\widehat{K^{\times}}.
\]

\textbf{例 1.6.} $K^{\times}$ 的一阶结构由短正合列
\begin{equation}\tag{1.1}
  1\to\mathcal{O}_K^{\times}\to K^{\times}\xrightarrow{v}\mathbb{Z}\to 0
\end{equation}
给出，其中 $v$ 是相对极大理想 $\mathfrak{p}_K\subseteq\mathcal{O}_K$ 取的赋值同态（即它把 $\mathcal{O}_K$ 的一个素元送到 $1$）；整数环的单位群 $\mathcal{O}_K^{\times}$ 是投射有限的，且在 $K^{\times}$ 中开。$K^{\times}$ 的二阶结构由逆向极限
\[
  \mathcal{O}_K^{\times}=\varprojlim_n \mathcal{O}_K^{\times}/(1+\mathfrak{p}_K^n)
\]
给出，于是取投射有限完备化后，(1.1) 变为
\[
  1\to\widehat{\mathcal{O}_K^{\times}}\to\widehat{K^{\times}}\xrightarrow{\hat v}\widehat{\mathbb{Z}}\to 0.
\]

现在，对任意有限 Galois 扩张 $L/K$，$\operatorname{Gal}(L/K)$ 作用在 $L$ 上并保持 $\mathcal{O}_L$ 与 $\mathfrak{p}_L$。因此它下降为作用在 $k_L:=\mathcal{O}_L/\mathfrak{p}_L$ 上、固定 $k_K:=\mathcal{O}_K/\mathfrak{p}_K$ 的作用；这些都是有限域，记 $k_L=\mathbb{F}_{q^n}$、$k_K=\mathbb{F}_q$（$q$ 为某个素幂）。于是我们有映射
\[
  \operatorname{Gal}(L/K)\to\operatorname{Gal}(k_L/k_K)=\mathbb{Z}/n\mathbb{Z},
\]
当 $L/K$ 非分歧时它是同构；群 $\operatorname{Gal}(k_L/k_K)$ 由 Frobenius 自同构 $x\mapsto x^q$ 生成。对所有这些 $L/K$ 取逆向极限，便得到同态
\[
  \operatorname{Gal}(K)\to\operatorname{Gal}(k)=\widehat{\mathbb{Z}},
\]
它由 Frobenius 元给出，并按交换化的泛性质穿过 $\operatorname{Gal}^{\mathrm{ab}}(K)$。在 lcft 之下，这个映射与上面的映射 $\hat v$ 一致。

现在回忆，数域 $F$ 的阿代尔环定义为正向极限
\[
  \mathbb{A}_F:=\varinjlim_{S}\Big(\prod_{v\in S}F_v\times\prod_{v\notin S}\mathcal{O}_{F_v}\Big),
\]
其中 $S$ 取遍 $F$ 的素位的有限集。它带有对角嵌入
\[
  F\hookrightarrow\mathbb{A}_F,
\]
$F$ 经由它是离散嵌入的（想想 $\mathbb{Z}\hookrightarrow\mathbb{R}$）。粗略地说，$\mathbb{A}_F$ 汇集了 $F$ 的全部局部信息，而这个嵌入则编码了它的整体面向。$\mathbb{A}_F$ 内部有单位群 $\mathbb{A}_F^{\times}$，其拓扑由正向极限
\[
  \mathbb{A}_F^{\times}=\varinjlim_{S}\Big(\prod_{v\in S}F_v^{\times}\times\prod_{v\notin S}\mathcal{O}_{F_v}^{\times}\Big)
\]
给出（$S$ 同上，且各项在 $\mathbb{A}_F^{\times}$ 中开），而不是更细的子空间拓扑。现在我们可以陈述整体类域论（global class field theory，gcft）的主定理了：

\textbf{定理 1.7（gcft 主定理）.} 对任意有限扩张 $F/\mathbb{Q}$，存在投射有限群的标准同构
\[
  \operatorname{Gal}^{\mathrm{ab}}(F)\cong\widehat{\mathbb{A}_F^{\times}/F^{\times}}.
\]

这两个主定理以下述方式相容。若 $v$ 是全球域 $F$ 的一个素位（$\bar F$ 为其代数闭包），则我们有映射 $F\hookrightarrow F_v$ 与 $\bar F\hookrightarrow\bar F_v$，它们诱导出（单射）映射
\[
  \operatorname{Gal}(F_v)\to\operatorname{Gal}(F),\qquad
  \operatorname{Gal}^{\mathrm{ab}}(F_v)\to\operatorname{Gal}^{\mathrm{ab}}(F).
\]
下图
\[
  \begin{CD}
    F_v^{\times} @>{x\mapsto(1,\ldots,1,x,1,\ldots)}>> \mathbb{A}_F^{\times}/F^{\times}\\
    @V{\mathrm{LCFT}}VV @V{\mathrm{GCFT}}VV\\
    \operatorname{Gal}^{\mathrm{ab}}(F_v) @>>> \operatorname{Gal}^{\mathrm{ab}}(F)
  \end{CD}
\]
的竖直箭头，由「各群到其投射有限完备化的自然映射」与「局部、整体 cft 主定理的相应等同」复合得到，并且此图交换。

我们将先花几周时间为局部域与 $\mathbb{Q}$ 的二次扩张搭建 cft，因为这很好地捕捉了这门学科令人兴奋之处，而且比我们之后要发展的上同调方法更贴近手算。为开端，设 $K$ 是任意特征 $\operatorname{char}(K)\neq 2$ 的域。令 $\operatorname{Gal}_2(K)$ 为 $\operatorname{Gal}(K)$ 的极大商群，使得其中每个元素的平方都为 $1$。它必然阿贝尔，故有满射
\[
  \operatorname{Gal}^{\mathrm{ab}}(K)\twoheadrightarrow\operatorname{Gal}_2(K),
\]
并且它带有 $\mathbb{F}_2$-向量空间的结构。

\textbf{断言 1.8.} 存在 $\mathbb{F}_2$-向量空间的标准同构
\[
  \operatorname{Gal}_2(K)\cong\big(K^{\times}/(K^{\times})^2\big)^{\vee}
  :=\operatorname{Hom}\big(K^{\times}/(K^{\times})^2,\mathbb{Z}/2\mathbb{Z}\big).
\]

\textit{证明.} 我们先如下构造这样的映射：给定 $\sigma\in\operatorname{Gal}_2(K)$，定义 $\chi_\sigma\in\big(K^{\times}/(K^{\times})^2\big)^{\vee}$ 为
\[
  \chi_\sigma\colon K^{\times}/(K^{\times})^2\to\mathbb{Z}/2\mathbb{Z},\qquad
  d\mapsto
  \begin{cases}
    0, & \sigma(\sqrt{d})=\sqrt{d},\\
    1, & \sigma(\sqrt{d})=-\sqrt{d}.
  \end{cases}
\]
容易看出它确实是同态：给定 $\sigma_1,\sigma_2\in\operatorname{Gal}_2(K)$ 与 $d\in K^{\times}/(K^{\times})^2$，我们有
\[
  (\sigma_1\sigma_2)(\sqrt{d})=(-1)^{\chi_{\sigma_1}(d)}(-1)^{\chi_{\sigma_2}(d)}\sqrt{d},
\]
% [原文疑误：抽出文本第二个指数写作 χ_{σ1}(d)，按上下文应为 χ_{σ2}(d)；此处按更正排，见报告。]
从而 $\chi_{\sigma_1\sigma_2}=\chi_{\sigma_1}+\chi_{\sigma_2}$。

现在，由于源与靶都是投射有限的 $2$-挠阿贝尔群，只需证明该映射在取连续对偶后是同构。按通常的做法，我们有
\[
  \operatorname{Hom}_{\mathrm{cts}}\Big(\big(K^{\times}/(K^{\times})^2\big)^{\vee},\mathbb{Z}/2\mathbb{Z}\Big)
  =K^{\times}/(K^{\times})^2,
\]
而给出连续映射 $\operatorname{Gal}_2(K)\to\mathbb{Z}/2\mathbb{Z}$，等同于给出一个二次扩张 $K\subset F\subset K^{\mathrm{sep}}$，它形如 $K(\sqrt{d})$，其中 $d\in K^{\times}$ 在乘以平方的意义下确定。\hfill$\square$

\textbf{例 1.9.} 在 $K=\mathbb{Q}$ 的情形，此结果给出同构
\[
  \mathbb{Q}^{\times}/(\mathbb{Q}^{\times})^2
  \xrightarrow{\ \sim\ }\bigoplus_p \mathbb{Z}/2\mathbb{Z}\oplus\mathbb{Z}/2\mathbb{Z},
\]
其中直和遍历所有素数 $p$。然而，如何把这个群与投射有限完备化
\[
  \big[\mathbb{Q}^{\times}\backslash\mathbb{A}_{\mathbb{Q}}^{\times}/(\mathbb{A}_{\mathbb{Q}}^{\times})^2\big]^{\wedge}
\]
相比较，则完全不显然——而 gcft 主定理正是要我们做这样的比较。

现在，若 $K$ 是局部域，则 lcft 预言 $K^{\times}/(K^{\times})^2$ 标准地自对偶；实现这一对偶的配对
\[
  (\cdot,\cdot)\colon K^{\times}/(K^{\times})^2\times K^{\times}/(K^{\times})^2\to\{1,-1\}
\]
称为希尔伯特符号。我们接下来几讲的目标就是构造它，并证明事情确实如此。

\addcontentsline{toc}{section}{第 2 讲　希尔伯特符号}
% 第 2 讲 希尔伯特符号
% 源：src/lecture02.txt（Raskin, 18.786 讲义，2017-08-21 版）
% 说明：原文页眉「2. HILBERT SYMBOLS 7/9」为 PDF 页眉残留，已删。

\section*{第 2 讲\quad 希尔伯特符号}

设 $K$ 是 $\mathbb{Q}_p$ 上的局部域（其实任意局部域都够用），且 $\operatorname{char}(K)\neq 2$。注意这包含 $\mathbb{Q}_2$ 上的域，因为我们关心的是域的特征，而不是剩余域的特征。回忆上一讲的配对
\begin{equation}\tag{2.1}
  \operatorname{Gal}_2(K):=\operatorname{Gal}^{\mathrm{ab}}(K)/\{g^2:g\in\operatorname{Gal}^{\mathrm{ab}}(K)\}
  \cong\operatorname{Hom}\big(K^{\times}/(K^{\times})^2,\mathbb{Z}/2\mathbb{Z}\big),
\end{equation}
其中 $\operatorname{Gal}_2(K)$ 与 $\operatorname{Gal}^{\mathrm{ab}}(K)$ 是投射有限群，后者是 $K$ 的极大阿贝尔扩张的 Galois 群；而 $K^{\times}/(K^{\times})^2$ 是二元域 $\mathbb{Z}/2\mathbb{Z}$ 上有限维或无限维的向量空间（取对偶时，构成 $K^{\times}/(K^{\times})^2$ 的 $\mathbb{Z}/2\mathbb{Z}$ 直和变成了乘积，这反映了左端的投射有限性）。

还要回忆，lcft 断言 $K^{\times}\to\operatorname{Gal}^{\mathrm{ab}}(K)$ 是投射有限完备化，因此与 (2.1) 相反地有 $\operatorname{Gal}_2(K)\cong K^{\times}/(K^{\times})^2$。于是 lcft 预言存在如下形式的标准配对：

\textbf{定义 2.1.} 令希尔伯特符号
\[
  (\cdot,\cdot)\colon K^{\times}/(K^{\times})^2\times K^{\times}/(K^{\times})^2\to\{1,-1\}
\]
由
\[
  (a,b):=
  \begin{cases}
    1, & \text{若存在 } x,y\in K \text{ 使 } ax^2+by^2=1,\\
    -1, & \text{否则},
  \end{cases}
\]
定义，其中 $a,b\in K^{\times}$。

\textbf{注 2.2.} 这个定义只在局部域上性质良好。还要注意，$(a,b)$ 确实只在 $a,b$ 乘以平方的意义下定义，因为这些平方因子可以吸收进 $x$ 与 $y$。

\textbf{命题 2.3.} 希尔伯特符号满足以下性质：

(1) \emph{双乘性.} 对一切 $a,b,c\in K^{\times}$，
\[
  (a,bc)=(a,b)\cdot(a,c).
\]

(2) \emph{非退化性.} 对一切 $a\in K^{\times}$，若对一切 $b\in K^{\times}$ 都有 $(a,b)=1$，则 $a\in(K^{\times})^2$。

注意 $(a,b)=(b,a)$ 是平凡的。双乘性说的是：$ax^2+bcy^2=1$ 可解，当且仅当 $ax^2+by^2=1$ 与 $ax^2+cy^2=1$ 两者分别可解，或者两者都不可解。这有点奇怪，而且事实证明一般只在局部域上成立。
% [原文疑误：抽出文本两处均作 "ax²+by²=1"，按 (a,bc) 的含义第一个方程应为 ax²+bcy²=1；已按数学含义改正。]

\textbf{例 2.4.} 取 $K:=\mathbb{R}$。只要 $a$ 与 $b$ 不同时为负，$ax^2+by^2=1$ 就可解。因此 $\mathbb{R}^{\times}/(\mathbb{R}^{\times})^2=\{1,-1\}$（因为 $(\mathbb{R}^{\times})^2=\mathbb{R}_{>0}$），于是配对 $\{1,-1\}\times\{1,-1\}\to\{1,-1\}$ 确实非退化。

我们现在问：$x\in K^{\times}$ 何时是平方？当 $K=\mathbb{R},\mathbb{C}$ 时答案清楚。例如当 $x\in\mathbb{Q}_2^{\times}$ 时，可写 $x=2^{v(x)}y$，其中 $y\in\mathbb{Z}_2^{\times}$，而 $x$ 是平方当且仅当 $v_2(x)$ 为偶数且 $y$ 是平方（后者由习题集 1 第 1(c) 题将证明，等价于 $y\equiv 1\bmod 8$）。

设 $\mathfrak{p}\subseteq\mathcal{O}_K$ 为唯一的极大理想，$k:=\mathcal{O}_K/\mathfrak{p}$ 为剩余域，$\operatorname{char}(k)=p$ 为奇素数，$\pi\in\mathfrak{p}$ 为素元，即 $\pi\notin\mathfrak{p}^2$。

\textbf{断言 2.5.} 设 $x\in K^{\times}$，并写 $x=\pi^{v(x)}y$，其中 $y\in\mathcal{O}_K^{\times}$。则以下诸条等价：

(1) $x$ 是平方；

(2) $v(x)$ 为偶数，且 $y$ 是平方；

(3) $y\bmod\mathfrak{p}$ 是 $k^{\times}$ 中的平方。
% [原文疑误：抽出文本作 "a square in K^×"，按上下文应为 k^×（剩余域的单位群），见报告。]

注意我们可以化归到 $x\in\mathcal{O}_K^{\times}$。我们给出两个证明：

\textit{证明（用 Hensel 引理）.} 除最后一个条件外，其余解释都清楚。于是设 $x\bmod\mathfrak{p}$ 是 $\mathcal{O}_K^{\times}$ 中的平方。由 Hensel 引理，多项式 $p(t)=t^2-x\in\mathcal{O}_K[t]$ 有根 $r\in\mathcal{O}_K$，只要它有模 $\mathfrak{p}$ 的根 $\bar r$ 使 $\bar p'(\bar r)\neq 0$。第一个条件由假设成立：$x\bmod\mathfrak{p}=\bar r^2$ 对某个 $\bar r$ 成立。此时 $p'(t)=2t$，而 $\bar r\neq 0$（因 $\bar r^2=\bar x\neq 0$），又 $\operatorname{char}(k)\neq 2$，故 $\bar p'(\bar r)=2\bar r\neq 0$，于是第二个条件也成立。
\hfill$\square$

\textit{证明（显式）.} 考虑映射 $x\mapsto x^2$，即
\[
  \mathcal{O}_K^{\times}\xrightarrow{\ \sigma\ }S\subseteq\mathcal{O}_K^{\times},
  \qquad S:=\{x\in\mathcal{O}_K^{\times}:x \text{ 是模 } \mathfrak{p} \text{ 的平方}\}.
\]
我们想证明 $\sigma$ 满射，即 $\mathcal{O}_K$ 中每个模 $\mathfrak{p}$ 为平方的元素，都是 $\mathcal{O}_K^{\times}$ 中的平方。注意 $\mathcal{O}_K^{\times}$ 是带完全滤过的滤过阿贝尔群（见下文定义 2.7）
\[
  \mathcal{O}_K^{\times}\supseteq 1+\mathfrak{p}\supseteq 1+\mathfrak{p}^2\supseteq 1+\mathfrak{p}^3\supseteq\cdots,
\]
其中 $1+\mathfrak{p}^n$ 都是 $\mathcal{O}_K^{\times}$ 的开子群。显然 $\mathcal{O}_K^{\times}/(1+\mathfrak{p})=k^{\times}$；类似地 $(1+\mathfrak{p})/(1+\mathfrak{p}^2)\cong k$：对任意 $1+a\pi,1+b\pi\in 1+\mathfrak{p}$（$a,b\in\mathcal{O}_K/\mathfrak{p}$），有
\[
  (1+a\pi)(1+b\pi)=1+(a+b)\pi+ab\pi^2,
\]
而 $ab\pi^2\in\mathfrak{p}^2$，故在相伴分次项中剩下 $1+(a+b)\pi$，于是乘法恰好对应 $k$ 中的加法。类似地，对每个 $n\geq 1$，由同样的论证可得 $(1+\mathfrak{p}^n)/(1+\mathfrak{p}^{n+1})\cong k$，因为 $n+1\leq 2n$。现在 $\sigma$ 如下作用在滤过上：
\[
  \begin{CD}
    \mathcal{O}_K^{\times} @<<< 1+\mathfrak{p} @<<< 1+\mathfrak{p}^2 @<<< \cdots\\
    @V{\sigma}VV @V{\sigma}VV @V{\sigma}VV\\
    \mathcal{O}_K^{\times} @<<< S @<<< 1+\mathfrak{p} @<<< \cdots
  \end{CD}
\]
其中包含 $1+\mathfrak{p}\subseteq S$ 成立，因为 $1$ 平凡地是平方。现在，$\mathrm{Gr}_0$ 上的映射
\[
  \mathcal{O}_K^{\times}/(1+\mathfrak{p})=k^{\times}\xrightarrow{\ \sigma\ }(k^{\times})^2=S/(1+\mathfrak{p})
\]
平凡地满射（核很小）。此外，对每个 $n\geq 1$，$\mathrm{Gr}_n$ 上的映射是
\[
  (1+\mathfrak{p}^n)/(1+\mathfrak{p}^{n+1})\xrightarrow{\ \sigma\ }(1+\mathfrak{p}^n)/(1+\mathfrak{p}^{n+1}),
\]
由于对任意 $x\in k$ 有
\[
  (1+\pi^n x)^2=1+2x\pi^n+x^2\pi^{2n}\equiv 1+2x\pi^n \pmod{\mathfrak{p}^{n+1}}
\]
（因 $2n\geq n+1$），它等同于映射 $k\to k$，$x\mapsto 2x$；而 $\#k=p$ 为奇素数，故这是同构。于是 $\sigma$ 在每个分次项上都满射，由命题 2.9 便知 $\mathcal{O}_K^{\times}\xrightarrow{\sigma}S$ 满射，如所欲证。
\hfill$\square$

\textbf{注 2.6.} 一般说来，我们处理 $\mathcal{O}_K^{\times}$ 的工具是 $p$ 进指数映射与这个滤过；后者虽然只是抽象的形式体系，却比 $\mathcal{O}_K^{\times}$ 更简单，因为其商全都同构于有限域。作为一般原则，我们可以通过相伴分次 $\mathrm{Gr}_*A$ 来理解关于 $A$ 的许多事情。

\textbf{定义 2.7.} 设 $A$ 为阿贝尔群。$A$ 上的一个\emph{滤过}是一列递降的子群
\[
  A=:F^0A\supseteq F^1A\supseteq F^2A\supseteq\cdots,
\]
若 $A\xrightarrow{\sim}\varprojlim_n A/F^nA$，则称该滤过\emph{完全}。群 $\mathrm{Gr}_nA:=F^nA/F^{n+1}A$ 称为滤过的\emph{相伴分次项}。

\textbf{例 2.8.} 群
\[
  \mathcal{O}_K\cong\varprojlim_n \mathcal{O}_K/\mathfrak{p}^n,
  \qquad
  \mathcal{O}_K^{\times}\cong\varprojlim_n \mathcal{O}_K^{\times}/(1+\mathfrak{p}^n)
\]
都是完全滤过。

\textbf{命题 2.9.} 设 $f\colon A\to B$ 是完全滤过阿贝尔群之间的同态，即对每个 $n\geq 0$ 有 $f(F^nA)\subset F^nB$。若诱导映射
\[
  F^nA/F^{n+1}A\to F^nB/F^{n+1}B
\]
满射（相应地，单射），则 $f$ 满射（相应地，单射）。

\textit{证明.} 假设相伴分次上的映射都满射，且两个滤过都完全，如断言 2.5 的显式证明中那样。设有 $x\in B$，我们想解方程 $f(y)=x$，$y\in A$。先可解 $f(y_0)\equiv x\bmod F^1B$，于是 $x-f(y_0)\in F^1B$。其次，由假设相伴分次映射满射，可解 $f(\varepsilon_1)\equiv x-f(y_0)\bmod F^2B$，其中 $\varepsilon_1\in F^1A$ 是从 $\mathrm{Gr}_1A$ 提升来的「误差项」。注意由 $f$ 是同态，有
\[
  f(y_1)=f(y_0+\varepsilon_1)=f(y_0)+f(\varepsilon_1)\equiv x \pmod{F^2B},
\]
其中我们记 $y_1:=y_0+\varepsilon_1$。这与之前的方程形式相同，故可迭代得到「相容」的 $y_n$，使 $f(y_n)\equiv x\bmod F^{n+1}B$ 对每个 $n\geq 0$ 成立；这里「相容」指 $y_n\equiv y_{n+1}\bmod F^{n+1}A$。但于是有诱导元素 $y\in\varprojlim A/F^nA=A$ 对应 $(y_0,y_1,\ldots)$（注意 $y_n$ 在 $n$ 充分大时模 $F^nA$ 稳定），它满足原方程 $f(y)=x$，因为两个滤过都由假设完全。\hfill$\square$

\textbf{注 2.10.} 尽管简单而抽象，许多事情（如上一个断言）都可用上述命题轻松证明。用 Hensel 引理的办法的优点在于：这里我们需要用平方映射 $\sigma$ 是同态这一事实，而这一般不成立。尽管如此，这个办法还是能告诉我们，在剩余特征为奇的局部域中 $\mathcal{O}_K^{\times}$ 的哪些元素是平方。

结论是：当 $K$ 是剩余特征为奇的局部域时，有 $[K^{\times}:(K^{\times})^2]=4$，因为 $[\mathcal{O}_K^{\times}:(\mathcal{O}_K^{\times})^2]=2$，而 $2\mathbb{Z}\subseteq\mathbb{Z}$ 同理；于是 $K^{\times}/(K^{\times})^2$ 同构于 $\mathbb{Z}/2\mathbb{Z}\times\mathbb{Z}/2\mathbb{Z}$，它以 $\{\pi,r\}\subset\mathcal{O}_K^{\times}$ 为基，其中 $\pi$ 是素元，$r$ 不是模 $\mathfrak{p}$ 的平方（因此 $\pi$ 与 $r$ 当然不会相差一个平方）。

现在我们把希尔伯特符号用扩张上的范数重新表述；与最初的定义相反，这里我们将不对称地看待它。设 $a$ 不是平方，于是 $K(\sqrt{a})$ 是 $K$ 的二次扩张（注意若 $a$ 是非零平方，则只需把 $K(\sqrt{a})$ 理解为 $K$ 的相应平展扩张，同构于 $K\times K$）。

\textbf{断言 2.11.} $(a,b)=1$ 当且仅当 $b$ 是扩张 $K(\sqrt{a})/K$ 的范数，即 $K(\sqrt{a})$ 中有元素其范数为 $b$。

\textit{证明.} 设 $b$ 是范数，即存在 $\alpha,\beta\in K$ 使
\[
  \alpha^2-\beta^2a=N(\alpha+\beta\sqrt{a})=b,
\]
从而 $\alpha^2=a\beta^2+b$。若 $\alpha\neq 0$，则
\[
  a\Big(\frac{\beta}{\alpha}\Big)^2+b\Big(\frac{1}{\alpha}\Big)^2=1,
\]
故 $(a,b)=1$。若 $\alpha=0$，则 $b+\beta^2a=0$。令
\[
  x:=\frac{1}{2}\Big(1+\frac{1}{a}\Big),
  \qquad
  y:=\frac{1}{2\beta}\Big(1-\frac{1}{a}\Big),
\]
我们有
\[
  ax^2+by^2
  =a\cdot\frac{1}{4}\cdot\frac{(a+1)^2}{a^2}
  +(-\beta^2a)\cdot\frac{1}{4\beta^2}\cdot\frac{(a-1)^2}{a^2}
  =\frac{(a+1)^2-(a-1)^2}{4a}=1,
\]
故仍有 $(a,b)=1$。反向蕴涵只是 $\alpha\neq 0$ 时论证的平凡倒推。\hfill$\square$

我们不加证明地陈述关于希尔伯特符号的主要结果。重要的是：$L^{\times}$ 在范映射下的像不能太大（不是全部），也不能太小。我们将看到，这个定理等价于希尔伯特符号的非退化性。

\textbf{定理 2.12.} 若 $L/K$ 是局部域的二次扩张，则范映射 $N\colon L^{\times}\to K^{\times}$ 是同态，且 $N(L^{\times})\subseteq K^{\times}$ 是指数为 $2$ 的子群。

\textbf{例 2.13.} 考虑 $\mathbb{C}/\mathbb{R}$。
% [原文在 src 抽取中于此句后截断（原 PDF 该例的后续文字未进入文本层），此处照存。]

\addcontentsline{toc}{section}{第 3 讲　带驯分歧的范数群}
% 第 3 讲 带驯分歧的范数群
% 源：src/lecture03.txt（Raskin, 18.786 讲义，2017-08-21 版）
% 说明：原文页眉「3. NORM GROUPS WITH TAME RAMIFICATION 11/13」为 PDF 页眉残留，已删。

\section*{第 3 讲\quad 带驯分歧的范数群}

设 $K$ 是 $\operatorname{char}(K)\neq 2$ 的域。则
\[
  K^{\times}/(K^{\times})^2
  \cong\{\text{连续同态 } \operatorname{Gal}(K)\to\mathbb{Z}/2\mathbb{Z}\}
  \cong\{K \text{ 上的二次平展代数}\},
\]
这与断言 1.8 中最初的陈述对偶（这个结果是 Kummer 理论的一个雏形）。注意 $K$ 上的平展代数或是 $K\times K$，或是二次扩张 $K(\sqrt{d})/K$；前者对应 $K^{\times}/(K^{\times})^2$ 中平凡的平方陪集，后者对应由 $d\in K^{\times}$ 确定的陪集。

若 $K$ 是局部的，则 lcft 说 $\operatorname{Gal}^{\mathrm{ab}}(K)\cong\widehat{K^{\times}}$ 标准地成立。结合起来（因为这样的同态当然穿过 $\operatorname{Gal}^{\mathrm{ab}}(K)$），我们得到 $K^{\times}/(K^{\times})^2$ 有限且标准自对偶。这等价于断言存在一个「足够好」的配对
\[
  (\cdot,\cdot)\colon K^{\times}/(K^{\times})^2\times K^{\times}/(K^{\times})^2\to\{1,-1\},
\]
即它是\emph{双乘的}，满足
\[
  (a,bc)=(a,b)(a,c),\qquad (ab,c)=(a,c)(b,c),
\]
并且是\emph{非退化的}，满足条件
\[
  \text{若对一切 } b \text{ 有 } (a,b)=1,\text{ 则 } a\in(K^{\times})^2.
\]

我们曾能给出这个配对的一个简单定义，即
\[
  (a,b)=1 \iff ax^2+by^2=1 \text{ 在 } K \text{ 中有解}.
\]
注意由这个定义显然有 $(a,b)=(b,a)$，但不幸的是双乘性与非退化性都不显然，尽管本讲中我们会在许多情形证明它们成立。我们在断言 2.11 中已经证明希尔伯特符号的一个较不对称的定义成立，即对一切 $a$，
\[
  (a,b)=1 \iff b \text{ 是 } K(\sqrt{a})/K=K[t]/(t^2-a) \text{ 中的范数},
\]
而当 $a$ 是平方时，后者同构于 $K\times K$，一切元素都是范数。

在第 2 讲末尾，我们提出了下述断言，并指出重要的是这个范数子群「不能太大」（不是全部）、也「不能太小」，且 $K$ 必须是局部的。

\textbf{断言 3.1.} 这些「好性质」，即双乘性与非退化性，对希尔伯特符号成立，当且仅当对一切二次扩张 $L/K$，$N(L^{\times})\subseteq K^{\times}$ 是指数为 $2$ 的子群，即 $K^{\times}/N(L^{\times})=\mathbb{Z}/2\mathbb{Z}$。

\textit{证明.} 假设对一切二次扩张 $L/K$，$N(L^{\times})\subseteq K^{\times}$ 都是指数 $2$ 的子群。设 $a\in K^{\times}$。我们想证明
\begin{equation}\tag{3.1}
  (a,\cdot)\colon K^{\times}\to\{1,-1\}
\end{equation}
是同态，这等价于双乘性的第一个等式（另一个由对称性得到）。若 $a$ 是平方，则这很清楚，因为其像恒为 $1$（可取 $(x,y)=(1/\sqrt{a},0)$）。若 $a$ 不是平方，则令 $L:=K(\sqrt{a})$；由断言 2.11，$(a,b)=1$ 当且仅当 $b\in N(L^{\times})$。现在，与 $a$ 缩并的希尔伯特符号分解为
\[
  K^{\times}\twoheadrightarrow K^{\times}/NL^{\times}\cong\{1,-1\},
\]
其中同构是标准的，因为两个群的阶都为 $2$；我们用了范数群的指数为 $2$ 这一事实来构造保单位的最后这个 $2$ 阶群之间的双射，否则商就太大了。投影平凡地是同态。

现在证明非退化性。设 $a\notin(K^{\times})^2$。则存在某个 $b\in K^{\times}$ 不是来自 $L:=K(\sqrt{a})$ 的范数，而这当且仅当 $(a,b)=-1$，故由逆否命题，非退化性成立。

反之，注意若 $a\notin(K^{\times})^2$，则由非退化性，(3.1) 中的映射满射；又由双乘性它是同态。于是其核，即 $N(K(\sqrt{a})^{\times})\subseteq K^{\times}$，指数必为 $\#\{1,-1\}=2$。\hfill$\square$

\textbf{例 3.2.} 同样，基本情形是 $\mathbb{C}/\mathbb{R}$，其中范数群就是 $\mathbb{R}_{>0}$，指数为 $2$。

现在还需证明：

\textbf{定理 3.3.} 若 $L/K$ 是局部域的二次扩张，且 $\operatorname{char}(K)\neq 2$，则 $NL^{\times}\subseteq K^{\times}$ 是指数为 $2$ 的子群。

注意下述证明没有覆盖剩余特征为 $2$ 的分歧情形。

\textit{证明.} 令 $L:=K(\sqrt{d})$（其中 $d$ 如此前的 $a$），于是 $L$ 只依赖于 $d$ 模平方的类。则有两种情形：$v(d)=0$，当且仅当 $d\in\mathcal{O}_K^{\times}$；以及 $v(d)=1$，当且仅当 $d$ 是素元（因为可以不断约去 $\pi^2$ 的因子；注意这里 $v$ 照常是赋值）。
% [原文疑误/抽取缺字：抽出文本作 "which is true if and only if O^×_K"，缺 d∈O^×_K；已补。]

\emph{情形 1.} 此时 $d$ 是单位，且 $\sqrt{d}\in\mathcal{O}_L^{\times}$。这个扩张不一定非分歧，但我们只做非分歧情形，把分歧情形留到下周。分歧的例子是 $\mathbb{Q}_2(\sqrt{3})/\mathbb{Q}_2$（或 $\mathbb{Q}_2(\sqrt{2})$、$\mathbb{Q}_2(\sqrt{-1})$ 等）；扩张 $\mathbb{Q}_2(\sqrt{5})/\mathbb{Q}_2$ 非分歧。我们需要：
% [原文疑误：抽出文本作 "Here d is a square, and √d ∈ O^×_K"；按 v(d)=0 与下文，应为 d 是单位、√d∈O^×_L，见报告。]

\textbf{断言 3.4.} $N(\mathcal{O}_L^{\times})=\mathcal{O}_K^{\times}$；更一般地，$x\in K^{\times}$ 是范数当且仅当 $v(x)$ 为偶数（故 $K$ 中的素元不是范数）。

\textit{证明.} 我们使用「滤过技巧」。我们有下述滤过，它们被范同态保持：
\[
  \begin{CD}
    \mathcal{O}_L^{\times} @<<< 1+\mathfrak{p}_L @<<< 1+\mathfrak{p}_L^2 @<<< \cdots\\
    @V{N}VV @V{N}VV @V{N}VV\\
    \mathcal{O}_K^{\times} @<<< 1+\mathfrak{p}_K @<<< 1+\mathfrak{p}_K^2 @<<< \cdots
  \end{CD}
\]
在相伴分次项上，我们首先有
\[
  \begin{CD}
    \mathcal{O}_L^{\times}/(1+\mathfrak{p}_L) @>{N}>> \mathcal{O}_K^{\times}/(1+\mathfrak{p}_K)\\
    @V{\cong}VV @V{\cong}VV\\
    k_L^{\times} @>{N}>> k_K^{\times}.
  \end{CD}
\]
由于我们在非分歧情形，$k_L/k_K$ 与 $L/K$ 一样是二次扩张。为证明范映射在相伴分次项上满射，只需证明它在剩余域上满射，即：

\textbf{断言 3.5.} 范映射
\[
  N\colon\mathbb{F}_{q^2}^{\times}\to\mathbb{F}_q^{\times},\qquad x\mapsto x^{q+1}=\mathrm{Frob}(x)\cdot x
\]
满射（注意：因 $\operatorname{Gal}(\mathbb{F}_{q^2}/\mathbb{F}_q)\cong\mathbb{Z}/2\mathbb{Z}$，$x\mapsto x^q$ 是固定 $\mathbb{F}_q$ 的自同构；在 Galois 扩张中，域范数定义为一个元素的所有 Galois 共轭之积）。

\textit{证明 1.} 有限域的单位群必循环，故该映射对应于
\[
  \mathbb{Z}/(q^2-1)\mathbb{Z}\to\mathbb{Z}/(q-1)\mathbb{Z},\qquad n\mapsto(q+1)n.
\]
\hfill$\square$

\textit{证明 2（对 $p\neq 2$）.} 若 $x\in\mathbb{F}_q^{\times}$，则 $x=N(\sqrt{-x})$，且 $\sqrt{-x}\in\mathbb{F}_{q^2}^{\times}$，因为 $\mathbb{F}_{q^2}$ 是 $\mathbb{F}_q$ 唯一的二次扩张。\hfill$\square$

\textit{证明 3.} 我们有 $\#\mathbb{F}_{q^2}^{\times}=q^2-1$、$\#\mathbb{F}_q^{\times}=q-1$，且 $\#\ker(N)\geq(q^2-1)/(q-1)=q+1$（因像含于 $\mathbb{F}_q^{\times}$）；但多项式 $x^{q+1}-1$ 在 $\mathbb{F}_{q^2}$ 中恰有 $q+1$ 个根，因为 $\mathbb{F}_{q^2}/\mathbb{F}_q$ 可分（有限域是完美域），故 $\#\ker(N)\leq q+1$。两相结合即得 $\#\ker(N)=q+1$，从而 $N$ 满射。\hfill$\square$
% [原文疑误：抽出文本作 "#Ker(N) ≤ (q²−1)/(q−1)"，方向反了——像含于 F_q^× 只给出 ≥；已按数学含义改为 ≥，并补出两侧夹逼的结论。]
% [原文疑误：抽出文本作 "roots in F_q"，按上下文应为 F_{q²}（核在 F_{q²}^× 中），见报告。]

于是第一个相伴分次项 $\mathrm{Gr}_0$ 上的映射满射。在后续各项上，我们有
\[
  \begin{CD}
    (1+\mathfrak{p}_L^n)/(1+\mathfrak{p}_L^{n+1}) @>{N}>> (1+\mathfrak{p}_K^n)/(1+\mathfrak{p}_K^{n+1})\\
    @V{\cong}VV @V{\cong}VV\\
    k_L @>{T}>> k_K.
  \end{CD}
\]
要验证此图交换，注意由于我们假设了 $L/K$ 非分歧，$\pi$ 也是 $L$ 的素元（例如 $\mathbb{Q}_p(\sqrt{p})/\mathbb{Q}_p$ 是分歧扩张，$p$ 不再是 $\mathbb{Q}_p(\sqrt{p})$ 的素元）。于是在范映射下，
\[
  1+a\pi^n\xrightarrow{\ N\ }(1+a\pi^n)(1+\sigma a\pi^n)
  =1+(a+\sigma a)\pi^n+a\sigma a\pi^{2n}
  \equiv 1+Ta\pi^n \pmod{\pi^{n+1}}.
\]
再次地，我们提出下述断言：

\textbf{断言 3.6.} 迹映射
\[
  T\colon\mathbb{F}_{q^2}\to\mathbb{F}_q,\qquad x\mapsto x+x^q=x+\mathrm{Frob}(x)
\]
满射。

\textit{证明 1.} 迹映射的核对应 Artin--Schreier 多项式 $x^q+x$ 的根，该多项式可分，因而有 $q$ 个根，故 $\#\ker(T)=q$，且 $\#\operatorname{coker}(T)=q^2/q=q=\#\mathbb{F}_q$。\hfill$\square$

\textit{证明 2.} 若 $q$ 与 $2$ 互素，则对任意 $x\in\mathbb{F}_q$ 有 $x=T(x/2)$（如上进行）。否则若 $q=2^r$，则在 $\mathbb{F}_2$ 上 $x^2+x+1$ 是唯一的首一不可约多项式，而 $\mathbb{F}_4$ 是它的分裂域。一般地，对扩张 $\mathbb{F}_{2^{n+1}}/\mathbb{F}_{2^n}$，分裂多项式是 $x^2+x+\alpha$，其中我们选取使该多项式不可约的 $\alpha$。这恰好对一半的 $\alpha$ 选择有效，因为加法同态
\[
  \mathbb{F}_{2^n}\xrightarrow{\ x\mapsto x^2+x\ }\mathbb{F}_{2^n}
\]
的核是 $\mathbb{F}_2$，故其像的指数为 $2$。这些多项式的任何根迹都为 $1$（因为它们是首一的），而要得到任何其它迹的元素，只需乘以 $\mathbb{F}_{2^n}$ 中任意元素（因为迹映射是 $\mathbb{F}_{2^n}$-线性的）。\hfill$\square$

\textit{证明 3.} 扩张 $L/K$ 上的迹映射满射当且仅当该扩张可分，本情形如此，因为有限域是完美域。\hfill$\square$

于是，由于有限域上迹与范映射都满射，范映射在所有相伴分次项上满射，由命题 2.9 便知范映射在 $\mathcal{O}_L^{\times}$ 上满射，这证明了断言。\hfill$\square$

现在，为完成情形 1 的证明，我们想证明 $x\in K^{\times}$ 是范数当且仅当其赋值为偶数。为此注意，对任意 $y\in L^{\times}$ 有 $N(y)=y\cdot\sigma y$，且 $v(y\sigma(y))=2v(y)$，因为 $\operatorname{Gal}(L/K)$ 保持赋值。反之，只需注意 $\pi^2$ 是范数。于是若 $v(d)=0$，则 $NL^{\times}\subseteq K^{\times}$ 是指数为 $2$ 的子群，如所欲证。

\emph{情形 2.} 此时 $v(d)=1$，且仍有 $\operatorname{char}(K)\neq 2$。这保证了驯分歧，因为我们处理的是二次扩张（分歧指数不被 $p$ 整除；野分歧情形（分歧指数被 $p$ 整除）我们将在下一讲处理）。我们断言 $N(\mathcal{O}_L^{\times})\subseteq\mathcal{O}_K^{\times}$ 的指数为 $2$（具体地，$N(\mathcal{O}_L^{\times})=(\mathcal{O}_K^{\times})^2$），且存在某个 $\pi\in\mathcal{O}_K$ 既是素元又是范数。显然，这足以说明 $L^{\times}$ 的范数群在 $K^{\times}$ 中指数为 $2$。

令 $L:=K(\sqrt{d})$，其中 $v(d)=1$，故 $d$ 不是平方。则 $N(\alpha+\beta\sqrt{d})=\alpha^2-d\beta^2$，且若 $x\in(\mathcal{O}_K^{\times})^2$，则 $x=N(\sqrt{x})$，故 $x$ 是范数。反之，
\[
  v(\alpha^2-d\beta^2)=0=\min\{v(\alpha^2),v(d\beta^2)\}=v(\alpha^2),
\]
因为前者为偶数、后者为奇数，故两者不相等。由此 $\alpha,\beta\in\mathcal{O}_K^{\times}$，于是 $x=\alpha^2-d\beta^2$ 是模 $\mathfrak{p}$ 的平方，而这当且仅当 $x$ 是 $\mathcal{O}_K^{\times}$ 中的平方。最后，因为 $-d=N(\sqrt{d})$，故存在 $K$ 的一个素元是范数。
% [原文疑误：抽出文本作 "min{v(α²), v(β²+1)}"，按上下文应为 min{v(α²), v(dβ²)}，见报告。]

于是结论是：$x\in K^{\times}$ 是 $K(\sqrt{d})$ 的范数，当且仅当 $(-d)^{-v(x)}x$（一个整单位）是模 $\mathfrak{p}$ 的平方，故定理在此情形成立。\hfill$\square$

我们总结：驯分歧情形（就我们的目的而言包含非分歧情形）可以通过显式猜出 $NL^{\times}\subseteq K^{\times}$ 来处理。野分歧要棘手得多。这一切都归结为希尔伯特符号的显式公式，正如我们在习题集 1 第 2 题所见。对 $\mathbb{Q}_2$ 也有这样的公式，费点力气便可证明希尔伯特符号的一切「好」性质在此情形也成立（例如见 [Ser73]）。

\addcontentsline{toc}{section}{第 4 讲　整体类域论与二次互反律}
% 第 4 讲：整体类域论与二次互反律
% 源：src/lecture04.txt（379 行，PDF 抽出文本，公式按数学含义重排）
% 块：R2 / reports/R-block02.md

\section*{第 4 讲\quad 整体类域论与二次互反律}

上一讲中，我们把希尔伯特符号 $(\cdot,\cdot)$ 的非退化性与双乘性，归结为证明：对一切二次扩张 $L/K$（$K$ 为局部域），$NL^{\times}\subseteq K^{\times}$ 都是指数为 $2$ 的子群。我们已经证明，当扩张非分歧、以及当 $p=\operatorname{char}(\OK_K/\mathfrak{p})$ 为奇数时结论成立（$p=2$ 的情形已在习题集 1 第 1 题中大致证明）。本讲中，我们将在整体情形，即 $F=\Q$ 的情形，做类似的分析。

我们比较以下两个事实。第一，
\[
\Gal^{\mathrm{ab}}(\Q)/2\cong \operatorname{Hom}\bigl(\Q^{\times}/(\Q^{\times})^2,\{1,-1\}\bigr)=\operatorname{Hom}(\Q^{\times},\{1,-1\}),
\]
其中全体素数连同 $-1$ 构成这个群作为 $\F_2$-向量空间的一组基。第二，cft（类域论）预测
\[
\mathbb{A}^{\times}_{\Q}/\Q^{\times}\cong \Gal^{\mathrm{ab}}(\Q)
\]
是投射有限完备化之间的典范同构。于是我们期待
\begin{equation*}\tag{4.1}
\operatorname{Hom}(\Q^{\times},\{1,-1\})=\Gal^{\mathrm{ab}}(\Q)/2\cong \Q^{\times}\backslash \mathbb{A}^{\times}_{\Q}/(\mathbb{A}^{\times}_{\Q})^2 .
\end{equation*}

回忆下面的定义：

\textbf{定义 4.1.} \emph{阿代尔环}定义为
\[
\mathbb{A}_{\Q}:=\varinjlim_{S}\Bigl(\prod_{p\notin S}\Z_p\times\prod_{p\in S}\Q_p\Bigr),
\]
其中 $S$ 取遍 $\Q$ 的有限素位集合（素数与 $\infty$），按包含关系定向，并约定 $\Q_{\infty}=\R$。\emph{伊代尔群}是阿代尔环的乘法群
\[
\mathbb{A}^{\times}_{\Q}=\varinjlim_{S}\Bigl(\prod_{p\notin S}\Z^{\times}_p\times\prod_{p\in S}\Q^{\times}_p\Bigr),
\]
其中 $S$ 同上。
% [原文疑误] 抽出的伊代尔群定义中第一个乘积印作 ∏_{p∈S} Z_p^×，与阿代尔环的定义不符，应为 p∉S；已按后者排。

\textbf{例 4.2.} 我们有 $\widehat{\Z}=\prod_p \Z_p$，其中 $p$ 遍历所有素数；而 $\widehat{\Z}\times\R\subseteq \mathbb{A}_{\Q}$ 作为子环嵌入，其中可以把任意 $n\neq 0$ 对角嵌入。类似地，$\widehat{\Z}^{\times}=\prod_p \Z^{\times}_p$，且 $\widehat{\Z}^{\times}\times\R^{\times}\subseteq \mathbb{A}^{\times}_{\Q}$。然而 $2$ 与 $1/2$ 都不在 $\widehat{\Z}^{\times}\times\R^{\times}$ 中，因为它们不在 $\Z^{\times}_2$ 中；任何有理数都是如此。若我们把有理数加进来补偿，即 $\widehat{\Z}^{\times}\times\Q^{\times}\times\R^{\times}\to \mathbb{A}^{\times}_{\Q}$，则 $-1$ 在 $\Q^{\times}$ 与 $\R^{\times}$ 中被重复计入，所以必须把 $\R^{\times}$ 换成 $\R_{>0}$（见习题集 2 第 1 题）。

我们希望有一个阿代尔配对 $\mathbb{A}^{\times}_{\Q}\times\Q^{\times}\to\{1,-1\}$，从伊代尔与有理数映入 $\Z/2\Z$，它应当经过平方分解：
\[
\begin{CD}
\mathbb{A}^{\times}_{\Q}\times\Q^{\times} @>>> \{1,-1\}\\
@VVV @| \\
\Q^{\times}\backslash \mathbb{A}^{\times}_{\Q}/(\mathbb{A}^{\times}_{\Q})^2\times \Q^{\times}/(\Q^{\times})^2 @>>> \{1,-1\}
\end{CD}
\]
这里乘积左项中的那份 $\Q^{\times}$ 就是对角的「主伊代尔」，为了实现 (4.1)，这个配对也应当经过它分解。

这个映射应当诱导同构（某种意义上这正是「非退化性」）
\[
\Q^{\times}\backslash \mathbb{A}^{\times}_{\Q}/(\mathbb{A}^{\times}_{\Q})^2 \xrightarrow{\ \sim\ } \operatorname{Hom}\bigl(\Q^{\times}/(\Q^{\times})^2,\{1,-1\}\bigr),
\]
也就是说，这个映射不该恒等于 $1$，也不该是「什么古怪的东西」。

于是固定
\[
x=(x_p)_p\in \mathbb{A}^{\times}_{\Q},
\]
其中 $p$ 既可以是素数也可以是 $\infty$；并如下定义所求的配对：
\[
y\mapsto \prod_p (x_p,y)_p,
\]
这里我们把 $y$ 看作 $\Q_p$ 中的元素，而 $(\cdot,\cdot)_p$ 表示 $p$ 处的希尔伯特符号，即在 $\Q_p$ 上的希尔伯特符号。然而，这个无穷乘积是否良定义，先验地并不清楚；为此我们引入下面的引理：

\textbf{引理 4.3.} 对除有限多个 $p$ 之外的所有 $p$，$(x_p,y)_p=1$。

\textit{证明.} 事实上，对除有限多个 $p$ 之外的所有 $p$，都有 $p\neq 2,\infty$、$x_p\in\Z^{\times}_p$ 且 $y\in\Z^{\times}_p$；由驯符号的恒等式（因为 $x_p$ 与 $y$ 的赋值都是 $0$），这蕴含 $(x_p,y)_p=1$。\hfill$\square$

现在，这个映射当然是双乘性的（每一项都是），同样也当然地模平方分解，即它是映射
\[
\mathbb{A}^{\times}_{\Q}/(\mathbb{A}^{\times}_{\Q})^2\to\{1,-1\},
\]
因为若在任一侧乘上平方，希尔伯特符号不变。于是，这个映射在第一个因子上模 $\Q^{\times}$ 分解，当且仅当下面的断言成立：

\textbf{断言 4.4.} 对一切 $x,y\in\Q^{\times}$，都有 $\prod_p (x,y)_p=1$（也就是说，这个映射在第一个因子乘以 $\Q^{\times}$ 中的元素时不变；由双乘性，这意味着我们只是多提了一个这样的因子出来）。

\textbf{注 4.5.} 这对一切数域都成立（用一般的素位）。本讲中我们将看到，这（近似地）是二次互反律的一种重新包装。这一性质是素数之间的一种「共谋」：「$p$ 进域在幕后彼此通气；它们虽然彼此分离，却协力保证乘积为 $1$。」形容这类「共谋」的词就是「互反律」。

\textit{证明（断言 4.4）.} 首先，由于映射在乘以平方下不变，可以假设 $x=\pm p_1\cdots p_r$、$y=\pm q_1\cdots q_s$。于是双乘性蕴含我们可以取 $x\in\{-1,2,p\}$、$y\in\{-1,2,q\}$，其中 $p$ 与 $q$ 表示奇素数。我们只在 $p$ 与 $q$ 为互异的奇素数时证明该断言。

我们要证明
\[
(p,q)_{\infty}\times\prod_{\ell}(p,q)_{\ell}=1,
\]
其中 $\ell$ 遍历所有素数。第一项为 $1$，因为 $p$ 与 $q$ 都是正数；同样，$\ell$ 取异于 $p,q$ 的奇素数时也可以忽略。于是问题归结为证明
\[
(p,q)_2\cdot (p,q)_p\cdot (p,q)_q=1 .
\]

现在，正如习题集 1 第 2(b) 题所示，
\[
(p,q)_p=\left(\frac{(-1)^{v_p(p)v_p(q)}\,q^{v_p(p)}\big/p^{v_p(q)}}{p}\right)=\left(\frac{q}{p}\right),
\]
类似地 $(p,q)_q=\left(\frac{p}{q}\right)$。这里回忆
\[
\left(\frac{n}{p}\right):=\begin{cases} 1, & n \text{ 是模 } p \text{ 的平方},\\ -1, & \text{否则},\end{cases}
\]
其中 $n$ 与 $p$ 互素。此外，
\[
(p,q)_2:=(-1)^{\varepsilon(p)\varepsilon(q)}=(-1)^{\frac{p-1}{2}\cdot\frac{q-1}{2}},
\]
除非 $p\equiv q\equiv 3 \pmod 4$，它都等于 $1$。这样我们就归结到了初等的同余条件，而这正是二次互反律的陈述。\hfill$\square$

\textbf{注 4.6.} 二次互反律使得我们可以通过逐次归约来高效计算勒让德符号。

\textit{证明（二次互反律）.} 把勒让德符号看作映射
\[
\left(\frac{\cdot}{p}\right):\F^{\times}_p\to\{1,-1\},
\]
并把 $\F^{\times}_p$ 重新解释为 $\Gal(\Q(\zeta_p)/\Q)$，于是它是该 Galois 群唯一的非平凡二次特征标。这个特征标被 $\Q(\zeta_p)$（在 $\Q$ 上）唯一的二次子域所编码。我们希望把这个域写成 $\Q(\sqrt{d})$ 的形式。我们想要某个
\[
x=\sum_{n=1}^{p-1}x_n\zeta_p^{\,n}\in\Q(\zeta_p),
\]
使得对一切 $m\in(\Z/p\Z)^{\times}$，
\[
\sum_{n=1}^{p-1}x_n\zeta_p^{\,mn}=m\cdot x=\left(\frac{m}{p}\right)\cdot x=\sum_{n=1}^{p-1}\left(\frac{m}{p}\right)x_n\zeta_p^{\,n},
\]
这是作为 Galois 作用的等式：因为 $x_n\in\Q$ 不动，而作用为 $\zeta_p\mapsto\zeta_p^{\,m}$。也就是说，Galois 群把 $x$ 平移 $\pm 1$，这蕴含 $x\in\Q(\zeta_d)$，即我们的二次子域。
% [原文疑误] 此处原文印作 Q(ζ_d)；上文说二次子域形如 Q(√d)，Q(ζ_d) 疑为 Q(√d) 之误。照译并登记。

这个等式蕴含 $x_{mn}=\left(\frac{m}{p}\right)x_m$。反复解这个方程，我们最终得到元素（称为「高斯和」）
% [原文疑误] 原文此处印作 x_{mn}=(m/p)x_m；由下文 x_m=(m/p)x_1 及 Gauss 和的形式，右端下标应为 x_n。照译并登记。
\[
x=G:=\sum_{n=1}^{p-1}\left(\frac{n}{p}\right)\zeta_p^{\,n}.
\]
由构造我们知道 $G^2\in\Q$，但现在我们想知道它究竟是哪一个（事实上我们将看到，它是 $p$ 或 $-p$）。

设 $\chi:k^{\times}\to\C^{\times}$ 是乘法特征标，$\psi:k\to\C^{\times}$ 是加法特征标，其中 $k$ 是任意有限域。令
\[
G_{\chi,\psi}:=\sum_{x\in k^{\times}}\chi(x)\psi(x).
\]

\textbf{注 4.7.} 一个有趣的类比：伽马函数定义为
\[
\Gamma(\chi_s):=\int_{\R_{>0}}e^{-t}t^{s}\,\frac{dt}{t},
\]
这就像一个高斯和，其中取 $\psi(t):=e^{-t}$、$\chi_s(t):=t^{s}$。

我们需要下面的引理：

\textbf{引理 4.8.} 若 $\chi$ 与 $\psi$ 都不是平凡特征标，则 $G_{\chi,\psi}\cdot G_{\chi^{-1},\psi^{-1}}=\#k$。

\textit{证明.} 我们有
\[
\begin{aligned}
G_{\chi,\psi}\cdot G_{\chi^{-1},\psi^{-1}}
&=\sum_{x,y\in k^{\times}}\chi(x)\psi(x)\chi^{-1}(y)\psi^{-1}(y)\\
&=\sum_{x,y\in k^{\times}}\chi(x/y)\psi(x-y)\\
&=\sum_{z,y\in k^{\times}}\chi(z)\psi\bigl(y(z-1)\bigr),
\end{aligned}
\]
其中作了变量替换 $z:=x/y$，即 $x=zy$。现在，若 $z\neq 1$，则 $y(z-1)$ 取遍 $k^{\times}$ 中所有值，于是（由非退化性）$\sum_{w\in k}\psi(w)=0$ 成立。因此我们得到
\[
\begin{aligned}
G_{\chi,\psi}\cdot G_{\chi^{-1},\psi^{-1}}
&=\sum_{z\in k^{\times}}\chi(z)\sum_{y\in k^{\times}}\psi\bigl(y(z-1)\bigr)\\
&=\sum_{z\in k^{\times}\setminus\{1\}}\chi(z)\bigl(-\psi(0)\bigr)+\chi(1)\sum_{y\in k^{\times}}\psi(0)\\
&=\chi(1)+1\cdot\#k^{\times}=1+\#k^{\times}=\#k,
\end{aligned}
\]
因为类似地 $\sum_{w\in k^{\times}}\chi(w)=0$。\hfill$\square$

现在，当 $\psi$ 对应 $\zeta_p$ 的某个幂、$\chi$ 取乘法勒让德特征标时，我们想知道 $G_{\chi^{-1},\psi^{-1}}$ 是什么。我们有
\[
G_{\chi^{-1},\psi^{-1}}=\sum_{n=1}^{p-1}\left(\frac{n}{p}\right)\zeta_p^{-n}
=\sum_{n=1}^{p-1}\left(\frac{-n}{p}\right)\zeta_p^{\,n}
=\left(\frac{-1}{p}\right)\sum_{n=1}^{p-1}\left(\frac{n}{p}\right)\zeta_p^{\,n}
=\left(\frac{-1}{p}\right)G .
\]
于是
\[
G_{\chi,\psi}\cdot G_{\chi^{-1},\psi^{-1}}=G\cdot\left(\frac{-1}{p}\right)G=p,
\]
从而
\[
G^2=\left(\frac{-1}{p}\right)\cdot p=(-1)^{(p-1)/2}\cdot p,
\]
而 $G$ 是 $p$ 或 $-p$ 的平方根，视条件而定。

现在回忆：$\Q(\zeta_p)/\Q$ 在 $q\neq p$ 处非分歧，并且有同构
\[
\Gal(\Q(\zeta_p)/\Q)\cong(\Z/p\Z)^{\times},\qquad \mathrm{Frob}_q\mapsto q \bmod p .
\]
于是 $\left(\frac{q}{p}\right)=1$ 当且仅当 $\mathrm{Frob}_q$ 固定 $G$；事实上，
\[
\begin{aligned}
\mathrm{Frob}_q(G)
&=\mathrm{Frob}_q\Bigl(\sum_{n=1}^{p-1}\left(\frac{n}{p}\right)\zeta_p^{\,n}\Bigr)
=\sum_{n=1}^{p-1}\left(\frac{n}{p}\right)\zeta_p^{\,qn}
=\sum_{n=1}^{p-1}\left(\frac{q^{-1}n}{p}\right)\zeta_p^{\,n}\\
&=\sum_{n=1}^{p-1}\left(\frac{qn}{p}\right)\zeta_p^{\,n}
=\sum_{n=1}^{p-1}\left(\frac{q}{p}\right)\left(\frac{n}{p}\right)\zeta_p^{\,n}
=\left(\frac{q}{p}\right)G .
\end{aligned}
\]
% 原文此处的中间一步抽出为 "( n/q" 套在勒让德符号里，按 q 的逆元排作 (q^{-1}n/p)。

此外，
\[
G^{q-1}=(G^2)^{(q-1)/2}=\bigl((-1)^{(p-1)/2}\cdot p\bigr)^{(q-1)/2}\equiv (-1)^{\frac{p-1}{2}\cdot\frac{q-1}{2}}\left(\frac{p}{q}\right)\pmod q,
\]
所以
\[
\left(\frac{q}{p}\right)G^2=G^{q+1}\equiv (-1)^{\frac{p-1}{2}\cdot\frac{q-1}{2}}\left(\frac{p}{q}\right)G^2 \pmod q .
\]
除以 $G^2=(-1)^{(p-1)/2}\cdot p$（它模 $q$ 可逆）之后，得到
\[
\left(\frac{q}{p}\right)\equiv (-1)^{\frac{p-1}{2}\cdot\frac{q-1}{2}}\left(\frac{p}{q}\right)\pmod q,
\]
即
\[
\left(\frac{p}{q}\right)\left(\frac{q}{p}\right)=(-1)^{\frac{p-1}{2}\cdot\frac{q-1}{2}},
\]
如所愿。\hfill$\square$

\addcontentsline{toc}{section}{第 5 讲　阿代尔配对的非退化性与正合列}
% 第 5 讲，源文件 src/lecture05.txt（317 行）；R3 块。
\section*{第 5 讲\quad 阿代尔配对的非退化性与正合列}

回忆一下，我们想要一个非退化的配对，而
\begin{equation}\tag{5.1}
\Q^\times\backslash\mathbb{A}_\Q^\times/(\mathbb{A}_\Q^\times)^2\times\Q^\times/(\Q^\times)^2\longrightarrow\{1,-1\},\qquad
((x_p),r)\longmapsto\prod_p (x_p,r)_p
\end{equation}
曾是候选者（与之前一样，$p$ 取遍所有素数以及 $\infty$）。这个配对的良定义性归结为互反律
\[
\prod_p (x,y)_p=1\qquad (x,y\in\Q^\times).
\]
我们看到，当 $x=p$、$y=q$ 为奇素数时，它归结为二次互反律
\[
\left(\frac{p}{q}\right)\left(\frac{q}{p}\right)=(-1)^{\frac{p-1}{2}\cdot\frac{q-1}{2}},
\]
而我们是通过考虑特征标
\[
\chi\colon\Gal(\Q(\zeta_p)/\Q)=(\Z/p\Z)^\times\xrightarrow{\ (\cdot/p)\ }\{1,-1\}
\]
来证明它的。我们看到，这对应于 $\Q(\zeta_p)$ 的一个唯一二次子扩张
\[
\Q(\sqrt{\pm p})=\Q\Bigl(\sqrt{\left(\tfrac{-1}{p}\right)p}\Bigr),
\]
其中的要点是高斯和
\[
G:=\sum_{a=1}^{p-1}\left(\frac{a}{p}\right)\zeta_p^a=\sqrt{\left(\frac{-1}{p}\right)p}.
\]

更一般地，若 $F/\Q$ 是 Galois 扩张、Galois 群为 $G$，且 $\Q$ 的一个素位 $q$ 非分歧，则 $[\mathrm{Frob}_q]\mapsto[1]\in G$（这里两边都是共轭类）当且仅当 $q$ 在 $F$ 中分裂。于是
\[
\left(\frac{q}{p}\right)=1\iff
\left(\frac{\left(\tfrac{-1}{p}\right)p}{q}\right)=1,
\]
因为右端等价于 $q$ 在扩张 $\Q(\sqrt{\pm p})$ 中分裂，而这又蕴含
\[
\left(\frac{q}{p}\right)=\left(\frac{\left(\tfrac{-1}{p}\right)p}{q}\right)
=\left(\frac{\left(\tfrac{-1}{p}\right)}{q}\right)\left(\frac{p}{q}\right)
=(-1)^{\frac{p-1}{2}\cdot\frac{q-1}{2}}\left(\frac{p}{q}\right),
\]
由此即得所需结论。

类似地，我们可以在其他情形下得到互反律的相应结论，例如：

\textbf{命题 5.1.} 我们有
\[
\prod_p (2,\ell)_p=1,
\]
其中 $\ell$ 为奇素数，$p$ 取遍所有素数（注意 $(2,\ell)_\infty=1$ 是平凡的）。

\textit{证明.} 与之前一样，若 $p\neq 2,\ell$ 则 $(2,\ell)_p=1$；而利用驯符号，
\[
(2,\ell)_\ell=\left(\frac{(-1)^{v(\ell)v(2)}2^{v(\ell)}/\ell^{v(2)}}{\ell}\right)
=\left(\frac{2}{\ell}\right).
\]
由习题集 1 第 2(d) 题所得的公式，我们有
% [已核：原文该式指数中原有两处在 PDF 抽取里成了不可识别字符（U+000F），按原 PDF 第 1177 行（Serre《A Course in Arithmetic》第 III 章的 Q_2 公式）读作 ε(1)ε(ℓ)+v(2)θ(ℓ)+v(ℓ)θ(1)；
%  由于 ε(1)=0、v(ℓ)=0，该式化为 (-1)^{θ(ℓ)}，与右端结论一致]
\[
(2,\ell)_2=(-1)^{\varepsilon(1)\varepsilon(\ell)+v(2)\theta(\ell)+v(\ell)\theta(1)}=(-1)^{\theta(\ell)},
\]
其中
\[
(-1)^{\theta(\ell)}:=
\begin{cases}
1, & \ell\equiv 1,-1\pmod 8,\\
-1, & \ell\equiv 3,-3\pmod 8,
\end{cases}
\]
这对应于 $\Z/16\Z$ 中 $\ell^2$ 到 $\Z/2\Z$ 的标准同构。因此，我们希望证明
\[
\left(\frac{2}{\ell}\right)=(-1)^{\theta(\ell)}.
\]
为了判断 $2$ 是否为模 $\ell$ 的平方，我们想要 $\sqrt2$ 的一个便利表达式，也就是 $\Q(\sqrt2)$ 的一个分圆嵌入（在其中 $\ell$ 分裂当且仅当 $\left(\frac{2}{\ell}\right)=1$）。回忆若 $\zeta_8$ 为本原八次单位根，则可取 $\zeta_8=\frac{\sqrt2}{2}+i\frac{\sqrt2}{2}$，于是
\[
\zeta_8+\zeta_8^{-1}=\zeta_8+\overline{\zeta}_8=2\operatorname{Re}(\zeta_8)=\sqrt2.
\]
在代数上，可以这样看出这一恒等式：
\[
(\zeta_8+\zeta_8^{-1})^2=2+\zeta_8^2+\zeta_8^{-2}=2+\zeta_4+\zeta_4^{-1}=2,
\]
因为 $\zeta_4$ 与 $\zeta_4^{-1}$ 恰为 $i$ 与 $-i$。这给出 $\Q(\sqrt2)\subseteq\Q(\zeta_8)$，以及特征标
\[
\Gal(\Q(\zeta_8)/\Q)=(\Z/8\Z)^\times\xrightarrow{\ \chi\ }\{1,-1\}.
\]
我们断言 $\left(\frac{2}{\cdot}\right)=\chi=(-1)^{\theta(\cdot)}$。显然 $\ker\bigl((-1)^{\theta(\cdot)}\bigr)=\{1,-1\}$，而 $n\in(\Z/8\Z)^\times$ 属于 $\ker(\chi)$ 当且仅当它固定 $\sqrt2=\zeta_8+\zeta_8^{-1}$，即
\[
\zeta_8+\zeta_8^{-1}\longmapsto\zeta_8^n+\zeta_8^{-n}=\zeta_8+\zeta_8^{-1},
\]
而这仅在 $n=1$（两项都不动）或 $n=-1$（两项互换）时成立。于是两个核相同，从而两个函数相等。\hfill$\square$

注意，我们也可以不用 Galois 群来论证。若假定 $\zeta_8+\zeta_8^{-1}\in\F_\ell$，那么要验证这一点，只需检查
\[
\zeta_8+\zeta_8^{-1}=(\zeta_8+\zeta_8^{-1})^\ell=\zeta_8^\ell+\zeta_8^{-\ell},
\]
也就是说它在 Frobenius 元的作用下不动；于是得到与前面相同的条件。

其他符号验证起来相对繁琐：例如 $\prod_p(2,2)_p=1$ 很简单，因为由习题集 1 第 2(d) 题有 $(2,2)_2=1$；而 $\prod_p(-1,\ell)_p=1$ 则由
\[
\left(\frac{-1}{p}\right)=(-1)^{\frac{p-1}{2}}
\]
看出。至此我们已验证了初始配对 (5.1) 的良定义性。现在我们要验证这个配对是非退化的。注意这一步其实并不需要互反律，因为论证反而更容易。

\textbf{命题 5.2.} (5.1) 中定义的映射
\[
\chi\colon\Q^\times\backslash\mathbb{A}_\Q^\times/(\mathbb{A}_\Q^\times)^2
\xrightarrow{\ \sim\ }\operatorname{Hom}(\Q^\times,\{1,-1\})\cong\Gal_2(\Q)
\]
是同构（注意：配对只从 $\Q^\times/(\Q^\times)^2$ 出发定义映射并无妨碍，因为右端的同态会吸收平方）。

\textit{证明.} 在习题集 2 第 1(a) 题中，我们证明了
\[
\mathbb{A}_\Q^\times=\widehat{\Z}^\times\times\R_{>0}\times\Q^\times,
\]
其中前两项经由局部素位嵌入，最后一项沿对角线嵌入。模去 $\Q^\times$ 消去最后一项，模去平方消去第二项，于是由中国剩余定理得到
\begin{equation}\tag{5.2}
\Q^\times\backslash\mathbb{A}_\Q^\times/(\mathbb{A}_\Q^\times)^2=\prod_p\Z_p^\times/(\Z_p^\times)^2,
\end{equation}
其中 $p$ 取遍所有素数；当 $p$ 为奇素数时，$\Z_p^\times/(\Z_p^\times)^2$ 是 2 阶群，由任一二次非剩余生成；而 $\Z_2^\times/(\Z_2^\times)^2\cong(\Z/8\Z)^\times$ 是 4 阶群，由 $-1$ 与 $5$ 生成。又有
\[
\Q^\times/(\Q^\times)^2=\{\pm p_1\cdots p_r:\ p_i\ \text{为素数}\}
=\Z/2\Z\oplus\bigoplus_p\Z/2\Z,
\]
其中 $p$ 同上，第一个 $\Z/2\Z$ 对应符号。我们将看到，在对偶之后，这些 $\Z/2\Z$ 的副本全都（几乎）一一对上。

设 $p$ 为奇素数，$r$ 为 $p$ 处的二次非剩余，即 $\Z_p^\times/(\Z_p^\times)^2$ 中的非平凡元。则
\[
(\chi(r)(\pm q))_p=(r,\pm q)_p
=\left(\frac{r^{v(\pm q)}/(\pm q)^{v(r)}}{p}\right)
=\begin{cases}1, & q\neq p,\\ -1, & q=p,\end{cases}
\]
是 $\chi$ 在 $\Q^\times\backslash\mathbb{A}_\Q^\times/(\mathbb{A}_\Q^\times)^2$ 的第 $p$ 个分量上所取的值，其中 $q$ 为素数。对最后一个基元素，有 $\chi(r)(-1)=1=(r,-1)_p$。于是 $p$ 处那些显然的（拓扑）基元素一一对上；现在我们必须问 $p=2$ 时会发生什么。一个自然的猜测是「在 $2$ 处取 $r=-1$ 或 $r=5$、其余各处取 $r=1$」所定义的伊代尔，因为由习题集 2 第 2(b) 题，$5$ 对应 $\Q_2$ 的唯一非分歧二次扩张。计算可得
\[
\begin{aligned}
\chi(5,1,1,\dots)(q)&=(5,q)_2=1,\\
\chi(5,1,1,\dots)(-1)&=(5,-1)_2=1,\\
\chi(5,1,1,\dots)(2)&=(5,2)_2=(-1)^{\theta(5)}=-1,
\end{aligned}
\]
可见这个基元素与 $2$ 处的基元素完美对应。这里我们用元组表示伊代尔，其坐标是就同构 (5.2) 而言取的，素数按通常顺序排列。于是
\[
\chi(-1,1,1,\dots)(-1)=-1=(-1,-1)_2
\]
完成了证明。不过，这两组基实际上并不完美对应，因为与另一个奇素数 $p$ 配对会给出与 $-1$ 是否为模 $p$ 的平方相应的符号；但我们可以用「上三角矩阵」轻易地把一组基用另一组表示出来，办法是对 $(-1,1,1,\dots)$ 这个基元素作修正（本质上，我们有一个对角线上全为 $1$ 的无限矩阵，只是在 $(-1,1,1,\dots)$ 处例外，它对应于 $\Q^\times/(\Q^\times)^2$ 的给定基中一个更复杂的元素）。

下面给一个稍微更严肃的论证。我们有如下行正合的交换图：
\[
\begin{CD}
0 @>>> \Z/2\Z\times\prod_{p\neq2}\Z/2\Z @>>>
\Q^\times\backslash\mathbb{A}_\Q^\times/(\mathbb{A}_\Q^\times)^2 @>>> \Z/2\Z @>>> 0\\[2.2ex]
@. @V{\sim}VV @V{\sim}VV @V{\varphi\mapsto\varphi(-1)}VV @.\\[2.2ex]
0 @>>> \{\varphi\colon\Q^\times\to\Z/2\Z\mid\varphi(-1)=1\} @>>>
\operatorname{Hom}(\Q^\times,\Z/2\Z) @>>> \Z/2\Z @>>> 0
\end{CD}
\]
第一个 $\Z/2\Z$ 对应伊代尔 $(5,1,1,\dots)$，其余的副本对应每个 $p$ 处的二次非剩余；在最右端的 $\Z/2\Z$ 中，我们得到 $(-1,1,1,\dots)$ 的像。两个映入 $\Z/2\Z$ 的映射都是商映射；右侧的竖直映射是同构，因为它非平凡；左侧的竖直映射也是同构，因为如我们先前所见，一切完美对应。于是中间的映射是同构，此即所求。\hfill$\square$

现在回到这个问题：证明对局部域的任何二次扩张 $L/K$，都有 $\#(K^\times/NL^\times)=2$。回忆这一陈述等价于希尔伯特符号的双乘性与非退化性；我们已经对奇素数以及非分歧、驯分歧扩张的情形证明了它，但对野分歧扩张、以及对 $\Q_2$ 上除 $\Q_2$ 本身之外的扩张还证不出来。我们现在的目标是更一般地证明它：即证明若 $L/K$ 是 $n$ 次循环扩张（也就是说，它是 Galois 群为 $\Z/n\Z$ 的 Galois 扩张），则 $\#(K^\times/NL^\times)=n$。为了把它放进更一般的框架：若 $L/K$ 是有限阿贝尔扩张，我们实际上期望有典范同构
\[
\Gal(L/K)\cong K^\times/NL^\times,
\]
所以期望的比数值相等更多。我们将用正合列与同调代数的方法来证明这一点，现在转入这一话题。

事实证明，短正合列是确定有限阿贝尔群阶数的极好工具。设我们有短正合列
\[
0\to M\overset{g}{\hookrightarrow}E\overset{f}{\twoheadrightarrow}N\to 0.
\]
则 $M=\ker(f)$，$N=\operatorname{coker}(g)=E/M$；用滤链的语言说，$M$ 与 $N$ 就像相伴分次项。确实，可以把 $M$ 与 $N$ 看作「原子」，把 $E$ 看作「分子」，而分子的精细结构决定它的「反应」。显然若 $M$ 与 $N$ 有限，则 $E$ 也有限，且 $\#E=\#M\#N$。野分歧带来的问题在于：$L^\times$ 上没有滤链。

一个问题是，许多运算并不保持短正合列。例如，若 $n\ge1$ 为整数，模去 $n$ 并不保持 $\#(E/nE)$。

\textbf{例 5.3.} (1) 若我们有正合列
\[
0\to\Z/n\Z\xrightarrow{x\mapsto(x,0)}\Z/n\Z\times\Z/n\Z
\xrightarrow{(x,y)\mapsto y}\Z/n\Z\to 0,
\]
则模去 $n$ 保持它。

(2) 若我们有正合列
\[
0\to\Z/n\Z\xrightarrow{x\mapsto xn}\Z/n^2\Z\xrightarrow{1\mapsto1}\Z/n\Z\to 0,
\]
则模去 $n$ 把它变成
\[
\Z/n\Z\to\Z/n\Z\xrightarrow{\operatorname{id}}\Z/n\Z\to 0.
\]
我们有同样的「原子」，但它们构成了不同的「分子」。在前一情形，模去 $n$ 之后的阶是 $n^2$，而在这里是 $n$。

同调代数的一个中心论点是：我们可以通过延长正合列来修正这一点。打个诗意的比方，被改变的正合列就像可见光；它缺了红外波段，而延长正合列就能让我们看到这一段。具体说来，这对应于 $n$-挠：$\Z/n^2\Z$ 没有 $\Z/n\Z\times\Z/n\Z$ 那么多的 $n$-挠，后者全是 $n$-挠。

对模 $M$，记 $M[n]\subseteq M$ 为 $\ker(n\colon M\to M)$。于是我们得到更长的正合列
\[
0\to M[n]\to E[n]\to N[n]\xrightarrow{\ \delta\ }M/n\to E/n\to N/n\to 0,
\]
其中对 $x\in N[n]$，$\delta(x)=ny$，这里 $y\in E$ 是满足 $f(y)=x$ 的任意元素；注意 $f(ny)=nx=0$，故 $\delta(x)\in M\subseteq E$，此即所求。我们将在下一讲证明这是一个正合列。

\addcontentsline{toc}{section}{第 6 讲　正合列与 Tate 上同调}
% 第 6 讲，源文件 src/lecture06.txt（208 行）；R3 块。
\section*{第 6 讲\quad 正合列与 Tate 上同调}

上一讲我们开始讨论一些简单的同调代数；动机是计算某些有限阿贝尔群的阶（特别是 $K^\times/N(L^\times)$，其中 $L/K$ 是局部域的循环扩张）。回忆下面的定义：

\textbf{定义 6.1.} 序列
\[
\cdots\to X_{n-1}\xrightarrow{d_n}X_n\xrightarrow{d_{n+1}}X_{n+1}\to\cdots
\]
称为正合的，若对每个 $n$ 都有 $\ker(d_{n+1})=\operatorname{im}(d_n)$；我们把 $d_i$ 称为微分。

要解这个方程，通常的做法是证明：若 $d_{n+1}$ 杀死某个元素，则该元素落在 $d_n$ 的像中。我们看到，对短正合列 $0\to M\hookrightarrow E\twoheadrightarrow N\to 0$，有 $M=E/N$ 与 $\#E=\#M\cdot\#N$，所以短正合列是度量阿贝尔群大小的有效手段。我们也看到，对任何这样的短正合列与 $n\ge1$，都有长正合列
\begin{equation}\tag{6.1}
0\to M[n]\to E[n]\to N[n]\xrightarrow{\ \delta\ }M/n\to E/n\to N/n\to 0,
\end{equation}
其中回忆
\[
M[n]:=\{x\in M:nx=0\}=\operatorname{Tor}_1(M,\Z/n)=H_1(M\otimes^{\mathbf{L}}\Z/n),
\]
分别表示 $n$-挠子群与第一同调群，$E$、$N$ 类似。边界映射 $\delta$ 把元素 $x\in N[n]$ 提升为 $\tilde x\in E$，使得 $n\tilde x\in M$（因为 $N$ 中 $nx=0$），然后把 $n\tilde x$ 映到它在 $M/n$ 中的等价类。还须验证以下断言：

\textbf{断言 6.2.} 边界映射 $\delta$ 良定义。

\textit{证明.} 设 $\tilde{\tilde x}$ 是 $x$ 的另一个提升。则 $\tilde x-\tilde{\tilde x}\in M$，因为它在 $N$ 中的像为零；于是 $n(\tilde x-\tilde{\tilde x})\in nM$，故在 $M/nM$ 中 $n\tilde x=n\tilde{\tilde x}$。\hfill$\square$

\textbf{断言 6.3.} (6.1) 中的序列正合。

% [原文疑误：原 PDF 第 28 页该证明首局作 "then x̃ ∈ M, and therefore x = 0 in N"，论证跳步
%  （由 δ(x)=nx̃=0 于 M/n 只能得 nx̃∈nM，即 x̃ 与 M 中某元素相差一个 E[n] 元素）；
%  译文按标准论证补全——从而 x 落在 E[n]→N[n] 的像中]
\textit{证明.} 除边界映射处之外，其余各处的正合性都是显然的。若对某个 $x\in N[n]$ 及其提升 $\tilde x\in E$ 有 $\delta(x)=n\tilde x=0$ 于 $M/n$，则 $n\tilde x\in nM$，故 $\tilde x$ 与 $M$ 中某个元素只相差一个 $E[n]$ 中的元素，从而 $x$ 落在 $E[n]\to N[n]$ 的像中。类似地，若 $x\in M/n$ 的像在 $E/n$ 中为零，则 $\tilde x=ny$，其中 $y\in E$，而 $\tilde x$ 是 $x$ 到 $M$ 的提升。投射到 $N$，由正合性得 $0=n\bar y$，从而 $\bar y\in N[n]$。于是再次由正合性得 $ny\in M$，且 $\delta(\bar y)=ny=x$（作为 $M/n$ 中的类），此即所求。\hfill$\square$

我们有下面这个有用的引理：

\textbf{引理 6.4.} 设
\[
0\to X_0\xrightarrow{d_1}X_1\xrightarrow{d_2}\cdots\xrightarrow{d_{n-1}}X_{n-1}\xrightarrow{d_n}X_n\to 0
\]
正合，且所有 $X_i$ 有限。则
\[
\#X_0\cdot\#X_2\cdots=\#X_1\cdot\#X_3\cdots.
\]

\textit{证明.} 对 $n$ 归纳。$n=1$ 时结论显然，故设它对 $n-1$ 成立。构造正合列
\[
0\to X_0\to\cdots\to X_{n-2}\to\ker(d_{n-1})\to 0
\]
与
\[
0\to\ker(d_{n-1})\to X_{n-1}\xrightarrow{d_{n-1}}X_n\to 0
\]
（后者的正合性来自 $d_{n-1}$ 满射到 $X_n$，因为原序列以 $X_n\to 0$ 结尾）。设 $n$ 为偶数。对第一个正合列用归纳假设，得
% [原文疑误：抽出文本作 "#X0···#Xn = #X0···#X_{n−1}·#X_{n−1}/#Im(d_{n−1}) = #X1···#Im(d_{n−1})·X_{n−1}/#Im(d_{n−1})"，
%  其中 #X_{n−1} 出现两次、指标链自相矛盾（按原文第一行，两端相差一个 #X_{n−1}/#Im 因子）。
%  此处按标准证明重写：改用 ker(d_{n−1}) 而非 Im(d_{n−1})，并把归纳用在长度为 n 的截断序列上。]
\[
\#X_0\#X_2\cdots\#X_{n-2}=\#X_1\#X_3\cdots\#X_{n-3}\cdot\#\ker(d_{n-1}),
\]
又由第二个正合列有 $\#X_{n-1}=\#\ker(d_{n-1})\cdot\#X_n$。于是
\[
\#X_0\#X_2\cdots\#X_n
=\#X_1\#X_3\cdots\#X_{n-3}\cdot\#\ker(d_{n-1})\cdot\#X_n
=\#X_1\#X_3\cdots\#X_{n-1},
\]
即所欲证。$n$ 为奇数的情形类似。\hfill$\square$

\textbf{定义 6.5.} 设 $M$ 为阿贝尔群，且 $M/n$ 与 $M[n]$ 都有限。则
\[
\chi(M):=\chi_n(M):=\frac{\#(M/n)}{\#(M[n])}
\]
称为 $M$ 的欧拉示性数。

\textbf{例 6.6.} (1) 若 $M$ 有限，则 $\chi(M)=1$。为此注意
\[
0\to M[n]\to M\xrightarrow{\ n\ }M\to M/n\to 0
\]
正合，于是由引理 6.4 得 $\#(M[n])\cdot\#M=\#M\cdot\#(M/n)$。

(2) 若 $M=\Z$，则 $\chi(M)=n$，因为 $M[n]=0$，而 $M/n=\Z/n\Z$ 的阶为 $n$。

下面这个引理是关于欧拉示性数的一个重要事实：

\textbf{引理 6.7.} 对短正合列 $0\to M\to E\to N\to 0$，若三个阿贝尔群中有两个的 $\chi$ 存在，则第三个的 $\chi$ 也存在，且 $\chi(M)\cdot\chi(N)=\chi(E)$；这里「存在」指的是（比方说对 $M$）$M/n$ 与 $M[n]$ 都有限。

\textit{证明.} 有正合列
\[
0\to M[n]\to E[n]\to N[n]\to M/n\to E/n\to N/n\to 0.
\]
更一般地，注意若 $X_{n-1}\xrightarrow{d_{n-1}}X_n\xrightarrow{d_n}X_{n+1}$ 正合，而 $X_{n-1}$ 与 $X_{n+1}$ 有限，则 $X_n$ 也有限，因为有短正合列
\[
0\to\operatorname{im}(d_{n-1})=\ker(d_n)\to X_n\to\operatorname{im}(d_n)\to 0,
\]
其中两端群有限，故 $\#X_n=\#\ker(d_n)\cdot\#\operatorname{im}(d_n)$ 也有限。于是序列中所有群都有限，且由引理 6.4，
\[
\#(M[n])\cdot\#(N[n])\cdot\#(E/n)=\#(E[n])\cdot\#(M/n)\cdot\#(N/n),
\]
由此即得所要的表达式。\hfill$\square$

作为应用，我们来计算 $\#(K^\times/(K^\times)^n)$。注意
\[
\chi(K^\times)=\frac{\#(K^\times/(K^\times)^n)}{\#(K^\times[n])},
\]
其中分母是 $K$ 中 $n$ 次单位根的个数。此外我们还有正合列
\[
0\to\mathcal{O}_K^\times\to K^\times\xrightarrow{\ v\ }\Z\to 0,
\]
于是由引理 6.7，$\chi(K^\times)=\chi(\mathcal{O}_K^\times)\chi(\Z)=n\chi(\mathcal{O}_K^\times)$。因此我们真正想算的是 $\chi(\mathcal{O}_K^\times)$。

一个好用的启发是：若 $\mathcal{O}_K^\times$ 包含某个开子群（即有限指数子群）$\Gamma$，则 $\Gamma\cong\mathcal{O}_K^+$；当 $\operatorname{char}(K)=0$ 时这可由 $p$ 进指数映射得到。于是
\begin{equation}\tag{6.2}
\chi(\mathcal{O}_K^\times)=\chi(\Gamma)\chi(\mathcal{O}_K^\times/\Gamma)=\chi(\Gamma)=\chi(\mathcal{O}_K)
\end{equation}
（末一个 $\chi$ 是对加法群取的），因为按假设 $\mathcal{O}_K^\times/\Gamma$ 有限。又加法群有 $\mathcal{O}_K[n]=0$（因为 $\mathcal{O}_K$ 是整环），且 $\chi(\mathcal{O}_K)=\#(\mathcal{O}_K/n)=|n|_K^{-1}$，其中 $|x|_K:=q^{-v(x)}$ 表示 $K$ 中正规化的绝对值（即对素元 $\pi$ 有 $v(\pi)=1$），$q$ 表示剩余域的阶。所得的公式
\begin{equation}\tag{6.3}
\#(K^\times/(K^\times)^n)=\frac{n\cdot\#(K^\times[n])}{|n|_K}
\end{equation}
重现了习题集 1 第 1(b) 题中当 $n=2$ 时已证的结果（不过同样的方法对一般的 $n$ 也适用）。不用指数映射的证明用到了这样的事实：对充分大的 $N$，
\[
1+\mathfrak{p}^N\xrightarrow{\ x\mapsto x^n\ }1+\mathfrak{p}^{N+v(n)}
\]
是同构（这可用滤链证明；它是我们先前那个加性陈述的乘性版本）。

现在我们引入循环群的 Tate 上同调。

\textbf{定义 6.8.} 若 $G$ 是（不必有限的）群，则 $G$-模 $A$ 是一个阿贝尔群，$G$ 通过群自同构作用在 $A$ 上。等价地，存在同态 $G\to\operatorname{Aut}(A)$，且该作用满足
\begin{enumerate}[label=(\arabic*),leftmargin=2.2em,itemsep=0pt,topsep=2pt]
\item $g\cdot(a+b)=g\cdot a+g\cdot b$，
\item $(gh)\cdot a=g\cdot(h\cdot a)$，
\end{enumerate}
对一切 $g,h\in G$、$a,b\in A$。

\textbf{例 6.9.} 若 $L/K$ 是域扩张，$G:=\Gal(L/K)$，则 $L$ 与 $L^\times$ 都是 $G$-模，因为域自同构同时保持这两种运算。这将是我们主要关心的例子。

现在假定 $G$ 有限，$A$ 为 $G$-模。

\textbf{定义 6.10.} 第一个 Tate 上同调群为
\[
\widehat{H}^0(G,A):=A^G/N(A),
\]
其中
\[
A^G:=\{a\in A:\ g\cdot a=a\ \text{对一切}\ g\in G\}
\]
表示不动点集。

注意范映射定义为
\[
N\colon A\to A,\qquad a\mapsto\sum_{g\in G}g\cdot a,
\]
所以我们确实需要 $G$ 有限的假设。此外，这个表达式表明范映射经由 $A^G\subseteq A$ 分解，即其像落在 $A^G$ 中。

\textbf{例 6.11.} (1) 回到例 6.9 并取 $A=L$，有 $A^G=K$，而 $N\colon L\to K$ 是域迹，故 $\widehat H^0(L/K)=K/T(L)=0$，因为 $L/K$ 必可分。

(2) 若 $A=L^\times$，则 $(L^\times)^G=K^\times$，且 $\widehat H^0(L^\times)=K^\times/N(L^\times)$。于是我们先前的问题现在重新表述为：对局部域的循环扩张 $L/K$，计算 $\widehat H^0(G,L^\times)$。

(3) 若 $A$ 是任意阿贝尔群，则称 $G$ 平凡地作用在 $A$ 上，若对一切 $g\in G$、$a\in A$ 有 $g\cdot a=a$。此时 $\widehat H^0(G,A)=A/(\#G)A$。于是 Tate 上同调的概念完全推广了我们先前的讨论。

\textbf{定义 6.12.} $G$-模的映射（或称 $G$-态射，或其他任何合理的叫法）$A\xrightarrow{f}B$ 是保持 $G$ 作用的群同态，即对一切 $g\in G$、$a\in A$ 有 $f(g\cdot a)=g\cdot f(a)$。

$G$-模的（短）正合列是指阿贝尔群的（短）正合列，但其中所有映射都是 $G$-态射。

\textbf{例 6.13.} $1\to\mathcal{O}_L^\times\to L^\times\xrightarrow{v}\Z\to1$ 是 $G$-模的短正合列，其中 $G:=\Gal(L/K)$，而 $G$ 平凡地作用在 $\Z$ 上、并通过 Galois 作用作用在 $\mathcal{O}_L^\times$ 上。

现在设 $0\to A\to B\to C\to 0$ 是 $G$-模的短正合列。则我们得到正合列
\begin{equation}\tag{6.4}
\widehat H^0(G,A)\xrightarrow{\ \alpha\ }\widehat H^0(G,B)\xrightarrow{\ \beta\ }\widehat H^0(G,C),
\end{equation}
其中 $\alpha$ 不必是单射（正如上一讲群作用平凡时所见到的那样），$\beta$ 不必是满射。这是因为 Tate 上同调涉及两种操作：其一取不动点，是左正合而非右正合的；其二取商，是右正合而非左正合的。

现在设 $G=\Z/n\Z$，$\sigma\in G$ 为生成元（即 $1$）。

\textbf{定义 6.14.} 第二个 Tate 上同调群为
\[
\widehat H^1(G,A):=\ker(N\colon A\to A)/(1-\sigma)A.
\]

注意我们之所以取商，是因为对任意形如 $x:=y-\sigma y$（$y\in A$）的元素，都有
\[
N(x)=x+\sigma x+\cdots+\sigma^{n-1}x
=y-\sigma y+\sigma y-\sigma^2 y+\cdots+\sigma^{n-1}y-\underbrace{\sigma^{n}y}_{y}=0,
\]
而我们希望把这些平凡情形从核中略去。

现在断言：对正合列 $0\to A\to B\to C\to 0$，存在正合列
\begin{equation}\tag{6.5}
\widehat H^0(A)\to\widehat H^0(B)\to\widehat H^0(C)\xrightarrow{\ \delta\ }
\widehat H^1(A)\to\widehat H^1(B)\to\widehat H^1(C),
\end{equation}
其边界映射 $\delta$ 把任意 $x\in C^G/N(C)$ 提升为 $\tilde x\in B$，然后取 $(1-\sigma)\tilde x$。由于 $x\in C^G$，在 $C$ 中有 $(1-\sigma)x=0$，故 $(1-\sigma)\tilde x\in A$。此外，$(1-\sigma)\tilde x$ 显然被 $A$ 中的范映射杀死，因此它给出 $\widehat H^1(G,A)$ 中的一个类。同样地，我们验证以下两点：

\textbf{断言 6.15.} 边界映射 $\delta$ 良定义，即它不依赖于 $\tilde x$ 的选取。

\textit{证明.} 若 $\tilde{\tilde x}$ 是另一个提升，则 $\tilde x-\tilde{\tilde x}\in A$（因为 $C\cong B/A$），故 $(1-\sigma)(\tilde x-\tilde{\tilde x})$ 在 $\widehat H^1(G,A)$ 中为零。\hfill$\square$

\textbf{断言 6.16.} 序列 (6.5) 可延展为正合的。

\textit{证明.} 与前面一样，我们只在边界映射处验证。取 $x\in B^G/N(B)$，它在 $\widehat H^1(A)$ 中的像是 $(1-\sigma)x=0$。若 $x\in\ker(\delta)$，则对 $x$ 的某个提升 $\tilde x$ 有 $\tilde x\in B^G$，从而落在 $\widehat H^0(B)$ 中。

取 $x\in C^G/N(C)$，它在 $\widehat H^1(A)$ 中的像是 $(1-\sigma)\tilde x$，其中 $\tilde x$ 是 $x$ 到 $B$ 的提升；按定义它在 $\widehat H^1(B)$ 中被杀死。若 $x\in\widehat H^1(A)$ 在 $\widehat H^1(B)$ 中为零，则 $x\in(1-\sigma)B$，从而 $x\in\operatorname{im}(\delta)$。\hfill$\square$

\addcontentsline{toc}{section}{第 7 讲　链复形与 Herbrand 商}
\section*{第 7 讲\quad 链复形与 Herbrand 商}

上一讲我们对循环群定义了 Tate 上同调群 $\HH^0(G,M)$ 与 $\HH^1(G,M)$。回忆：若 $G=\Z/n\Z$、生成元为 $\sigma$，则一个 $G$-模就是带自同构 $\sigma:M\xrightarrow{\sim}M$ 且满足 $\sigma^n=\mathrm{id}_M$ 的交换群 $M$。我们的主要例子来自域扩张 $L/K$，其中 $\Gal(L/K)=G$，此时 $L$ 与 $L^\times$ 都是 $G$-模。于是
\[
\HH^0(G,M):=M^G/N(M)=\ker(1-\sigma)\Big/\Big\{\sum_{i=0}^{n-1}\sigma^i m:m\in M\Big\},
\]
\[
\HH^1(G,M):=\ker(N)/(1-\sigma),
\]
因为 $\ker(1-\sigma)$ 中的元素在 $\sigma$ 作用下不动，从而在 $G$ 作用下不动。我们的目标是在上述例子中证明 $\#\HH^0=n$，所用的工具是长正合列。

我们看到，若
\[
0\to M\to E\to N\to 0
\]
是 $G$-模的短正合列（即 $M,E,N$ 都是带 $n$ 阶自同构的交换群，这些自同构与上述映射相容，且 $N=E/M$，于是 $M$ 在 $N$ 的自同构下不动），则得到长正合列
\[
\HH^0(G,M)\to\HH^0(G,E)\to\HH^0(G,N)\xrightarrow{\delta}\HH^1(G,M)\to\HH^1(G,E)\to\HH^1(G,N),
\]
其中边缘映射 $\delta$ 把 $x\in\HH^0(N)=N^G/N(N)$ 提升为 $\tilde x\in E$，使得 $(1-\sigma)\tilde x\in\ker(N)\subseteq M$，从而给出 $\HH^1(G,M)$ 中的一个类。

现在定义第二个边缘映射
\begin{equation}\tag{7.1}
\HH^1(G,M)\to\HH^1(G,E)\to\HH^1(G,N)\xrightarrow{\partial}\HH^0(G,M)\to\HH^0(G,E)\to\HH^0(G,N),
\end{equation}
它把 $x\in\HH^1(G,N)$ 提升为元素 $\tilde x\in E$。此时 $N(\tilde x)=\sum_{i=0}^{n-1}\sigma^i\tilde x\in M^G$，因为它被 $1-\sigma$ 杀死，于是它定义了 $\HH^0(G,M)$ 中的一个类。我们验证：

\textbf{断言 7.1.} 边缘映射 $\partial$ 是良定义的。

\textit{证明.} 若 $\tilde x'$ 是 $x$ 的另一个提升，则 $\tilde x-\tilde x'\in M$（因为 $N=E/M$），从而 $\sum_{i=0}^{n-1}\sigma^i(\tilde x-\tilde x')\in N(M)$，在 $\HH^0(G,M)$ 中被消掉。\hfill$\square$

\textbf{断言 7.2.} (7.1) 中的序列正合。

\textit{证明.} 若 $x\in\HH^1(G,E)$，则 $N(x)=0$，故 $\partial(x)=0$ 于 $\HH^0(G,M)$。若 $x\in\ker(\partial)$，则对 $x$ 的某个提升 $\tilde x\in E$ 有 $N(\tilde x)=0$，而 $x$ 正是 $\tilde x$ 的像。

若 $x\in\HH^1(G,N)$、取提升 $\tilde x\in E$，则由定义 $\partial(x)=N(\tilde x)$ 在 $\HH^0(G,E)$ 中为零。若 $x\in\HH^0(G,M)$ 在 $\HH^0(G,E)$ 中为零，则 $x\in N(E)$，从而 $x\in\im(\partial)$。\hfill$\square$

于是我们得到循环群 Tate 上同调的「$2$-周期」正合列，这引出如下定义。

\textbf{定义 7.3.} 对每个 $i\in\Z$（正、负均可），定义
\[
\HH^i(G,M):=
\begin{cases}
\HH^0(G,M), & i\equiv 0 \pmod 2,\\
\HH^1(G,M), & i\equiv 1 \pmod 2.
\end{cases}
\]
这一良好性质对非循环群不成立，因此我们常设法把上同调化归到循环群的情形。

作为另一种写法，置
\begin{equation}\tag{7.2}
\cdots\xrightarrow{\sum_{i=0}^{n-1}\sigma^i}M\xrightarrow{1-\sigma}M\xrightarrow{\sum_{i=0}^{n-1}\sigma^i}M\xrightarrow{1-\sigma}\cdots,
\end{equation}
并注意它构成我们将要称为链复形的东西：

\textbf{定义 7.4.} 链复形 $X_\bullet$ 是一列
\[
\cdots\xrightarrow{d_{-2}}X_{-1}\xrightarrow{d_{-1}}X_0\xrightarrow{d_0}X_1\xrightarrow{d_1}X_2\xrightarrow{d_2}\cdots,
\]
% [原文疑误：原文作 $\ker(d_{i+1})\supset\ker(d_i)$，按 $d_{i+1}d_i=0$ 应为 $\mathrm{im}(d_i)\subseteq\ker(d_{i+1})$；译文径改]
满足对一切 $i\in\Z$ 有 $d_{i+1}d_i=0$（也就是说 $\im(d_i)\subseteq\ker(d_{i+1})$，但不必像正合列那样取等号）。于是定义 $X_\bullet$ 的第 $i$ 个上同调为
\[
H^i(X_\bullet):=\ker(d_i)/\im(d_{i-1}).
\]
因此，长正合列是一种特殊的链复形。注意 (7.2) 满足该定义，因为
\[
(1-\sigma)\sum_{i=0}^{n-1}\sigma^i x=\sum_{i=0}^{n-1}\sigma^i x-\sum_{i=0}^{n-1}\sigma^{i+1}x=Nx-Nx=0,
\]
且两个映射显然交换。Tate 上同调群于是就是该链复形的上同调，由此可以清楚地看出它们具有 $2$-周期性。

\textbf{定义 7.5.} $G$-模 $M$ 的 \textbf{Herbrand 商}或\textbf{欧拉示性数}为
\[
\chi(M):=\frac{\#\HH^0(G,M)}{\#\HH^1(G,M)},
\]
它只在分子与分母都有限时才有定义。

此定义推广了我们先前对平凡 $G$-模的讨论：此时 $\HH^0(G,M)=M/n$、$\HH^1(G,M)=M[n]$；不过要注意，从偶上同调到奇上同调的边缘映射为零。

\textbf{引理 7.6.} 设
\[
0\to M\to E\to N\to 0
\]
是 $G$-模的短正合列。若 $\chi$ 对三个 $G$-模中的两个有定义，则它对三个都有定义，此时
\[
\chi(M)\cdot\chi(N)=\chi(E).
\]

\textit{证明.} 作出长正合列
\[
0\to\ker(\alpha)\to\HH^0(M)\xrightarrow{\alpha}\HH^0(E)\to\HH^0(N)\xrightarrow{\delta}\HH^1(M)\to\HH^1(E)\xrightarrow{\beta}\HH^1(N)\to\coker(\beta)\to 0 .
\]
由于第二个边缘映射给出正合列
\[
\HH^1(E)\xrightarrow{\beta}\HH^1(N)\xrightarrow{\partial}\HH^0(M)\xrightarrow{\alpha}\HH^0(E),
\]
我们有
\[
\ker(\alpha)=\im(\partial)=\HH^1(N)/\ker(\partial)=\HH^1(N)/\im(\beta)=\coker(\beta).
\]
应用引理 6.4 并消去 $\#\ker(\alpha)$ 与 $\#\coker(\beta)$，即得欲证之结论（与引理 6.7 的情形相同）。\hfill$\square$

关于有限性，插入一段简短的讨论：

\textbf{断言 7.7.} 群 $\HH^0(G,M)$ 与 $\HH^1(G,M)$ 都是 $n$-挠的。

\textit{证明.} 设 $x\in M^G$。则 $N(x)=\sum_{i=0}^{n-1}\sigma^i x=\sum_{i=0}^{n-1}x=nx$。于是 $nx\in N(M)$，故 $\HH^0(G,M)$ 是 $n$-挠的。现设 $x\in\ker(N)$。则
\[
nx=nx-Nx=\sum_{i=1}^{n}(1-\sigma^i)x=(1-\sigma)\sum_{i=1}^{n}(1+\cdots+\sigma^{i-1})x,
\]
因此 $nx\in(1-\sigma)M$，$\HH^1(G,M)$ 同样是 $n$-挠的。\hfill$\square$

于是，$\HH^0(G,M)$ 与 $\HH^1(G,M)$ 的有限生成性蕴含有限性。现在回忆我们的目标：对局部域的循环 $n$ 次扩张 $L/K$，证明 $\#\HH^0(L^\times)=n$。我们有如下更精细的论断：

\textbf{断言 7.8.} 保持上述设定，则
\begin{enumerate}
\item[(1)] $\HH^1(L^\times)=0$（从而 $\chi(L^\times)=\#\HH^0(L^\times)$）；
\item[(2)] $\chi(\OK_L^\times)=1$；
\item[(3)] $\chi(L^\times)=n$。
\end{enumerate}

\textit{证明.} 先证 (2) 蕴含 (3)。有正合列
\[
1\to\OK_L^\times\to L^\times\xrightarrow{v}\Z\to 0,
\]
其中 $v$ 表示赋值。于是由引理 7.6，
\[
\chi(L^\times)=\chi(\OK_L^\times)\cdot\chi(\Z)=1\cdot n=n
\]
（用到 (2)），这里注意
\[
\HH^0(\Z)=\Z^G/N\Z=\Z/n\Z,\qquad
\HH^1(\Z)=\ker(N)/(1-\sigma)=0 .
\]
下面证明 (2)。

\textbf{引理 7.9.} 若 $M$ 是有限 $G$-模，则 $\chi(M)=1$。

\textit{证明.} 有正合列
\[
0\to M^G\xrightarrow{1-\sigma}\ker(N)\to\HH^1(G,M)\to 0,
\]
\[
0\to\ker(N)\to M\xrightarrow{\sum_{i=0}^{n-1}\sigma^i}M^G\to\HH^0(G,M)\to 0,
\]
于是由引理 7.6，
\[
\#\ker(N)\cdot\#M^G=\#M\cdot\#\HH^0(G,M),\qquad
\#M^G\cdot\#\ker(N)=\#M\cdot\#\HH^1(G,M),
\]
故 $\#\HH^0(G,M)=\#\HH^1(G,M)$，即 $\chi(M)=1$，如所欲证。\hfill$\square$

类似地有 $\chi(\OK_L)=1$，其中把 $\OK_L$ 视为加法群。事实上，更容易建立的说法是 $\chi(L)=1$。直观上，这是因为我们工作在 $p$ 进数上，一切都必须是 $\Q$-向量空间，故 $n$ 可逆；而由断言 7.7，我们的上同调群都是 $n$-挠的，因此两个上同调群都必须消没，从而 $\chi(L)=1$。

由正规基定理，若 $L/K$ 是有限 Galois 扩张，则作为 $K[G]$-模有
\[
L\simeq\prod_{g\in G}K=K[G],
\]
其中 $G$ 通过置换坐标作用。这是因为 $K$ 的作用（数乘；$L$ 是 $K$-向量空间）与 $G$ 的作用（$G$ 作为 $K$-向量空间作用于 $L$）交换，故 $L$ 上有 $K[G]$-作用。

\textbf{断言 7.10.} 设 $A$ 是任意交换群，$A[G]:=\prod_{g\in G}A$，$G$ 通过置换坐标作用于此。若 $G$ 循环，则
\[
\HH^0(G,A[G])=\HH^1(G,A[G])=0 .
\]

\textit{证明.} 我们把断言改述如下：设 $R$ 是交换环，则 $R[G]$ 通过乘法运算
\[
\Big(\sum_{g\in G}x_g[g]\Big)\Big(\sum_{h\in G}y_h[h]\Big):=\sum_{g,h\in G}x_gy_h[gh]
\]
成为（可能非交换的）$R$-代数，其中我们用 $[h]\in\prod_{g\in G}R$ 记第 $h$ 个坐标为 $1$、其余坐标为 $0$ 的元素。于是 $R[G]$-模等价于带同态 $G\to\operatorname{Aut}_R(M)$ 的 $R$-模。特别地，$\Z[G]$-模等价于 $G$-模。

现在 $\HH^0(G,A[G])=A[G]^G/N$，其中 $A[G]^G$ 等同于对角嵌入的 $A\subset\prod_{g\in G}A$，而 $N((a,0,\dots,0))=\sum_{g\in G}a[g]$ 等于 $A$ 的对角嵌入，故 $\HH^0(G,A[G])=0$。

同理，$\HH^1(G,A[G])=\ker(N)/(1-\sigma)$，且
\[
A[G]\supseteq\ker(N)=\Big\{\sum_{g\in G}a_g[g]\in A[G]:\sum_{g\in G}a_g=0\Big\}.
\]
现在，一般元素可写成 $\sum_{i=0}^{n-1}a_i[\sigma^i]$，并对每个 $i$ 选取 $b_i$ 使得
\[
(1-\sigma^{n-i})a_i=(1-\sigma)b_i .
\]
于是
\[
(1-\sigma)\sum_{i=0}^{n-1}b_i[\sigma^i]
=\sum_{i=0}^{n-1}(1-\sigma^{n-i})a_i[\sigma^i]
=\sum_{i=0}^{n-1}a_i[\sigma^i]-\sum_{i=0}^{n-1}a_i[1]
=\sum_{i=0}^{n-1}a_i[\sigma^i],
\]
故 $\ker(N)\subset(1-\sigma)A[G]$，从而 $\HH^1(G,A[G])=0$，如所欲证。\hfill$\square$

由此可见，用这种方式得不到有趣的 Tate 上同调。

现在我们回到 $\chi(\OK_L)=1$ 的证明。困难在于正规基定理对 $\OK_L$ 并不适用：虽然 $L=K[G]$，我们却不一定有 $\OK_L\simeq\OK_K[G]$。

不过，$\OK_L$ 确实存在具有正规基的开子群。取 $L/K$ 的一组正规基 $\{e_1,\dots,e_n\}$。对充分大的 $N$，有 $\pi^Ne_1,\dots,\pi^Ne_n\in\OK_L$，其中 $\pi$ 是 $L$ 的素元；于是它们自由张成 $\OK_L$ 的某个开子群。由于这个子群（记作 $\Gamma$）具有有限指数，由 (6.2) 得
\[
\chi(\OK_L)=\chi(\Gamma)=\chi(\OK_K[G])=1 .
\]

为证 $\chi(\OK_L^\times)=1$（更完整的证明将在下一讲给出），注意 $\OK_L^\times\supseteq\Gamma\simeq\OK_L^+$ 是 $G$-等价的，其中 $\Gamma$ 是开子群（这一事实的证明用到 $p$ 进指数映射）。于是 $\chi(\OK_L^\times)=\chi(\Gamma)=1$，如所欲证。\hfill$\square$

\textbf{注 7.11.} 本课程中，所有环与模均假定含幺。

\addcontentsline{toc}{section}{第 8 讲　Tate 上同调与逆向极限}
\section*{第 8 讲\quad Tate 上同调与逆向极限}

回忆：对局部域的扩张 $L/K$、Galois 群 $G:=\Z/n\Z$，我们试图证明 $\#\HH^0(G,L^\times)=n$。我们曾断言 $\chi(L^\times)=n$ 当且仅当 $\chi(\OK_L^\times)=1$，其中 $\chi$ 表示 Herbrand 商 $\#\HH^0/\#\HH^1$；并回忆有限群的 Herbrand 商等于 $1$，且 $\chi$ 对短正合列具有乘性。

上一次我们证明了 $\chi(\OK_L)=1$：用正规基定理说明 $L\simeq K[G]$、$\OK_L$ 含有一个有限指数的开子群 $\Gamma$ 使得 $\OK_L^\times\supset\Gamma\simeq\OK_L^+$ 是 $G$-等价的（从而 $\Gamma$ 在 $G$-作用下封闭），并且 $\OK_K[G]\simeq\Gamma\subseteq\OK_L$。随后我们用断言 7.10 说明对每个 $i$ 都有 $\HH^i(G,\Gamma)=0$，从而 $\chi(\OK_L)=\chi(\Gamma)=1$。

现在我们要给出更好、即更代数化（不用 $p$ 进指数映射！）的 $\chi(\OK_L^\times)=1$ 的证明。为此固定一个同构于 $\OK_K[G]$（作为 $\OK_K[G]$-模）的开子群 $\Gamma\subseteq\OK_L$。

\textbf{断言 8.1.} 对充分大的 $N$，$1+\mathfrak p_K^N\Gamma$ 是 $\OK_L^\times$ 的子群，其中 $\mathfrak p_K$ 是 $\OK_K$ 的极大理想。

\textit{证明.} 对 $x,y\in\Gamma\subseteq\OK_L$，有
\begin{equation}\tag{8.1}
(1+\pi^Nx)(1+\pi^Ny)=1+\underbrace{\pi^N(x+y)}_{\in\,\mathfrak p_K^N\Gamma}+\underbrace{\pi^{2N}(xy)}_{\in\,\mathfrak p_K^{2N}\OK_L},
\end{equation}
其中 $\pi$ 是 $\mathfrak p_K$ 的素元。于是，若取 $N$ 充分大使得 $\mathfrak p_K^N\OK_L\subseteq\Gamma$（这是可能的，因为 $\Gamma$ 是 $\OK_L$ 的开子群），则该乘积落在 $1+\pi^N\Gamma$ 中，从而 $1+\mathfrak p_K^N\Gamma$ 是 $\OK_L^\times$ 的子群。\hfill$\square$

\textbf{断言 8.2.} $\Gamma$ 的各阶上同调全部消没。

取 $N$ 使得 $\mathfrak p_K^N\OK_L\subseteq\mathfrak p_K\Gamma$。则 (8.1) 中最后一项落在 $\mathfrak p_K^{2N}\OK_L\subseteq\mathfrak p_K^{N+1}\Gamma$ 中。这表明我们应当用加法商群来过滤 $1+\mathfrak p_K^N\Gamma$，即用 $1+\mathfrak p_K^{N+i}\Gamma$ 过滤，使得由上述计算，对一切 $i\ge 0$ 有加法群的同构
\[
(1+\mathfrak p_K^{N+i}\Gamma)/(1+\mathfrak p_K^{N+i+1}\Gamma)\simeq\Gamma/\mathfrak p_K\Gamma\simeq k_K[G],
\]
其中 $k_K$ 表示 $K$ 的剩余域。此外，这些同构是 Galois 等变的，即 $G$-等变的：$G$-作用保持所有项（$\Gamma$ 由假设保持），从而作用于两边，且被同构保持。于是由断言 7.10，
\[
\HH^j\big(G,(1+\mathfrak p_K^{N+i}\Gamma)/(1+\mathfrak p_K^{N+i+1}\Gamma)\big)=\HH^j(G,k_K[G])=0
\]
对每个 $i\ge 0$ 与 $j$ 成立。作为推论（为此需要下面的引理），
\begin{equation}\tag{8.2}
\HH^j\big(G,(1+\mathfrak p_K^N\Gamma)/(1+\mathfrak p_K^{N+i}\Gamma)\big)=0
\end{equation}
对一切 $i\ge 0$ 与 $j$ 成立。

\textbf{引理 8.3.} 对任意 $G$-模的短正合列
\[
0\to M\to E\to N\to 0,
\]
若 $\HH^i(G,M)=\HH^i(G,N)=0$，则对每个 $i$ 有 $\HH^i(G,E)=0$。

\textit{证明.} 由 (6.4) 有正合列
\[
\HH^i(G,M)\xrightarrow{\alpha}\HH^i(G,E)\xrightarrow{\beta}\HH^i(G,N),
\]
故 $\HH^i(E)=\ker(\beta)=\im(\alpha)=0$，如所欲证。\hfill$\square$

现在有正合列
\[
0\to\frac{1+\mathfrak p_K^{N+i}\Gamma}{1+\mathfrak p_K^{N+i+1}\Gamma}\to\frac{1+\mathfrak p_K^{N}\Gamma}{1+\mathfrak p_K^{N+i+1}\Gamma}\to\frac{1+\mathfrak p_K^{N}\Gamma}{1+\mathfrak p_K^{N+i}\Gamma}\to 0,
\]
于是 (8.2) 可由对 $i$ 的归纳与引理 8.3 得到。

剩下的问题是证明 $\HH^j(1+\mathfrak p_K^N\Gamma)=0$。在理想的情形下，对任意带 $G$-作用与 $G$-等变结构映射的模序列，我们都应有
\[
\HH^j\Big(G,\varprojlim_n M_n\Big)=\varprojlim_n\HH^j(G,M_n).
\]
这样便可推出
\[
\HH^j(1+\mathfrak p_K^N\Gamma)
=\HH^j\Big(\varprojlim_{i\ge 0}(1+\mathfrak p_K^N\Gamma)/(1+\mathfrak p_K^{N+i}\Gamma)\Big)
=\varprojlim_{i\ge 0}\HH^j\big((1+\mathfrak p_K^N\Gamma)/(1+\mathfrak p_K^{N+i}\Gamma)\big)
=\varprojlim_{i\ge 0}0=0
\]
（由 (8.2)，因为我们的过滤是完备的）。因此，我们需要设法说明 Tate 上同调与逆向极限可以交换。

\textbf{引理 8.4.} 设有一列模
\begin{equation}\tag{8.3}
\begin{CD}
\vdots @. \vdots @. \vdots @. @.\\
0 @>>> M_{n+1} @>>> E_{n+1} @>>> N_{n+1} @>>> 0\\
@. @VVV @VVV @VVV @.\\
0 @>>> M_{n} @>>> E_{n} @>>> N_{n} @>>> 0\\
\vdots @. \vdots @. \vdots @. @.
\end{CD}
\end{equation}
其各行正合。则我们有正合列
\[
0\to\varprojlim_n M_n\xrightarrow{\varphi}\varprojlim_n E_n\xrightarrow{\psi}\varprojlim_n N_n .
\]
% [原文疑误：原文作 "if $M_n\to M_{n+1}$ is surjective for each $n$"，与 (8.3) 阶梯图的方向相反；译文按图示取 $M_{n+1}\to M_n$]
此外，若对每个 $n$，$M_{n+1}\to M_n$ 都是满射，则 $\psi$ 也是满射（否则可能不是！）。

\textit{证明.} 显然 $\varphi$ 是单射：若 $x\in\ker(\varphi)$，则其像的每个坐标都是 $0$，故由每个 $M_n\to E_n$ 的单射性与相容性，$x$ 的每个坐标都是 $0$，从而 $x=0$。同理，由 (8.3) 每行正合可得 $\ker(\psi)=\im(\varphi)$。

为（直观地）看出 $\psi$ 的满射性何以会失败，考虑相容系统 $(x_n)\in\varprojlim_n N_n$。我们可以把每个 $x_n$ 提升为某个 $y_n\in E_n$，但如何相容地做到这一点、使 $(y_n)\in\varprojlim_n E_n$，并不清楚。

现在假设每个映射 $M_{n+1}\to M_n$ 都是满射。取 $(x_n)\in\varprojlim_n N_n$，并设对某个 $n$ 已构造出 $y_n\in E_n$。任取 $\tilde y_{n+1}\mapsto x_{n+1}$（经 $E_{n+1}\to N_{n+1}$），令 $\alpha_{n+1}$ 为 $\tilde y_{n+1}$ 在 $E_n$ 中的像。则 $y_n-\alpha_{n+1}\in M_n$（因为它在 $N_n$ 中消没），且由假设它可提升为 $\beta_{n+1}\in M_{n+1}$。若令 $y_{n+1}:=\tilde y_{n+1}+\beta_{n+1}$，则它在 $E_n$ 中映到 $\alpha_{n+1}+y_n-\alpha_{n+1}=y_n$，在 $N_{n+1}$ 中映到 $x_{n+1}$（因为 $\beta_{n+1}$ 在 $N_{n+1}$ 中映到 $0$）；于是由归纳法，存在相容系统 $(y_n)\in\varprojlim_n E_n$ 经 $\psi$ 映到 $(x_n)$，如所欲证（注意，这一结果可表述为满射 $E_{n+1}\twoheadrightarrow E_n\times_{N_n}N_{n+1}$，即到纤维积）。\hfill$\square$

\textbf{命题 8.5.} 若
\begin{equation}\tag{8.4}
\cdots\twoheadrightarrow M_{n+1}\twoheadrightarrow M_n\twoheadrightarrow\cdots\twoheadrightarrow M_0
\end{equation}
是 $G$-模的序列（各映射均为满射；证明中要用到这一点，见下文），且对一切 $n$ 与 $i$ 有 $\HH^i(M_{n+1})\twoheadrightarrow\HH^i(M_n)$，则对一切 $i$ 有
\[
\HH^i\Big(\varprojlim_n M_n\Big)=\varprojlim_n\HH^i(M_n).
\]

\textit{证明.} 我们对 $\HH^0$ 给出证明。令 $M:=\varprojlim_n M_n$，于是我们要比较 $\HH^0(M)=M^G/N(M)$ 与 $\varprojlim_n\HH^0(M_n)=\varprojlim_n M_n^G/N(M_n)$。这归结为证明自然映射
\[
\Big(\varprojlim_n M_n^G\Big)\Big/\Big(\varprojlim_n N(M_n)\Big)\xrightarrow{\sim}\varprojlim_n\big(M_n^G/N(M_n)\big)
\]
是同构。我们有交换图
\[
\begin{CD}
0 @>>> N(M_{n+1}) @>>> M_{n+1}^G @>>> \HH^0(G,M_{n+1}) @>>> 0\\
@. @V{\alpha_n}VV @VVV @VVV @.\\
0 @>>> N(M_{n}) @>>> M_{n}^G @>>> \HH^0(G,M_{n}) @>>> 0,
\end{CD}
\]
并断言每个 $\alpha_n$ 都是满射。事实上，取 $x\in N(M_n)$，则对某个 $y\in M_n$ 有 $x=N(y)$。把 $y$ 提升为元素 $z\in M_{n+1}$，便有 $\alpha_n(N(z))=x$，如所欲证。于是由引理 8.4，
\[
\varprojlim_n\HH^0(M_n)=\Big(\varprojlim_n M_n^G\Big)\Big/\Big(\varprojlim_n N(M_n)\Big).
\]
现在
\[
\HH^0\Big(\varprojlim_n M_n\Big)=\Big(\varprojlim_n M_n\Big)^G\Big/N\Big(\varprojlim_n M_n\Big)=M^G/N(M).
\]
由于 $G$ 作用于 $\varprojlim_n M_n$ 的每个坐标，显然 $(\varprojlim_n M_n)^G=\varprojlim_n M_n^G$；于是只剩下证明 $N(\varprojlim_n M_n)\xrightarrow{\sim}\varprojlim_n N(M_n)$。对每个 $n$ 令 $K_n:=\ker(N:M_n\to M_n)$，则有交换图
\[
\begin{CD}
0 @>>> K_{n+1} @>>> M_{n+1} @>>> N(M_{n+1}) @>>> 0\\
@. @VVV @VVV @V{\beta_n}VV @.\\
0 @>>> K_{n} @>>> M_{n} @>>> N(M_{n}) @>>> 0.
\end{CD}
\]
我们希望证明 $\beta_n$ 是满射。回忆 $\HH^1(G,M_n)=K_n/(1-\sigma)M_n$，于是有交换图
\[
\begin{CD}
0 @>>> (1-\sigma)M_{n+1} @>>> K_{n+1} @>>> \HH^1(G,M_{n+1}) @>>> 0\\
@. @V{\gamma_n}VV @V{\beta_n}VV @V{\delta_n}VV @.\\
0 @>>> (1-\sigma)M_{n} @>>> K_{n} @>>> \HH^1(G,M_{n}) @>>> 0.
\end{CD}
\]
现在 $\gamma_n$ 是满射（证明与 $\alpha_n$ 的情形类似），而 $\delta_n$ 由假设是满射。于是由蛇引理，$\beta_n$ 是满射，从而引理 8.4 给出
\[
\varprojlim_n N(M_n)=\Big(\varprojlim_n M_n\Big)\Big/\Big(\varprojlim_n K_n\Big)=N\Big(\varprojlim_n M_n\Big).
\]
由此即得 $\HH^0(\varprojlim_n M_n)=\varprojlim_n\HH^0(M_n)$，如所欲证。\hfill$\square$

\textbf{推论 8.6.} 对序列 (8.4)，若对一切 $n$ 与 $i$ 有 $\HH^i(M_n)=0$，则
\[
\HH^i\Big(\varprojlim_n M_n\Big)=0 .
\]
特别地，$\HH^i(1+\mathfrak p_K^N\Gamma)=0$，其中我们置
\[
M_i:=(1+\mathfrak p_K^N\Gamma)/(1+\mathfrak p_K^{N+i}\Gamma)\qquad(i\ge 0).
\]
于是，由于 $1+\mathfrak p_K^N\Gamma\subseteq\OK_L^\times$ 是有限指数子群（作为加法群带有正规基），我们有
\[
\chi(\OK_L^\times)=\chi(1+\mathfrak p_K^N\Gamma)=1,
\]
这就建立了断言 7.8 的 (2)。

\addcontentsline{toc}{section}{第 9 讲　希尔伯特定理 90 与上链复形}
\section*{第 9 讲\quad 希尔伯特定理 90 与上链复形}

一如既往，$G=\Z/n\Z$，$L/K$ 是局部域的 Galois 扩张，$\Gal(L/K)=G$，生成元为 $\sigma\in G$。上一讲我们证明了：

\textbf{定理 9.1.} $\chi(L^\times)=n$，其中 $\chi$ 表示 Herbrand 商。

注意，我们的方法其实很容易推广到非阿基米德的情形。本讲我们将证明：

\textbf{定理 9.2（希尔伯特定理 90）.} $\HH^1(G,L^\times)=0$。\footnote{原文此处印作 $\HH^1(G,K^\times)=0$；由下文 $\HH^0(G,L^\times)=K^\times/N(L^\times)$、以及推论 9.3 给出的显式陈述可知应为 $\HH^1(G,L^\times)=0$，此处径改。}

两者合起来便推出 $\HH^0(G,L^\times)=K^\times/N(L^\times)$ 的基数为 $n$。希尔伯特定理 90 的另一个推论是：若 $L$ 与 $K$ 都是有限域，则对一切 $i$ 有 $\HH^i(G,L^\times)=0$，因为 $L$ 有限故 $\chi(L^\times)=1$（于是由周期性所有上同调都消没）。这与 [Wei74] 中的第一个结果类似。具体地，
\[
\HH^1(G,L^\times)=\ker(N:L^\times\to K^\times)\big/\{y/\sigma y:y\in L^\times\},
\]
其中每个 $y/\sigma y$ 的范数都是 $1$，因为 $y/\sigma y\cdot\sigma y/\sigma^2y\cdots\sigma^{n-1}y/\sigma^ny=1$。于是希尔伯特定理 90 蕴含：

\textbf{推论 9.3.} 若 $L/K$ 是循环扩张，$x\in L^\times$ 且 $N(x)=1$，则存在 $y\in L^\times$ 使得 $x=y/\sigma y$。

\textbf{例 9.4.} 取 $L:=\Q(i)$、$K:=\Q$。选 $x\in\Q(i)$ 使 $N(x)=1$。则对某些满足 $a^2+b^2=c^2$ 的 $a,b,c\in\Z$ 有 $x=a/c+(b/c)i$。于是希尔伯特定理 90 给出勾股数组的通常参数化：
\[
(r-s)^2+(2rs)^2=(r+s)^2 .
\]

对 $n=2$，证明很简单。此时 $N(x)=x\cdot\sigma x=1$。若 $x\neq-1$，令 $y:=x+1$，则
\[
x\cdot\sigma y=x(\sigma x+1)=N(x)+x=1+x=y,
\]
故 $x=y/\sigma y$，如所欲证。若 $x=-1$，令 $y:=\sqrt d$，其中 $L=K(\sqrt d)$，则同样有 $y/\sigma y=\sqrt d/(-\sqrt d)=-1=x$。注意，这就完成了对局部域的二次扩张 $L/K$ 证明 $\#(K^\times/NL^\times)=2$，从而完成了希尔伯特符号良好性质的证明！事实上，回忆对域 $L:=K(\sqrt a)$（$a\in K^\times$ 但非平方元），$(a,b)=1$ 当且仅当 $b\in N(L^\times)$。

现在我们转向希尔伯特定理 90 的一般情形。主要引理如下：

\textbf{引理 9.5.} 对每个 $x\in L$，令
\[
H_x:L\to L,\qquad y\mapsto x\cdot\sigma(y),
\]
它是 $K$-向量空间的线性映射。则 $H_x$ 的特征多项式为 $t^n-N(x)\in K[t]$，其中我们把特征多项式的定义规范化为首一的。

注意，这个特征多项式比 $y\mapsto xy$ 的特征多项式简单：只要 $T(x)\neq 0$，后者就会含有 $t^n-1$ 的非零倍，而迹非退化（任何可分扩张都如此）时就会出现这种情况。

\textit{证明（由 9.5 推出 9.2）.} 取 $x\in L$ 并设 $N(x)=1$。则 $H_x$ 的特征多项式为 $t^n-1$，故 $1$ 是它的根，从而是 $H_x$ 的特征值。于是 $\ker((H_x-1)\otimes_K\bar K)\neq 0$；而对域而言张量积与取核交换，故 $\ker(H_x-1)\neq 0$。于是存在 $y\in L^\times$ 使得 $H_x(y)=x\cdot\sigma(y)=y$，即 $x=y/\sigma y$，如所欲证。\hfill$\square$

\textit{证明（引理 9.5）.} 首先注意 $H_x^n$ 对应乘以 $N(x)$，因为对任意 $y\in L$，
\[
H_x^n(y)=x\cdot\sigma(x\cdot\sigma(x\cdots\sigma(y)))=x\cdot\sigma(x)\cdot\sigma^2(x)\cdots\sigma^{n-1}(x)\cdot\sigma^n(y)=N(x)y .
\]
于是 $H_x$ 的极小多项式整除 $t^n-N(x)$。现在回忆：线性算子 $T$ 的极小多项式总整除其特征多项式，而后者次数为 $n$，故只需证明两者相等。这等价于 $T$ 的若尔当分解中没有共享特征值的块，而这又等价于对一切 $\lambda\in K$ 有 $\dim_K(\ker(T-\lambda I))\le 1$。

下面给出一个并不完全成立的证明。设 $H_x(y_1)=\lambda y_1$、$H_x(y_2)=\lambda y_2$，且 $y_1,y_2\neq 0$（即两者都是「真正的」特征向量）。我们想证明 $y_2$ 是 $y_1$ 的倍数，即 $y_2/y_1\in K$，也就是被 $\Gal(L/K)$ 固定。事实上，由 $\sigma y_2=H_x(y_2)/x$（对 $y_1$ 同理）可得
\[
\frac{\sigma y_2}{\sigma y_1}=\frac{x\lambda y_2}{x\lambda y_1}=\frac{y_2}{y_1}.
\]
然而问题在于：这个证明发生在 $L$ 中，而不在 $K$ 中——而我们的特征值实际住在 $K$ 里！因此，我们需要在 $L\otimes_K\bar K\simeq\prod_{g\in G}\bar K$ 中工作，而它不一定是域。

我们可以在基变换之后再计算特征多项式。回忆
\[
L\otimes_K\bar K\xrightarrow{\sim}\prod_{i=0}^{n-1}\bar K,\qquad a\otimes b\mapsto\big((\sigma^ia)\cdot b\big)_{i=0}^{n-1}.
\]
这个同构非典范地延拓为 $\bar K$ 的一个自同构，除此之外一切都是典范的，其中群通过左乘作用于坐标集合。映射
\[
\sigma\otimes\mathrm{id}:L\otimes_K\bar K\to L\otimes_K\bar K
\]
对应置换坐标；此外还有映射
\[
\mu_x\otimes\mathrm{id}:L\otimes_K\bar K\to L\otimes_K\bar K,\qquad
(y_0,\dots,y_{n-1})\mapsto\big(xy_0,(\sigma x)y_1,\dots,(\sigma^{n-1}x)y_{n-1}\big),
\]
其中 $\mu_x$ 表示乘以 $x$。现在设 $\lambda\in\bar K$ 是 $H_x$ 的特征值，对应的特征向量为 $(y_0,\dots,y_{n-1})$。则
\[
H_x(y)=\big(xy_1,(\sigma x)y_2,(\sigma^2x)y_3,\dots,(\sigma^{n-1}x)y_0\big)
=(\lambda y_0,\lambda y_1,\lambda y_2,\dots,\lambda y_{n-1}),
\]
于是 $xy_1=\lambda y_0$，即 $y_1=(\lambda/x)y_0$；同理 $y_2=(\lambda/\sigma x)y_1=(\lambda^2/(x\cdot\sigma x))y_0$。一般地有
\[
y_i=\frac{\lambda^i}{x\cdots\sigma^{i-1}x}\,y_0,
\]
故所有坐标都由 $y_0$ 唯一确定，即
\[
(y_0,\dots,y_{n-1})=y_0\Big(1,\frac{\lambda}{x},\frac{\lambda^2}{x\cdot\sigma x},\dots,\frac{\lambda^{n-1}}{\prod_{i=0}^{n-2}\sigma^ix}\Big).
\]
因此每个特征空间确实都是一维的，如所欲证。注意，只有当
\[
\frac{\lambda^n}{x\cdots\sigma^{n-1}x}=\frac{\lambda^n}{N(x)}=1,
\]
即 $\lambda^n=N(x)$ 时，这才真正定义出特征向量，这与我们的预期一致（并且 $n$ 个 $n$ 次根都以重数一出现）。\hfill$\square$

现在回忆我们的目标是：对局部域的交换扩张 $L/K$，典范地（且是在很强的意义下）证明
\[
K^\times/NL^\times\simeq\Gal(L/K).
\]
我们已经证明 $K^\times/NL^\times$ 的阶正确，但将用上同调对非循环群一般地证明它。下面我们引入同调代数的语言，这将是我们的核心工具。

\textbf{定义 9.6.} 交换群的（上链）复形 $X$ 是一列
\[
\cdots\to X^{-1}\xrightarrow{d^{-1}}X^0\xrightarrow{d^0}X^1\xrightarrow{d^1}\cdots,
\]
使得微分满足对每个 $i$ 有 $d^{i+1}d^i=0$。

\textbf{记号 9.7.} 我们常把整个复形记作 $X^\bullet$，其中「$\bullet$」处于指标的位置。我们也常省略指标，例如把 $d^i$ 写成 $d$，并写 $d\cdot d=d^2=0$。注意，有些作者写作 $H_i:=H^{-i}$，对 $X^i$ 也类似，这样微分降低次数；而我们的约定是微分提升次数。

\textbf{定义 9.8.} 第 $i$ 个上同调群为 $H^i(X):=\ker(d^i)/\im(d^{i-1})$。

这些上同调群其实正是我们要的不变量；但 $X$ 是更「丰富」的对象，因此最好在我们过程的最后一步才过渡到上同调。下面引入复形映射的零伦这一重要想法。

\textbf{定义 9.9.} 若下图交换
\[
\begin{CD}
\cdots @>>> X_{-1} @>d_{-1}>> X_0 @>d_0>> X_1 @>>> \cdots\\
@. @V{f_{-1}}VV @V{f_0}VV @V{f_1}VV @.\\
\cdots @>>> Y_{-1} @>d_{-1}>> Y_0 @>d_0>> Y_1 @>>> \cdots
\end{CD}
\]
则称 $f$ 为\textbf{复形映射}。注意，$f$ 诱导上同调之间的映射，因为由交换性，$X_\bullet$ 中微分的核与像都被保持到 $Y_\bullet$ 中。若下图中的映射 $h$
\[
\begin{CD}
\cdots @>>> X_{-1} @>d_{-1}>> X_0 @>d_0>> X_1 @>>> \cdots\\
@. @V{f_{-1}}VV @V{f_0}VV @V{f_1}VV @.\\
\cdots @>>> Y_{-1} @>d_{-1}>> Y_0 @>d_0>> Y_1 @>>> \cdots
\end{CD}
\]
（图中另有斜向连接相邻两行的映射 $h_0,h_1$，此处从略）满足 $dh+hd=f$，更精确地说，对每个 $i$ 有
\[
% [原文疑误：原印 $d_ih_{i+1}+h_id_{i+1}=f_{i+1}$ 指标位移有误，与引理 9.10 的证明不自洽；已按标准约定改写，其中 $h_i:X_i\to Y_{i-1}$]
f_i=d_{i-1}h_i+h_{i+1}d_i
\]
则称 $h$ 为 $f$ 的\textbf{零伦}。

\textbf{引理 9.10.} 若 $f$ 是零伦的，则其诱导的上同调映射 $H^i(X)\xrightarrow{H^i(f)}H^i(Y)$ 对一切 $i$ 为零。

\textit{证明.} 取 $x\in X_i$ 使得 $d_ix=0$。则 $f(x)=d_{i-1}h_i(x)+h_{i+1}d_ix=d(h_i(x))$，故 $f(x)\in\im(d_{i-1})$，从而 $f(x)=0$ 于 $H^i(Y)$。\hfill$\square$

现在，我们在此处的指导原则是：对代数而言，同构是比相等好得多的概念，因为相等指的是没有结构的集合。因此，若 $f\simeq g$，即 $f$ 同伦于 $g$（意思是 $f-g$ 存在零伦），则任何真正的数学检验都无法再区分 $f$ 与 $g$。

我们希望为复形映射定义某种「余核」的概念。一个糟糕的想法是：对复形映射 $f:X\to Y$ 直接取 $\coker(f)$。更好的想法如下：

\textbf{定义 9.11.} \textbf{同伦余核}或\textbf{锥} $h\coker(f)=\operatorname{Cone}(f)$ 具有如下泛性质：复形映射 $h\coker(f)\to Z$ 等价于映射 $Y\to Z$ 连同 $X\to Z$ 的一个零伦的数据；注意这给出如下交换图
\[
\begin{CD}
X @>>> Z\\
@V{f}VV @.\\
Y @.
\end{CD}
\]
注意它与通常余核的泛性质之间的相似性。

\addcontentsline{toc}{section}{第 10 讲　同伦、拟同构与余不变商}
% 第 10 讲 同伦、拟同构与余不变商
% 源：src/lecture10.txt（199 行）
\section*{第 10 讲\quad 同伦、拟同构与余不变商}

请注意，本讲中许多断言的证明留作习题集 5。

回忆：带微分 $d$ 的阿贝尔群序列若满足 $d^2=0$，就称为\textbf{复形}；若 $df=fd$，则 $f\colon X\to Y$ 称为链复形的\textbf{态射}；若 $dh+hd=f$，则 $h$ 称为（$f$ 的）\textbf{零伦}。下图示意这一情形：
\[
\begin{array}{ccccccccc}
\cdots & \longrightarrow & X_{-1} & \xrightarrow{\ d\ } & X_0 & \xrightarrow{\ d\ } & X_1 & \longrightarrow & \cdots\\[5pt]
& & {\scriptstyle f}\big\downarrow & & {\scriptstyle f}\big\downarrow & & {\scriptstyle f}\big\downarrow & & \\[5pt]
\cdots & \longrightarrow & Y_{-1} & \xrightarrow{\ d\ } & Y_0 & \xrightarrow{\ d\ } & Y_1 & \longrightarrow & \cdots
\end{array}
\]
其中竖直箭头为 $f$；零伦 $h$ 由一组斜箭头（从 $X_i$ 指向 $Y_{i-1}$）给出，条件 $dh+hd=f$ 说的正是：「先沿上排走 $d$、再沿 $h$ 下行」与「先沿 $h$ 下行、再沿下排走 $d$」二者之和等于 $f$（原图用 $h$ 标出这些斜箭头，此处从略）。

链复形的不变量是同调群
\[
H_i(X):=\ker(d\colon X_i\to X_{i+1})/\operatorname{Im}(d\colon X_{i-1}\to X_i),
\]
而对 $f,g\colon X\to Y$，我们称 $f\simeq g$（即 $f$ 与 $g$ \textbf{同伦}），当且仅当 $f-g$ 存在零伦；由引理 9.10，这迫使 $f$ 与 $g$ 在上同调上诱导相同的映射。

设 $G$ 为有限群，$L/K$ 为局部域扩张且 $G=\Gal(L/K)$。按定义有 $\HH^0(G,L^\times)=K^\times/NL^\times$。我们的目标是证明 $\HH^0(G,L^\times)\cong G^{\mathrm{ab}}$（典范地），即 $G$ 的交换化。本讲的计划是：对每个 $i\in\Z$ 定义 Tate 上同调群 $\HH^i$（对非循环群要复杂一些），然后用它们开始走向这一事实的证明。

回忆我们的基本原则：给定同伦 $h\colon f\simeq g$，则 $f$ 与 $g$ 在实践中已无法区分（这一点我们直接采信）。这一原则的一个应用是锥（或称同伦余核）的构造：

\textbf{断言 10.1.} 若 $f\colon X\to Y$ 是复形的映射，则同伦余核 $\operatorname{hCoker}(f)$（又称锥 $\operatorname{Cone}(f)$）存在；刻画它的万有性质是：链复形映射 $\operatorname{hCoker}(f)\to Z$ 等价于映射 $g\colon Y\to Z$ 连同 $g\circ f\colon X\to Z$ 的一个零伦 $h$。

\textit{证明.} 我们断言下列链复形就是 $\operatorname{hCoker}(f)$：
\[
(10.1)\qquad \cdots\to X_0\oplus Y_{-1}\to X_1\oplus Y_0\to X_2\oplus Y_1\to\cdots
\]
其微分为
\[
X_i\oplus Y_{i-1}\ni\begin{pmatrix}x\\ y\end{pmatrix}
\xrightarrow{\ d\ }
\begin{pmatrix}-dx\\ f(x)+dy\end{pmatrix}\in X_{i+1}\oplus Y_i,
\]
注意其次数恰好增加 $1$。可将该微分写成矩阵 $\begin{pmatrix}-d&0\\ f&d\end{pmatrix}$；它平方为零：
\[
\begin{pmatrix}-d&0\\ f&d\end{pmatrix}\begin{pmatrix}-d&0\\ f&d\end{pmatrix}
=\begin{pmatrix}d^2&0\\ -fd+df&d^2\end{pmatrix}
=\begin{pmatrix}0&0\\0&0\end{pmatrix},
\]
这里用了链复形态射的定义（$df=fd$）以及 $X$、$Y$ 均为复形这一点。

现在验证该链复形满足 $\operatorname{hCoker}(f)$ 的万有性质。设有映射 $\operatorname{hCoker}(f)\to Z$，即下图交换：
\[
\begin{array}{ccccccc}
\cdots & \longrightarrow & X_{i+1}\oplus Y_i & \longrightarrow & X_{i+2}\oplus Y_{i+1} & \longrightarrow & \cdots\\
& & \big\downarrow & & \big\downarrow & & \\
\cdots & \longrightarrow & Z_i & \longrightarrow & Z_{i+1} & \longrightarrow & \cdots
\end{array}
\]
若把该映射写成 $(x,y)\mapsto h(x)+g(y)$，则上述交换性意味着
\[
dh(x)+dg(y)=d(h(x)+g(y))=h(-dx)+g(f(x)+dy)=-h(dx)+gf(x)+g(dy).
\]
取 $x=0$ 得 $dg=gd$；于是必须有 $dh+hd=g\circ f$，即 $h$ 是 $g\circ f$ 的零伦，如所欲证。\hfill$\square$

\textbf{推论 10.2.} 复合 $X\to Y\to\operatorname{hCoker}(f)$ 典范地零伦（作为练习，请把这个零伦显式构造出来！）。

\textbf{例 10.3.} 设
\[
X:=(\cdots\to0\to A\to0\to\cdots),\qquad Y:=(\cdots\to0\to B\to0\to\cdots),
\]
其中 $A,B$ 是位于次数 $0$ 的有限阿贝尔群，并设 $f\colon A\to B$。则
\[
\operatorname{hCoker}(f)=(\cdots\to0\to A\xrightarrow{\ f\ }B\to0\to\cdots),
\]
$B$ 位于次数 $0$。于是
\[
H_0\operatorname{hCoker}(f)=\operatorname{Coker}(f),\qquad H_{-1}\operatorname{hCoker}(f)=\ker(f),
\]
可见链复形的语言推广了先前的概念。

\textbf{记号 10.4.} 对链复形 $X$，记 $X[n]$ 为 $X$ 平移 $n$ 位所得的链复形，即第 $i$ 次为 $X_{i+n}$、微分为 $(-1)^nd$（$d$ 仍记 $X$ 的微分）的链复形。例如 $X[1]=\operatorname{hCoker}(X\to0)$。其含义是：给出 $0\colon X\to Y$ 的一个零伦，等价于给出映射 $X[1]\to Y$。

\textbf{引理 10.5.} 对任意映射 $f\colon X\to Y$，序列
\[
H_iX\to H_iY\to H_i\operatorname{hCoker}(f)
\]
对一切 $i$ 正合。

\textit{证明.} 复合为零，因为 $X\to Y\to\operatorname{hCoker}(f)$ 零伦（引理 9.10）。为证正合性，取 $y\in Y_i$ 使 $dy=0$，并设它在 $H_i\operatorname{hCoker}(f)$ 中的像为零，即
\[
\begin{pmatrix}0\\ y\end{pmatrix}
=\begin{pmatrix}-d&0\\ f&d\end{pmatrix}\begin{pmatrix}\alpha\\ \beta\end{pmatrix}
=\begin{pmatrix}-d\alpha\\ f(\alpha)+d\beta\end{pmatrix}
\]
其中 $\alpha\in X_i$、$d\alpha=0$，$\beta\in Y_{i-1}$。于是 $f(\alpha)+d\beta=y$ 给出 $f(\alpha)=y$（在 $H_iY$ 中），如所欲证。\hfill$\square$

\textbf{断言 10.6.} 还有同伦核 $\operatorname{hKer}(f)$ 的概念，由下述万有性质刻画：映射 $Z\to\operatorname{hKer}(f)$ 等价于映射 $Z\to X$ 连同复合 $Z\to X\to Y$ 的一个零伦。特别地，$\operatorname{hKer}(f)=\operatorname{hCoker}(f)[-1]$。

\textbf{例 10.7.} 设 $f\colon A\to B$ 是阿贝尔群的映射（仍位于次数 $0$）。则
\[
\begin{aligned}
\operatorname{hCoker}(f)&=(\cdots\to0\to A\xrightarrow{\ f\ }B\to0\to0\to\cdots),\\
\operatorname{hKer}(f)&=(\cdots\to0\to0\to A\xrightarrow{\ f\ }B\to0\to\cdots),
\end{aligned}
\]
其中 $\operatorname{hKer}(f)_0=A$。同伦余核也在其同调中还原出核与余核。

\textbf{断言 10.8.} 复合
\[
X\xrightarrow{\ f\ }Y\to\operatorname{hCoker}(f)
\]
零伦，故存在典范映射
\[
X\to\operatorname{hKer}\bigl(Y\to\operatorname{hCoker}(f)\bigr),
\]
后者称为「映射柱」。该映射是同伦等价。

\textbf{定义 10.9.} 映射 $f\colon X\to Y$ 称为\textbf{同伦等价}，若存在映射 $g\colon Y\to X$ 以及同伦 $gf\simeq\operatorname{id}_X$、$fg\simeq\operatorname{id}_Y$，此时记 $X\simeq Y$。若 $H_i(f)\colon H_i(X)\xrightarrow{\sim}H_i(Y)$ 对每个 $i$ 都是同构（即 $X$ 与 $Y$ 在上同调层面相等），则称 $f$ 为\textbf{拟同构}。

\textbf{断言 10.10.} 若 $f\colon X\to Y$ 是同伦等价，则它是拟同构。

\textit{证明.} 这是引理 9.10 的直接推论：该引理保证 $f$ 与 $g$ 在上同调层面互为逆。\hfill$\square$

\textbf{推论 10.11.} 给定 $f\colon X\to Y$，有长正合列
\[
\cdots\to H_{i-1}\operatorname{hCoker}(f)\to H_iX\to H_iY\to H_i\operatorname{hCoker}(f)\to H_{i+1}X\to\cdots.
\]

\textit{证明.} 记 $g$ 为映射 $Y\to\operatorname{hCoker}(f)$，则复合
\[
Y\xrightarrow{\ g\ }\operatorname{hCoker}(f)\to\operatorname{hCoker}(g)=\operatorname{hKer}(g)[1]\simeq X[1]
\]
由推论 10.2 零伦，而同伦等价由断言 10.8 给出。故由引理 10.5，序列
\[
H_iY\to H_i\operatorname{hCoker}(f)\to H_iX[1]=H_{i+1}X
\]
正合；再用一次引理 10.5 即得结论。\hfill$\square$

\textbf{断言 10.12.} 设 $f_i\colon X_i\to Y_i$ 对每个 $i$ 都是单射。则 $\operatorname{hCoker}(f)\to Y/X$（即第 $i$ 次为 $Y_i/X_i$ 的复形）是拟同构。

\textbf{例 10.13.} 若 $f\colon A\hookrightarrow B$ 是位于次数 $0$ 的阿贝尔群映射，则映射 $\operatorname{hCoker}(f)\to B/A$ 形如
\[
\begin{array}{ccccccccc}
\cdots & \longrightarrow & A & \longrightarrow & B & \longrightarrow & 0 & \longrightarrow & \cdots\\
& & \big\downarrow & & \big\downarrow & & \big\downarrow & & \\
\cdots & \longrightarrow & 0 & \longrightarrow & B/A & \longrightarrow & 0 & \longrightarrow & \cdots
\end{array}
\]
容易看出这确实是拟同构。注意有对偶的论断：若每个 $f_i$ 都是满射，则同伦核拟同构于朴素核。

\textbf{注 10.14.} 若 $A$ 是结合代数（例如 $\Z$ 或 $\Z[G]$），则可考虑 $A$-模的链复形
\[
\cdots\to X_{-1}\xrightarrow{\ d\ }X_0\xrightarrow{\ d\ }X_1\to\cdots,
\]
其中 $X_i$ 是 $A$-模、$d$ 是 $A$-模映射。此时上同调也自动是 $A$-模。

现在回到原问题：对有限群 $G$ 在某个 $A$ 上的作用定义 Tate 上同调。注意
\[
\HH^0(G,A)=A^G/N(A)=\operatorname{Coker}(N\colon A\to A^G).
\]
事实上，我们可以做得比 $N\colon A\to A^G$ 更好：范映射经过我们将称为余不变商的模分解。

\textbf{定义 10.15.} $A$ 的\textbf{余不变商}为 $A_G:=A/\sum_{g\in G}(g-1)A$；它满足万有性质：$A_G$ 是满足「对一切 $x\in A$、$g\in G$ 都有 $gx=x$」的 $A$ 的最大商模。

注意可把不变子模 $A^G$ 看成各 $(g-1)A$ 的核的交，故它同样是对一切 $g$ 满足 $gx=x$ 的 $A$ 的最大子模。于是范映射分解为 $A\to A_G\xrightarrow{N}A^G$，即
\[
\begin{array}{ccc}
A & \xrightarrow{\ N\ } & A^G\\
& \searrow{\scriptstyle N} & \big\uparrow\\[2pt]
& A_G &
\end{array}
\]

现在的计划是为 $A^G$ 与 $A_G$ 定义导出（复形）版本，记作 $A^{hG}$ 与 $A_{hG}$，并考虑范映射 $N\colon A_{hG}\to A^{hG}$；Tate 上同调将是此映射的同伦余核。基本观察是：$\Z$ 以平凡方式成为 $G$-模（即 $\Z[G]$ 通过增广作用在 $\Z$ 上），$G$ 的每个元素都作为恒等自同构作用。若 $M$ 是 $G$-模，则 $M^G=\operatorname{Hom}_G(\Z,M)$（因为 $1\in\Z$ 的像必须是 $G$-不变的，且与 $M^G$ 的元素一一对应），并且 $M_G=M\otimes_{\Z[G]}\Z$。事实上，设 $I\subseteq A$ 是作用在 $M$ 上的理想，则由张量积的右正合性有 $A/I\otimes_AM=M/IM$。这里 $\Z=\Z[G]/I$，其中 $I$ 是由元素 $g-1$ 生成的「增广理想」，因此 $M_G=M/IM$，如所欲证。

现在面临一般问题：$A$ 是结合代数、$M$ 是结合 $A$-模，我们想「导出」函子 $-\otimes_AM$ 与 $\operatorname{Hom}_A(M,-)$。它们应把 $A$-模的链复形变成阿贝尔群的复形，并保持锥与拟同构。我们将在下一讲开始处理。

\addcontentsline{toc}{section}{第 11 讲　映射复形与投射消解}
% 第 11 讲 映射复形与投射消解
% 源：src/lecture11.txt（264 行）
\section*{第 11 讲\quad 映射复形与投射消解}

本讲始终设 $A$ 是结合代数（可以非交换），例如 $A=\Z$ 或 $\Z[G]$，其中 $G$ 是（通常是有限的）群。回忆我们想要的规则是 $X\mapsto X_{hG},\,X^{hG}$：$X$ 是 $G$-模的复形，而 $X_{hG}$、$X^{hG}$ 是阿贝尔群的复形。这一构造的指导性「公理」将是：

\begin{enumerate}[label=(\arabic*),leftmargin=2.4em]
\item 若 $X$ \textbf{无环}，即对一切 $i\in\Z$ 有 $H_i(X)=0$，则希望 $X_{hG}$ 与 $X^{hG}$ 也无环。
\item $X_{hG}$ 与 $X^{hG}$ 都应与锥交换：若 $f\colon X\to Y$ 是 $G$-模复形的映射，则 $\operatorname{hCoker}(f)^{hG}\simeq\operatorname{hCoker}(X^{hG}\to Y^{hG})$。
\end{enumerate}

第二条条件相对容易满足，因为由断言 10.1 的证明，它只要求与有限直和及平移交换。

\begin{enumerate}[label=(\arabic*),leftmargin=2.4em,start=3]
\item 该构造应与不变子模、余不变商有关。即若 $X=(\cdots\to0\to M\to0\to\cdots)$ 只在次数 $0$ 有项，则 $H^0(X^{hG})=M^G$，且
\[
H^0(X_{hG})=M_G=\Z\otimes_{\Z[G]}M=M\Big/\sum_{g\in G}(g-1)M.
\]
\end{enumerate}

一个朴素而不正确的尝试是定义
\[
X^{hG}:=\Bigl(\cdots\to(X^{-1})^G\xrightarrow{\ d\ }(X^0)^G\xrightarrow{\ d\ }(X^1)^G\to\cdots\Bigr),
\]
其中
\[
X:=(\cdots\to X^{-1}\xrightarrow{\ d\ }X^0\xrightarrow{\ d\ }X^1\to\cdots).
\]
它平凡地满足 (2) 与 (3)；并注意它确实良定义，因为微分与群自同构交换。它还满足 (1) 的一个弱化版本：若 $X\simeq0$（即零复形），则 $X^{hG}$ 与 $X_{hG}$ 也同伦等价于 $0$。然而，这个构造不保持无环复形。具体地，若 $G:=\Z/2\Z$ 通过乘以 $1$ 与 $-1$ 作用在 $\Z$ 上，则
\[
\Bigl(\cdots\to0\to\Z\xrightarrow{\times2}\Z\to\Z/2\to0\to\cdots\Bigr)^{hG}
=\Bigl(\cdots\to0\to0\to0\to\Z/2\to0\to\cdots\Bigr),
\]
而后者不是无环的！

一个更省事的论证是：若该构造保持无环复形，则由 (10.1) 的构造，任意两个无环复形之间映射的锥必无环；又由假设它与锥交换，故由推论 10.11 它也会保持拟同构。但由断言 10.12，对 $G$-模的单射 $M\hookrightarrow N$，有 $\operatorname{hCoker}(i)\xrightarrow{\mathrm{qis}}\operatorname{Coker}(i)=N/M$（此后「qis」均指拟同构）。于是，若朴素不变子模保持无环复形，它就会保持余核，而我们知道这是错的。

注意，对结合代数 $A$，若 $A=\Z[G]$，则 $M\mapsto M^G=\operatorname{Hom}_{\Z[G]}(\Z,M)$，其中 $G$ 以平凡作用作用于 $\Z$。于是，对每个结合代数 $A$ 与 $A$-模 $M$ 都有一类一般的问题。

\textbf{定义 11.1.} 设 $X,Y$ 是 $A$-模的复形。\textbf{映射复形} $\operatorname{Hom}_A(X,Y)$ 是由
\[
\operatorname{Hom}_A^i(X,Y):=\prod_{j\in\Z}\operatorname{Hom}_A(X^j,Y^{j+i}),\qquad d^if:=df-(-1)^ifd
\]
定义的阿贝尔群复形（符号交替以保证微分平方为零）。可图示如下：
\[
\operatorname{Hom}^{-1}(X,Y)\xrightarrow{\ h\mapsto dh+hd\ }\operatorname{Hom}^0(X,Y)\xrightarrow{\ f\mapsto df-fd\ }\operatorname{Hom}^1(X,Y).
\]
其中各项按 $\operatorname{Hom}_A^i(X,Y)=\prod_j\operatorname{Hom}_A(X^j,Y^{j+i})$ 展开为分量：$\operatorname{Hom}^{-1}$ 的分量形如 $h^{-1}\colon X^{-1}\to Y^{-2}$、$h^0\colon X^0\to Y^{-1}$、$h^1\colon X^1\to Y^0$；$\operatorname{Hom}^0$ 的分量形如 $f^{-1}\colon X^{-1}\to Y^{-1}$、$f^0\colon X^0\to Y^0$、$f^1\colon X^1\to Y^1$；$\operatorname{Hom}^1$ 的分量形如 $g^{-1}\colon X^{-1}\to Y^0$、$g^0\colon X^0\to Y^1$、$g^1\colon X^1\to Y^2$。原图把这三列分量并排画出，并用 $h\mapsto dh+hd$、$f\mapsto df-fd$ 标出相邻两列之间的微分；其中 $d$ 分别指 $X$ 与 $Y$ 的微分。

\textbf{断言 11.2.} 对任意 $A$-模复形 $X,Y$，$\operatorname{Hom}_A(X,Y)$ 是复形。

注意，复形映射 $f\colon X\to Y$ 等价于元素 $f=(f^i)\in\operatorname{Hom}^0_A(X,Y)$ 且 $df=0$，其中 $d$ 指 $\operatorname{Hom}(X,Y)$ 上的微分。$f$ 的零伦同样等价于元素 $h\in\operatorname{Hom}^{-1}_A(X,Y)$ 且 $dh=f$，$d$ 同上。于是 $H^0\operatorname{Hom}(X,Y)$ 等价于 $X\to Y$ 的映射模同伦的等价类。因此这一构造推广了同调代数中许多重要概念。

\textbf{例 11.3.} 若 $X:=(\cdots\to0\to A\to0\to\cdots)$，$A$ 位于次数 $0$，则 $\operatorname{Hom}_A(X,Y)=Y$。因此这样的 $X$ 就是我们即将称为投射的对象。

\textbf{定义 11.4.} $A$-模的复形 $P$ 称为\textbf{投射的}（或称同伦投射、$K$-投射等；此概念由 Spaltenstein 定义），若对每个无环的 $A$-模复形 $Y$，$\operatorname{Hom}_A(P,Y)$ 也无环。

上面的困难在于：$\Z$ 作为 $\Z[G]$-模的复形并不投射。我们将说明，在某种意义下可以把 $\Z$「唯一地」替换为投射对象。

\textbf{引理 11.5.} 若 $P$ 是 $A$-模复形，$P^i=0$ 对一切 $i\gg0$（即 $P$ 的非零项在指标上有上界），且每个 $P^i$ 作为 $A$-模是投射的，则 $P$（作为复形）是投射的。

回忆下述定义：

\textbf{定义 11.6.} $A$-模 $P^i$ 作为 $A$-模\textbf{投射}，若以下等价条件之一成立：
\begin{enumerate}[label=(\arabic*),leftmargin=2.4em]
\item $\operatorname{Hom}_A(P^i,-)$ 保持满射；
\item $\operatorname{Hom}_A(P^i,-)$ 是正合函子；
\item $P^i$ 是自由模的直和项；
\item 对任意满射 $N\twoheadrightarrow M$，每个 $A$-模映射 $P^i\to M$ 都可提升为映射 $P^i\to N$，使下图交换：
\[
\begin{array}{ccc}
& & N\\
& \nearrow & \big\downarrow\\
P^i & \longrightarrow & M\longrightarrow0.
\end{array}
\]
\end{enumerate}

我们简要说明这些等价性。(1) 与 (4) 显然等价，因为 (4) 说的正是：若 $N\twoheadrightarrow M$，则 $\operatorname{Hom}_A(P^i,N)\twoheadrightarrow\operatorname{Hom}_A(P^i,M)$。条件 (2) 与 (1) 平凡地等价。为证 (3)，取 $M:=P^i$，$N$ 为某个满射到 $P^i$ 的自由模（例如以 $P^i$ 的全部元素为基，把对应元素送到对应元素）。于是 (4) 给出 $N\twoheadrightarrow P^i$ 的一个分裂，把 $P^i$ 实现为 $N$ 的直和项（原文此处印作 $M$ 的直和项，但此处 $M=P^i$，应为 $N$；已径改）。容易看出投射模的直和项仍投射；至于反向，只需注意自由模是投射的。

\textbf{断言 11.7.} $A$-模 $P$ 作为 $A$-模投射，当且仅当它作为次数 $0$ 的复形投射。

\textit{证明.} 若 $P$ 作为次数 $0$ 的复形投射，取满射 $f\colon N\twoheadrightarrow M$，并作无环复形
\[
X:=(\cdots\to0\to\ker(f)\to N\to M\to0\to\cdots).
\]
则 $\operatorname{Hom}_A(P,X)$ 为
\[
\cdots\to0\to\operatorname{Hom}_A(P,\ker(f))\to\operatorname{Hom}_A(P,N)\to\operatorname{Hom}_A(P,M)\to0\to\cdots,
\]
故 $\operatorname{Hom}_A(P,-)$ 保持满射，按定义 $P$ 作为 $A$-模投射。

反之，若 $X:=(\cdots\to X^{-1}\xrightarrow{d^{-1}}X^0\xrightarrow{d^0}X^1\to\cdots)$ 是无环复形且 $P$ 作为 $A$-模投射，则
\[
\operatorname{Hom}_A(P,X)=(\cdots\to\operatorname{Hom}_A(P,X^{-1})\to\operatorname{Hom}_A(P,X^0)\to\operatorname{Hom}_A(P,X^1)\to\cdots),
\]
它无环：因为若 $X^{i-1}\twoheadrightarrow\ker(X^i\to X^{i+1})$，则
\[
\operatorname{Hom}(P,X^{i-1})\twoheadrightarrow\operatorname{Hom}(P,\ker(X^i\to X^{i+1}))=\ker\bigl(\operatorname{Hom}(P,X^i)\to\operatorname{Hom}(P,X^{i+1})\bigr),
\]
这里用了假设 $\operatorname{Hom}(P,-)$ 正合、从而保持核。故 $\operatorname{Hom}_A(P,X)$ 无环，$P$ 作为次数 $0$ 的复形投射，如所欲证。\hfill$\square$

\textit{证明（引理 11.5）.} 设 $Y$ 是无环的 $A$-模复形。我们需要下述断言：

\textbf{断言 11.8.} 每个映射 $P\to Y$ 都零伦。

\textit{证明.} 设 $f\colon P\to Y$。我们用降序归纳构造 $f$ 的零伦 $h$。基例：对一切 $i\gg0$（其含义见引理陈述），有 $P^i=0$，故 $f^i=0$，于是可取 $h^i:=0$。现设已对一切 $j>i$ 有映射 $h\colon P^j\to Y^{j-1}$，使下图交换：
\[
\begin{array}{ccccccccc}
\cdots & \longrightarrow & P^{i-1} & \longrightarrow & P^i & \longrightarrow & P^{i+1} & \longrightarrow & \cdots\\
& & {\scriptstyle f}\big\downarrow & & {\scriptstyle f}\big\downarrow & & {\scriptstyle f}\big\downarrow & & \\
\cdots & \longrightarrow & Y^{i-1} & \longrightarrow & Y^i & \longrightarrow & Y^{i+1} & \longrightarrow & \cdots
\end{array}
\]
（图中标 $?$ 的正是待构造的 $h\colon P^i\to Y^{i-1}$）。我们希望构造 $h\colon P^i\to Y^{i-1}$ 使 $dh+hd=f$。注意，对一切 $x\in P^i$，
\[
(d(f-hd))(x)=(df-dhd)(x)=\bigl(df-(f-hd)d\bigr)(x)=(df-fd)(x)=0,
\]
这里中间一步用了归纳假设 $dh=f-hd$（作用在 $P^{i+1}$ 上）。由于 $(f-hd)\colon P^i\to\ker(d\colon Y^i\to Y^{i+1})$，且 $d\colon Y^{i-1}\to\ker(d\colon Y^i\to Y^{i+1})$ 是满射，而 $P^i$ 按假设投射，映射 $f-hd$ 可提升为映射 $h\colon P^i\to Y^{i-1}$ 使 $dh=f-hd$。于是 $dh+hd=f$，$h$ 给出 $f$ 的零伦，如所欲证。\hfill$\square$

由该断言，$H^0(\operatorname{Hom}_A(P,Y))$ 中只有零元（其元素即 $P\to Y$ 的同伦类）；且对每个 $i$，有
$H^i(\operatorname{Hom}_A(P,Y))=H^0(\operatorname{Hom}_A(P,Y[i]))=0$，因为 $Y[i]$ 也无环。故各阶上同调消没，$\operatorname{Hom}_A(P,Y)$ 无环，从而 $P$ 是投射的，如所欲证。\hfill$\square$

我们的计划大致是：证明每个 $X$ 都拟同构于某个投射复形 $P$，即 $P\xrightarrow{\mathrm{qis}}X$，称为 $X$ 的\textbf{投射消解}。此外，$P$ 在某种稍后定义的意义下是「唯一」或「导出」的。于是得到「修正」版本 $\operatorname{Hom}^{\mathrm{der}}_A(X,Y):=\operatorname{Hom}_A(P,Y)$。取 $A:=\Z[G]$ 并选取某个投射消解 $P\xrightarrow{\mathrm{qis}}\Z$（它相当典范，对有限群甚至更简单，虽非严格唯一，但对上同调没有影响），便可定义 $X^{hG}:=\operatorname{Hom}_G(P,X)$。这将满足我们的全部公理：它与不变子模有关（因为 $P$ 类似于 $\Z$），又保持无环复形（因为 $P$ 投射）！

下述命题足以证明第一点，因为我们关心的复形是次数 $0$ 处的 $\Z$，它显然有上界。

\textbf{命题 11.9.} 设 $X$ 是 $A$-模复形，且 $X$ 有上界，即如前 $X^i=0$ 对一切 $i\gg0$。则存在投射消解 $P\xrightarrow{\mathrm{qis}}X$。

\textit{证明.} 不妨设 $X$ 在指标 $0$ 处有上界。取自由模 $P^0$ 经映射 $\alpha^0$ 满射到 $X^0$（如前，取生成元即可，使核恰为生成元之间的关系）。再取 $P^{-1}$ 为自由模，满射到纤维积 $P^0\times_{X^0}X^{-1}$：
\[
\begin{array}{ccccccc}
P^{-1} & \longrightarrow & P^0 & \longrightarrow & 0 & \longrightarrow & \cdots\\
\big\downarrow & & \big\downarrow{\scriptstyle\alpha^0} & & & & \\
X^{-1} & \longrightarrow & X^0 & \longrightarrow & 0 & \longrightarrow & \cdots
\end{array}
\]
该构造保持上同调，因为
\[
H^0X=X^0/\operatorname{Im}(X^{-1})=P^0/P^{-1}=H^0P,
\]
这是由于 $P^{-1}$ 满射到 $\ker(\alpha^0)$，且它在 $P^0/\ker(\alpha^0)\cong X^0$ 中的像等于 $X^{-1}$（因为 $P^0\twoheadrightarrow X^0$）。又因 $P^{-1}\twoheadrightarrow X^{-1}$，可以迭代此过程，构造出 $X$ 的投射消解 $P$，如所欲证。\hfill$\square$

第二个论断是关于投射消解的唯一性，由下述引理给出：

\textbf{引理 11.10.} 设 $P_1$ 与 $P_2$ 都是 $A$-模复形 $X$ 的投射消解。则存在同伦等价 $\gamma$，使下图在同伦意义下交换，即 $\beta\gamma\simeq\alpha$：
\[
\begin{array}{ccccc}
& P_1 & \xrightarrow{\ \gamma\ } & P_2 & \\
& {\scriptstyle\alpha}\big\downarrow & & \big\downarrow{\scriptstyle\beta} & \\
& & X & &
\end{array}
\]

\textit{证明.} 考虑下图：
\[
\begin{array}{ccccc}
& P_1 & \xrightarrow{\ \gamma\ } & P_2 & \\
& {\scriptstyle\alpha}\big\downarrow & & \big\downarrow{\scriptstyle\beta} & \\
& X & \xrightarrow{\ \delta\ } & \operatorname{hCoker}(\beta) &
\end{array}
\]
由于 $\beta$ 按假设是拟同构，$\operatorname{hCoker}(\beta)$ 由推论 10.11 无环。由断言 11.8，复合 $\delta$ 零伦；故由断言 10.8，存在典范映射
\[
P_1\to\operatorname{hKer}\bigl(X\to\operatorname{hCoker}(\beta)\bigr)\simeq P_2,
\]
且为同伦等价，如所欲证。由对称性，反方向也存在这样的映射，故 $\gamma$ 是同伦等价，而该图在同伦意义下平凡地交换。\hfill$\square$

现在可以问这里的 $\gamma$ 有多唯一。答案如下：

\textbf{断言 11.11.} 所有这样的 $\gamma$ 都同伦。

\textit{证明.} 我们模仿引理 11.10 的证明，把单个态射换成 $\operatorname{Hom}$-复形。我们有映射
\[
\operatorname{Hom}(P_1,P_2)\xrightarrow{\ \beta_*\ }\operatorname{Hom}(P_1,X)\to\operatorname{Hom}(P_1,\operatorname{hCoker}(\beta))=\operatorname{hCoker}(\beta_*),
\]
其中 $\beta_*$ 由与 $\beta$ 复合给出，最后的等同出于形式理由。由于 $P_2$ 投射，最后的复形无环（按定义），故由推论 10.11，映射 $\psi\colon\operatorname{Hom}(P_1,P_2)\to\operatorname{Hom}(P_1,X)$ 是拟同构，从而是映射同伦类上的同构。特别地，
% [原文疑误：原印 $H^0\operatorname{Hom}(P_1,P_2)=H^0\operatorname{Hom}(P_2,X)$；但 $\psi:\operatorname{Hom}(P_1,P_2)\to\operatorname{Hom}(P_1,X)$ 是拟同构，且给定的映射是 $\alpha\in H^0\operatorname{Hom}(P_1,X)$，故右端应为 $P_1$。此处径改。]
\[
H^0\operatorname{Hom}(P_1,P_2)=H^0\operatorname{Hom}(P_1,X)
\]
于是既然 $H^0\operatorname{Hom}(P_1,X)$ 中给定了映射，$H^0\operatorname{Hom}(P_1,P_2)$ 中的诱导映射便在同伦意义下良定义（如定义 11.1 之后的讨论所述）。\hfill$\square$

事实上可以证明，同伦之间的所有同伦也都同伦，依此类推；故在建立唯一性方面，这已是我们能期望的最佳结果。

\textbf{定义 11.12.} 两个 $A$-模链复形 $M$ 与 $N$ 的第 $i$ 个 \textbf{Ext}-群定义为 $\operatorname{Ext}^i_A(M,N):=H^i\operatorname{Hom}(P,N)$，其中 $P$ 是 $M$ 的某个投射消解。正如我们刚才所示，此定义与 $P$ 的选取无关。

\addcontentsline{toc}{section}{第 12 讲　导出函子与显式投射消解}
% 第 12 讲 导出函子与显式投射消解
% 源：src/lecture12.txt（244 行）
\section*{第 12 讲\quad 导出函子与显式投射消解}

设 $X$ 与 $Y$ 是 $A$-模复形。回忆上一讲定义了 $\operatorname{Hom}_A(X,Y)$，以及
\[
\operatorname{Hom}^{\mathrm{der}}_A(X,Y):=\operatorname{Hom}_A(P,Y),
\]
其中 $P\xrightarrow{\mathrm{qis}}X$ 是投射复形，即 $X$ 的一个投射消解。我们还定义了 \textbf{Ext}-群
\[
\operatorname{Ext}^i_A(X,Y):=H^i\operatorname{Hom}^{\mathrm{der}}_A(X,Y).
\]
这里最重要的例子是 $A:=\Z[G]$，此时 $X^{hG}:=\operatorname{Hom}^{\mathrm{der}}_A(\Z,X)$ 就是 $X$ 的同伦不变量。这一构造具有如下基本性质：$\operatorname{Hom}^{\mathrm{der}}_A(X,-)$ 保持拟同构；且 $P\xrightarrow{\mathrm{qis}}X$ 在同伦意义下唯一，这样的同伦又在同伦意义下唯一，如此递推。

附带说明，我们其实可以更典范地定义导出函子 $\operatorname{Hom}^{\mathrm{der}}(X,-)$，而不必选取特定的投射消解：
\[
\operatorname{Hom}^{\mathrm{der}}(X,Y):=\varinjlim_{P\xrightarrow{\mathrm{qis}}X,\ P\ \text{投射}}\operatorname{Hom}(P,Y),
\]
其中 $P$ 按链复形映射 $P'\xrightarrow{\mathrm{qis}}P\xrightarrow{\mathrm{qis}}X$ 排序，这便强制消除了对 $P$ 的选取。

\textbf{断言 12.1.} 设 $f\colon X_1\to X_2$ 是链复形映射，$X_1$、$X_2$ 分别有投射消解 $P_1$、$P_2$。则存在映射 $\phi\colon P_1\to P_2$，使下图在同伦意义下交换：
\[
\begin{array}{ccc}
X_1 & \xrightarrow{\ f\ } & X_2\\
{\scriptstyle\mathrm{qis}}\big\uparrow & & \big\uparrow{\scriptstyle\mathrm{qis}}\\
P_1 & \xrightarrow{\ \phi\ } & P_2
\end{array}
\]
此外，这样的映射在同伦意义下唯一。

\textit{证明.} 由于导出函子保持拟同构，诱导出的映射复形之间的映射
\[
\operatorname{Hom}(P_1,P_2)\xrightarrow{\mathrm{qis}}\operatorname{Hom}(P_1,X_2)
\]
是拟同构。我们已给定一个映射，即复合 $P_1\to X_1\to X_2$；它被微分杀死（因为它是链复形映射），从而定义了 $H^0\operatorname{Hom}(P_1,X_2)$ 中的一个上同调类。于是 $H^0\operatorname{Hom}(P_1,P_2)$ 中存在某个上同调类，它是该映射经 $P_2$ 的提升，且良定义、在同伦意义下唯一。\hfill$\square$

由此可见，对每个链复形 $Y$，沿 $\phi$ 拉回便得到映射
\[
\operatorname{Hom}^{\mathrm{der}}_A(X_2,Y)\to\operatorname{Hom}^{\mathrm{der}}_A(X_1,Y).
\]
一个快速的「应用」如下：

\textbf{断言 12.2.} 若 $H\subseteq G$ 是子群，$X$ 是 $G$-模复形，则得到限制映射 $X^{hG}\to X^{hH}$（它在同伦意义下良定义）。

直观地说，$G$-不变的东西当然也是 $H$-不变的。

\textit{证明.} 考虑
\[
\Z[G/H]=\{f\colon G/H\to\Z\mid f\ \text{只在有限多个点处非零}\},
\]
它通过平移带有 $G$-作用，并且等同于从 $H$ 到 $G$ 由平凡模诱导出的模 $\Z[G]\otimes_{\Z[H]}\Z$。

\textbf{断言 12.3.} $\operatorname{Hom}^{\mathrm{der}}_G(\Z[G/H],X)\simeq X^{hH}$ 是拟同构。

\textit{证明.} 设 $P_H\xrightarrow{\mathrm{qis}}\Z$，其中 $P_H$ 是 $H$-模的投射复形。则得到诱导的 $G$-模 $\Z[G]\otimes_{\Z[H]}P_H$。注意 $\Z[G]$ 作为 $\Z[H]$-模是自由的：取 $G/H$ 的陪集代表元即得一基。因此 $\Z[G]\otimes_{\Z[H]}-$ 保持拟同构（我们已知 $\Z[H]$ 的情形，而 $\Z[G]$ 可看作 $\Z[H]$ 的若干拷贝的直和）。这意味着
\[
\Z[G]\otimes_{\Z[H]}P_H\xrightarrow{\mathrm{qis}}\Z[G/H],
\]
而后者作为 $\Z[G]$-模的复形是投射的：这是因为 $\Z[G]$ 与 $P_H$ 都有界，故它也有界，而诱导到 $\Z[G]$ 保持投射模，因为我们仍得到自由模的直和项。或者也可用下述万有性质：到无环复形的每个映射都零伦，而从诱导复形出发的 $G$-等变映射等同于从 $P_H$ 出发的 $H$-等变映射。于是由定义
\[
\operatorname{Hom}^{\mathrm{der}}_G(\Z[G/H],X):=\operatorname{Hom}_G(\Z[G]\otimes_{\Z[H]}P_H,X)=\operatorname{Hom}_H(P_H,X)=:X^{hH},
\]
即得断言。\hfill$\square$

由此可见，我们经映射
\[
\varepsilon\colon\Z[G/H]\to\Z,\qquad \sum_{g_i\in G/H}n_ig_i\mapsto\sum_i n_i
\]
得到映射 $X^{hG}\to X^{hH}$；当 $\Z$ 取平凡作用时，$\varepsilon$ 显然是 $G$-等变的。由前述讨论，我们得到导出函子之间的限制映射
\[
X^{hG}=\operatorname{Hom}^{\mathrm{der}}_G(\Z,X)\to\operatorname{Hom}^{\mathrm{der}}_G(\Z[G/H],X)=X^{hH},
\]
它在同伦意义下良定义（同伦意义下良定义，如此递推，即我们「惯常的错误」）。\hfill$\square$

回忆此处一切都是阿贝尔群的复形，故不存在「类型不匹配」。事实上，若 $H\le G$ 是有限指数的，则有 $G$-等变映射
\[
\Z\xrightarrow{\ \kappa\ }\Z[G/H]\xrightarrow{\ \varepsilon\ }\Z,\qquad 1\mapsto\sum_{g\in G/H}g,
\]
其复合对应乘以指数 $[G:H]$。这给出膨胀映射 $X^{hH}\to X^{hG}$，使得复合 $X^{hG}\to X^{hH}\to X^{hG}$ 同伦于乘以 $[G:H]$。

更具体地说，设有 $H$-不变的对象与 $G$-不变的对象。取 $G/H$ 的陪集代表元，便可对任一 $H$-不变元素取「相对范数」，得到一个 $G$-不变元素。这正是上面的映射所做的，也解释了复合为何乘以 $[G:H]$。

\textbf{定义 12.4.} $H^i(G,X):=H^i(X^{hG})$ 称为 $G$ 以 $X$ 为系数的（超）上同调。

前缀「超」过去用来指复形上的运算；若该复形只在次数 $0$ 有项，则称为「群上同调」。

\textbf{断言 12.5.} 若 $X$ 只在非负次数有项，即 $X^i=0$ 对一切 $i<0$，则 $H^0(G,X)=H^0(X)^G$，且 $H^i(G,X)=0$ 对 $i<0$。

\textit{证明.} 取 $\Z$ 作为 $G$-模的某个投射消解 $P$；由命题 11.9 的构造，可取 $P$ 只在非正次数有项。按定义，$H^0(G,X)$ 等同于映射 $f\colon P\to X$ 的同伦类，而所有这些映射都形如
\[
\begin{array}{ccccccccc}
\cdots & \longrightarrow & P^{-2} & \longrightarrow & P^{-1} & \longrightarrow & P^0 & \longrightarrow & 0\\
& & {\scriptstyle f}\big\downarrow & & \big\downarrow & & \big\downarrow & & \\
\cdots & \longrightarrow & 0 & \longrightarrow & 0 & \longrightarrow & X^0 & \longrightarrow & X^1
\end{array}
\]
于是 $f$ 的任何同伦都是 $0$，同理 $H^i(G,X)=0$ 对 $i<0$。由交换性必有 $df=fd=0$。由此可见，这样的映射 $f$ 等同于 $G$-等变映射
\[
\Z=P^0/dP^{-1}=\operatorname{Coker}(P^{-1}\to P^0)\to\ker(X^0\to X^1)=H^0(X),
\]
而后者按拟同构等同于 $H^0(X)$ 中的一个 $G$-不变向量（即经由 $1$ 的像给出）。\hfill$\square$

下面转向构造 $\Z$ 作为 $G$-模的显式投射消解。

\textbf{例 12.6.} 设 $G:=\Z/n\Z$，生成元为 $\sigma$。我们断言下述映射是拟同构：
\[
\begin{array}{ccccccccccc}
\cdots & \longrightarrow & \Z[G] & \xrightarrow{1-\sigma} & \Z[G] & \xrightarrow{\sum_i\sigma^i} & \Z[G] & \xrightarrow{1-\sigma} & \Z[G] & \xrightarrow{\varepsilon} & 0\\
& & \big\downarrow & & \big\downarrow & & \big\downarrow & & \big\downarrow & & \\
\cdots & \longrightarrow & 0 & \longrightarrow & 0 & \longrightarrow & 0 & \longrightarrow & \Z & \longrightarrow & 0
\end{array}
\]
负阶上同调的消没来自我们早先关于 Tate 上同调的结果；而 $\varepsilon$ 的核，即坐标和为零的元素的集合，正是 $1-\sigma$ 的像。

\textbf{推论 12.7.} 若 $M$ 是 $G$-模（视为次数 $0$ 处的复形），则 $M_{hG}$ 拟同构于复形
\[
\cdots\to0\to0\to M\xrightarrow{1-\sigma}M\xrightarrow{\sum_i\sigma^i}M\xrightarrow{1-\sigma}\cdots,
\]
其中第一个 $M$ 位于次数 $0$。

注意，从 $\Z[G]$ 到任一对象的 $G$-等变映射就是该对象的一个元素。确实，不变子模正是零阶上同调群，正如所愿。因此这一构造只给出了我们早先关于 Tate 上同调所学内容的「一半」。

现在我们想对每个 $G$ 构造显式消解。贯穿始终的「格言」是：「这类东西都来自 bar 构造。」设 $A$ 是交换环，$B$ 是 $A$-代数；最重要的情形是 $A:=\Z$、$B:=\Z[G]$。

\textbf{定义 12.8.} 对上述 $A$ 与 $B$，\textbf{bar 复形} $\operatorname{Bar}_A(B)$ 为
\[
\cdots\to B\otimes_AB\otimes_AB\xrightarrow{b_1\otimes b_2\otimes b_3\mapsto b_1b_2\otimes b_3-b_1\otimes b_2b_3}B\otimes_AB\xrightarrow{b_1\otimes b_2\mapsto b_1b_2}B\to0\to\cdots,
\]
其中 $B$ 位于次数 $0$。一般地，$\operatorname{Bar}^{-n}_A(B):=B^{\otimes_A n+1}$，微分为
\[
b_1\otimes\cdots\otimes b_{n+1}\mapsto b_1b_2\otimes b_3\otimes\cdots\otimes b_{n+1}
-b_1\otimes b_2b_3\otimes\cdots\otimes b_{n+1}
+b_1\otimes b_2\otimes b_3b_4\otimes\cdots\otimes b_{n+1}-\cdots.
\]

容易看出这一微分的平方为零：只需有选择地删去张量因子并检查符号，故它确实是链复形。

\textbf{断言 12.9.} $\operatorname{Bar}_A(B)$ 同伦等价于零。

\textit{证明.} 我们要找恒等映射的零伦，即映射 $h$ 使 $hd+dh=\mathrm{id}$：
\[
\begin{array}{ccccccccc}
\cdots & \longrightarrow & B\otimes_AB\otimes_AB & \longrightarrow & B\otimes_AB & \longrightarrow & B & \longrightarrow & 0\\
& & {\scriptstyle h_2}\big\uparrow{\scriptstyle\mathrm{id}} & & {\scriptstyle h_1}\big\uparrow{\scriptstyle\mathrm{id}} & & {\scriptstyle h_0}\big\uparrow{\scriptstyle\mathrm{id}} & & \\
\cdots & \longrightarrow & B\otimes_AB\otimes_AB & \longrightarrow & B\otimes_AB & \longrightarrow & B & \longrightarrow & 0
\end{array}
\]
于是定义 $h_0(b):=1\otimes b$，$h_1(b_1\otimes b_2):=1\otimes b_1\otimes b_2$。确实，此时有
\[
(dh_1+h_0d)(b_1\otimes b_2)=d(1\otimes b_1\otimes b_2)+1\otimes b_1b_2=b_1\otimes b_2-1\otimes b_1b_2+1\otimes b_1b_2=b_1\otimes b_2,
\]
正如所愿。容易证明，对一切 $n$ 类似地定义 $h_n$ 便给出恒等映射的零伦。\hfill$\square$

换个说法，考虑图
\[
\begin{array}{ccccccc}
\cdots & \longrightarrow & B\otimes_AB\otimes_AB & \longrightarrow & B\otimes_AB & \longrightarrow & 0\\
& & & & {\scriptstyle\mathrm{id}}\big\downarrow & & \big\downarrow{\scriptstyle b\otimes b'\mapsto bb'}\\
\cdots & \longrightarrow & 0 & \longrightarrow & B & \longrightarrow & 0
\end{array}
\]
其中微分中的 $d$ 就是乘法映射。这是同伦等价，因为它的锥正是 $\operatorname{Bar}_A(B)$。

把每一项看作双模（即左右各有交换作用的模），其中在首项上从左边乘以 $B$，在末项上从右边乘以 $B$。于是这些微分都是双模同态。给定（左）$B$-模 $M$，与 $M$ 在 $B$ 上作张量，便得到图
\[
\begin{array}{ccccccc}
\cdots & \longrightarrow & B\otimes_AB\otimes_AM & \longrightarrow & B\otimes_AM & \longrightarrow & 0\\
& & & & {\scriptstyle b\otimes m\mapsto bm}\big\downarrow & & \\
\cdots & \longrightarrow & 0 & \longrightarrow & M & \longrightarrow & 0
\end{array}
\]
它也是同伦等价（映射 $B\otimes_AM\to M$ 就是「作用映射」；并注意这些张量积有意义，因为 $B$ 是 $A$-模）。微分不变，只是末项换成了 $M$ 的元素，例如
\[
b_1\otimes b_2\otimes m\mapsto b_1b_2\otimes m-b_1\otimes b_2m.
\]
换言之，$M$ 典范地同伦等价于一个复形，其中每一项都形如 $B\otimes_AN$；此处 $N$ 代表 $B\otimes_A\cdots\otimes_AB\otimes_AM$，即由某个 $A$-模诱导出的模。

现在把这应用于 $A:=\Z$、$B:=\Z[G]$、$M:=\Z$（即平凡模）的情形。注意这只是 $B$-模复形之间的拟同构，而并非同伦等价：因为逆映射只是 $A$-线性的，而不是 $B$-线性的！确实，这样的逆映射应为
\[
B\otimes_AM\xrightarrow{b\otimes m\mapsto bm}M\xrightarrow{m\mapsto1\otimes m}B\otimes_AM,
\]
即 $b\otimes m\mapsto bm\mapsto1\otimes bm\neq b\otimes m=b(1\otimes m)$，因为作用在 $B$ 上而不在 $M$ 上，这并不是 $B$-模的作用。一般地，单向拟同构的存在并不蕴含反方向拟同构的存在；而同伦等价按定义包含反方向的映射，故是对称的。另外回忆，同伦等价蕴含拟同构。

于是我们得到 $\Z$ 的自由 $G$-模典范投射消解
\[
\begin{array}{ccccccccc}
\cdots & \longrightarrow & \Z[G^3] & \longrightarrow & \Z[G\times G] & \longrightarrow & \Z[G] & \longrightarrow & 0\\
& & \big\downarrow & & \big\downarrow & & \big\downarrow & & \\
\cdots & \longrightarrow & 0 & \longrightarrow & 0 & \longrightarrow & \Z & \longrightarrow & 0
\end{array}
\]
其微分为 $(g_1,g_2)\mapsto g_1g_2-g_1$，等等，因为 $G$ 平凡地作用于 $\Z$。注意 $\Z[G\times G]\cong\Z[G]\otimes_\Z\Z[G]$，因为两者都以乘积群的元素为基。

这是计算群上同调的一个很好的显式投射消解！最终得到形如
\[
\cdots\to0\to M\to\Z[G]\otimes M\to\Z[G\times G]\otimes M\to\cdots
\]
的复形，其中 $M$ 位于次数 $0$，且 $G$ 有限。
% [原文疑误：原印「位于 $\Z[G]\otimes M$ 的微分核中的元素」——按 $n$-上闭链的定义应为 $\Z[G^n]\otimes M$（该复形的第 $n$ 项是 $\Z[G^n]\otimes M$，其中 $\Z[G\times G]=\Z[G^2]$），此处径改]
位于 $\Z[G^n]\otimes M$ 的微分核中的元素称为以 $M$ 为系数的\textbf{群 $n$-上闭链}；位于微分像中的元素称为 \textbf{$n$-上边缘}。

\addcontentsline{toc}{section}{第 13 讲　同伦余不变商、交换化与 Tate 上同调}
% 第 13 讲：同伦余不变商、交换化与 Tate 上同调
% 源：src/lecture13.txt（235 行，含 2 处页眉）
\section*{第 13 讲\quad 同伦余不变商、交换化与 Tate 上同调}

回忆上一讲我们显式构造了 $G$-模复形 $X$ 的同伦不变量 $X^{hG}$。为此我们构造了棒消解 $P_G^{\mathrm{can}}\xrightarrow{\ \sim\ }\Z$，其中 $P_G^{\mathrm{can}}$ 是一个典范的、位于非正次数的自由 $G$-模复形。于是有拟同构
\[
X^{hG}\simeq\operatorname{Hom}_G\bigl(P_G^{\mathrm{can}},X\bigr).
\]
具体地，对 $P_G^{\mathrm{can}}$ 有
\[
\begin{array}{ccccccc}
\cdots & \longleftarrow & \Z[G^{3}] & \longleftarrow & \Z[G\times G] & \longleftarrow & \Z[G]\longleftarrow 0\longleftarrow\cdots\\[2pt]
\cdots & \longleftarrow & 0 & \longleftarrow & 0 & \longleftarrow & \Z\longleftarrow 0\longleftarrow\cdots
\end{array}
\]
两行由增广映射 $\Z[G]\to\Z$ 连接，微分为 $d_{-1}(g_1,g_2)=g_1g_2-g_1$（$G$-作用始终作用在第一个分量上）。注意若 $G$ 有限，则这些都是有限秩 $G$-模。

对每个 $G$-模 $M$，我们有
\[
\cdots\to 0\to M\xrightarrow{\ m\mapsto(gm-m)_{g\in G}\ }\underbrace{\prod_{g\in G}M}_{\{\varphi\colon G\to M\}}\to\prod_{g,h\in G}M\to\cdots,
\]
% [原文疑误：此段三处的 $M_{hG}$ 均应为 $M^{hG}$——该复形是 $\operatorname{Hom}_G(P_G^{\mathrm{can}},M)$（同伦不变子模），其零阶上同调为不变子模 $M^G$、一阶为通常的 $H^1(G,M)$；已径改]
经由某个进一步的微分，得到 $M^{hG}$。我们可以用这个表达式显式计算 $M^{hG}$ 的一阶上同调。结果表明，函数 $\varphi\colon G\to M$ 被该微分杀掉，当且仅当它是 $1$-上闭链（有时称为扭同态），即对一切 $g,h\in G$ 有
\[
\varphi(gh)=\varphi(g)+g\cdot\varphi(h)
\]
（群作用意义下）。类似地，$\varphi$ 是 $1$-上边缘，如果存在 $m\in M$ 使得对一切 $g\in G$ 有 $\varphi(g)=g\cdot m-m$。要点是
\[
H^{1}(G,M):=H^{1}(M^{hG})=\{1\text{-上闭链}\}/\{1\text{-上边缘}\}.
\]
作为推论，若 $G$ 平凡地作用在 $M$ 上，则 $H^{1}(G,M)=\operatorname{Hom}_{\mathrm{Group}}(G,M)$：因为此时 $1$-上边缘全为零，而 $1$-上闭链恰好是通常的群同态。这也说明零阶上同调恰好就是不变子模，正如上一讲所示。

现在，我们的目标（很久以前就定下的）是对任意有限群 $G$ 定义 Tate 上同调与 Tate 复形。我们希望
\[
\widehat{H}^{0}(G,M)=M^{G}/N(M)=\operatorname{Coker}\bigl(M^{G}\xrightarrow{\ N\ }M^{G}\bigr),
\]
因为它推广了局部域扩张的局部类域论的核心问题。回忆 $M_{G}=M/(g-1)M$（等价于与平凡模作张量，且与不变子模对偶——我们更愿意把后者看作子模），于是这个映射穿过 $M_{G}$，并诱导出上面的范映射。

我们的计划是：对 $G$-模复形 $X$，构造
\[
X_{hG}\xrightarrow{\ N\ }X^{hG}\longrightarrow X_{tG}:=\operatorname{hCoker}(N).
\]
因此我们首先需要定义同伦余不变商 $X_{hG}$。

注意若 $M$ 是 $G$-模，则 $M_{G}=M\otimes_{\Z[G]}\Z$。令 $I_{G}:=\ker(\varepsilon)$，于是有短正合列
\[
0\to I_{G}\to\Z[G]\xrightarrow{\ \varepsilon\ }\Z\to 0,\qquad \sum_{i}n_{i}g_{i}\mapsto\sum_{i}n_{i}.
\]
我们断言 $I_{G}$ 作为 $\Z$-模由 $\{g-1:g\in G\}$ 张成（留作习题）。一个推论是
\[
\Z[G]^{\oplus G}\to\Z[G]\to\Z\to 0
\]
正合，因为 $\Z[G]^{\oplus G}\cong I_{G}$，其中第 $g$ 个坐标上的 $1$ 映到 $g-1$。

\textbf{注 13.1.} 计算张量积的正确算法如下：回忆张量积是右正合的，即它保持满射，且与代数作张量给出原来的模。要与一个模作张量，先取该模的生成元与关系，用它写出如上的一条消解，再与这条消解作张量，得到一个定义在该模的若干拷贝的直和上的矩阵，然后取余核。

如果我能用类似的过程对链复形定义 $M_{hG}$，那将非常方便。

\textbf{定义 13.2.} 若 $X$ 与 $Y$ 是链复形，则
\[
(X\otimes Y)_{i}:=\bigoplus_{j\in\Z}X_{j}\otimes Y_{i-j},
\]
微分为
\[
d(x\otimes y):=dx\otimes y+(-1)^{j}x\otimes dy.
\]
若 $X$ 是右 $A$-模复形、$Y$ 是左 $A$-模复形，则 $X\otimes_{A}Y$ 类似地定义。

注意因子 $(-1)^{j}$ 保证微分的平方为零。另外，对代数而言不必担心左、右 $A$-模之别，因为把乘法次序反过来就给出左代数与右代数之间的同构；对 $G$-模而言，这意味着把每个元素换成它的逆。

现在，对 $X_{hG}$ 的一个糟糕猜测是 $X\otimes_{\Z[G]}\Z$，因为它不保持无环复形，等价地不保持拟同构。更好的猜测是取投射消解 $P_{G}\simeq\Z$（例如 $P_G^{\mathrm{can}}$），并用它作张量：
\[
X_{hG}:=X\otimes_{\Z[G]}P_{G}.
\]

\textbf{定义 13.3.} 左 $A$-模复形 $F$ 称为\textbf{平坦的}，如果对每个右 $A$-模的无环复形 $Y$，$Y\otimes_{A}F$ 也无环，即 $-\otimes_{A}F$ 保持单射。

我们现在问 $P_{G}$ 是否平坦。事实上：

\textbf{断言 13.4.} 任何投射复形都平坦。

一个更容易的断言是：

\textbf{断言 13.5.} 任何上有界且各项 $F^{i}$ 都平坦的复形 $F$ 都平坦。

为证明这个断言，我们要用到投射模平坦这一事实：它们是自由模的直和项，而自由模显然平坦（即若 $F=F_{1}\oplus F_{2}$，则 $F\otimes M=(F_{1}\otimes M)\oplus(F_{2}\otimes M)$）。

\textit{证明.} 情形 1. 设 $F$ 只出现在 $0$ 次，即 $i\neq 0$ 时 $F^{i}=0$。对每个复形 $Y=Y^{\bullet}$，$Y\otimes_{A}F$ 为
\[
\cdots\to Y^{i}\otimes_{A}F\xrightarrow{\ d^{i}\otimes\mathrm{id}_{F}\ }Y^{i+1}\otimes_{A}F\to\cdots.
\]
由于 $F$ 平坦，对每个 $i$ 有 $H^{i}(Y\otimes_{A}F)=H^{i}(Y)\otimes_{A}F$（因为 $F$ 平坦意味着与 $F$ 作张量与取核、余核、像交换），故若 $Y$ 无环，则 $Y\otimes_{A}F$ 也无环。

情形 2. 设 $F$ 只出现在 $0$ 与 $-1$ 次，即 $F$ 形如
\[
\cdots\to 0\to F^{-1}\to F^{0}\to 0\to\cdots,
\]
从而 $F^{\bullet}=\operatorname{hCoker}(F^{-1}\to F^{0})$。由于张量积与同伦余核交换，我们得到
\[
Y\otimes_{A}F=\operatorname{hCoker}\bigl(Y\otimes_{A}F^{-1}\to Y\otimes_{A}F^{0}\bigr),
\]
于是由情形 1，若 $Y$ 无环，则 $Y\otimes_{A}F^{0}$ 与 $Y\otimes_{A}F^{-1}$ 也无环，从而由上同调长正合列，$Y\otimes_{A}F$ 无环。类似的（归纳）论证给出 $F$ 有界的情形。

情形 3. 一般情形，考虑截断复形
\[
\begin{array}{l}
F^{[0]}=\bigl(\cdots\to 0\to 0\to F^{0}\to 0\to\cdots\bigr),\\[2pt]
F^{[1]}=\bigl(\cdots\to 0\to F^{-1}\to F^{0}\to 0\to\cdots\bigr),\\[2pt]
F^{[2]}=\bigl(\cdots\to F^{-2}\to F^{-1}\to F^{0}\to 0\to\cdots\bigr),\ \dots
\end{array}
\]
相邻截断之间的自然映射 $F^{[i]}\to F^{[i+1]}$ 在 $-i+1,\dots,0$ 次为恒等映射、在 $-i$ 次为微分 $d$；原图中的所有方块显然交换，故这些都是复形的态射，且 $F=\varinjlim_{i}F^{[i]}$。由于正向极限与张量积交换（注意逆向极限不然，因为满射性），有 $Y\otimes_{A}F=\varinjlim_{i}Y\otimes_{A}F^{[i]}$。由情形 2，每个 $Y\otimes_{A}F^{[i]}$ 无环；又上同调与正向极限交换（因为后者保持核、余核与像），故若 $Y$ 无环，则 $Y\otimes_{A}F$ 也无环。\hfill$\square$

\textbf{注 13.6.} 设 $Y$ 是 $A$-模复形，取拟同构 $F\xrightarrow{\ \sim\ }Y$，其中 $F$ 平坦，并定义
\[
Y\otimes_{A}^{\mathrm{der}}X:=F\otimes_{A}X.
\]
则它在拟同构意义下良定义，而拟同构又在同伦意义下良定义，如此以往（真是一路到底！）。

\textbf{定义 13.7.}（$Y$ 相对于 $X$ 的）第 $i$ 个\textbf{挠群}为
\[
\operatorname{Tor}_{i}^{A}(Y,X):=H^{-i}\bigl(Y\otimes_{A}^{\mathrm{der}}X\bigr).
\]

\textbf{定义 13.8.} 链复形 $X$ 的\textbf{同伦余不变商}是复形
\[
X_{hG}:=X\otimes_{\Z[G]}^{\mathrm{der}}\Z\simeq X\otimes_{\Z[G]}P_{G}
\]
（注意它只在拟同构意义下良定义）。

\textbf{定义 13.9.} $H_{i}(G,X):=H^{-i}(X_{hG})$（注意这里更偏好下标记号，因为 $X_{hG}$ 一般只在非正次数上有定义）。

下面做一些基本计算。

\textbf{断言 13.10.} 若 $X$ 被 $0$ 从上界住，则 $H_{0}(G,X)=H_{0}(X)_{G}$（证明与断言 12.5 类似）。

\textbf{断言 13.11.} $H_{1}(G,\Z)=G^{\mathrm{ab}}$，其中 $G^{\mathrm{ab}}$ 表示 $G$ 的交换化。

注意这可以说是与本讲开头所见结论相对偶的陈述；$H^{1}(G,M)$ 关乎映射 $G\to M$，这等同于映射 $G^{\mathrm{ab}}\to M$，而这里 $H_{1}(G,\Z)$ 由走出 $G$ 的映射决定。

\textit{证明.} 回忆短正合列
\[
0\to I_{G}\to\Z[G]\xrightarrow{\ \varepsilon\ }\Z\to 0.
\]
上同调长正合列给出正合列
\[
H_{1}(G,\Z[G])\to H_{1}(G,\Z)\to H_{0}(G,I_{G})\to H_{0}(G,\Z[G])\to H_{0}(G,\Z).
\]
我们有
\[
H_{0}(G,\Z[G])=H_{0}(\Z[G])_{G}=\Z[G]\otimes_{\Z[G]}\Z=\Z
\]
（由断言 13.10）。当然 $H_{0}(G,\Z)=H_{0}(\Z)_{G}=\Z$，且 $H_{1}(G,\Z[G])=0$，因为
\[
\Z[G]_{hG}:=\Z[G]\otimes_{\Z[G]}P_{G}=P_{G}\simeq\Z
\]
是拟同构。于是我们的正合列实际上是
\[
0\to H_{1}(G,\Z)\xrightarrow{\ \sim\ }H_{0}(G,I_{G})\to\Z\xrightarrow{\ \sim\ }\Z,
\]
这给出所说的同构。要点是
\[
H_{1}(G,\Z)=(I_{G})_{G}=I_{G}/I_{G}^{2},
\]
因为 $M_{G}=M/I_{G}\cdot M$。

\textbf{断言 13.12.} 映射
\[
\Z[G]/I_{G}^{2}\to G^{\mathrm{ab}}\times\Z,\qquad g\mapsto(\bar g,1)
\]
是同构。

这将推出 $I_{G}/I_{G}^{2}=\ker(\varepsilon)/I_{G}^{2}=G^{\mathrm{ab}}$，正是所要。

\textit{证明.} 首先注意上述映射是同态。事实上，记 $[g]\in\Z[G]$ 为 $g$ 的类，则对任意 $g,h\in G$ 有
\[
[g]+[h]\mapsto(\bar g\bar h,2),\qquad [g]\mapsto(\bar g,1),\qquad [h]\mapsto(\bar h,1),
\]
而后两个像之和恰为第一个。我们断言该映射有逆，由
\[
G\times\Z\to\Z[G]/I_{G}^{2},\qquad (g,n)\mapsto[g]+n-1
\]
诱导。它是同态，因为
\[
([g]-1)([h]-1)=[gh]-[g]-[h]+1\in I_{G}^{2},
\]
从而
\[
([g]-1)+([h]-1)\equiv[gh]-1\pmod{I_{G}^{2}},
\]
正是所要。最后，二者互逆，因为
\[
(\bar g,1)\mapsto[g]+1-1=[g],\qquad [g]+n-1\mapsto(\bar g,1)(1,n-1)=(\bar g,n),
\]
正是所要。\hfill$\square$

这证明了该断言。

最后，我们定义范映射 $X_{hG}\xrightarrow{\ N\ }X^{hG}$ 为如下复合：
\[
X_{hG}=X\otimes_{\Z[G]}P_{G}\to X\otimes_{\Z[G]}\Z\to\operatorname{Hom}_{\Z[G]}(\Z,X)\to\operatorname{Hom}_{\Z[G]}(P_{G},X)=X^{hG},
\]
其中第二个映射经由逐项的范映射（用张量–同调伴随）。然后置
\[
X_{tG}:=\operatorname{hCoker}\bigl(X_{hG}\xrightarrow{\ N\ }X^{hG}\bigr),
\]
我们断言它在拟同构意义下推广了先前对循环群所得者，于是可以定义
\[
\widehat{H}^{i}(G,X):=H^{i}(X_{tG}).
\]

很快我们将证明：

\textbf{断言 13.13（lcft）.} 对有限群 $G$ 与局部域扩张 $L/K$，$P_{G}\to L^{\times}[2]$ 在 Tate 上同调上是同构。

这给出
\[
\widehat{H}^{-2}(G,\Z)\cong\widehat{H}^{0}(G,L^{\times})=K^{\times}/N(L^{\times}).
\]
我们有正合列
\[
0=H^{-2}(\Z_{hG})\to\widehat{H}^{-2}(G,\Z)\xrightarrow{\ \sim\ }\underbrace{H^{-1}(\Z_{hG})}_{H_{1}(G,\Z)=G^{\mathrm{ab}}}\to H^{-1}(\Z^{hG})=0,
\]
因为 $\Z_{hG}$ 位于非正次数。于是，对 Galois 群为 $G$ 的局部域扩张 $L/K$，有
\[
K^{\times}/N(L^{\times})\cong G^{\mathrm{ab}}.
\]

\addcontentsline{toc}{section}{第 14 讲　Tate 上同调与 $K^{\mathrm{unr}}$}
% 第 14 讲：Tate 上同调与 K^unr
% 源：src/lecture14.txt（230 行，含 2 处页眉）
\section*{第 14 讲\quad Tate 上同调与 $K^{\mathrm{unr}}$}

设 $G$ 是有限群，$X$ 是 $G$-模复形。设 $P\xrightarrow{\ \sim\ }\Z$ 是 $G$-模的投射复形（拟同构）。则
\begin{itemize}[leftmargin=2.2em,itemsep=1pt,topsep=2pt]
\item $X^{hG}:=\operatorname{Hom}_G(P,X)$ 是\textbf{同伦不变量}；
\item $X_{hG}:=P\otimes_{\Z[G]}X$ 是\textbf{同伦余不变商}；
\item $X_{tG}:=\operatorname{hCoker}\bigl(X_{hG}\xrightarrow{\ N\ }X^{hG}\bigr)$ 是 \textbf{Tate 复形}。
\end{itemize}
前两个构造保持拟同构，把无环复形送到无环复形，因为投射复形平坦。此外，回忆记号
\begin{itemize}[leftmargin=2.2em,itemsep=1pt,topsep=2pt]
\item $H^{i}(G,X):=H^{i}(X^{hG})$ 是群上同调；
\item $H_{i}(G,X):=H^{-i}(X_{hG})$ 是群同调；
\item $\widehat{H}^{i}(G,X):=H^{i}(X_{tG})$ 是 Tate 上同调。
\end{itemize}
最后一个构造在 $G$ 为循环群且 $X$ 集中在 $0$ 次时推广了我们先前的构造。

现在让我们考虑模（而非复形）的 Tate 上同调。设 $G$ 作用在 $M$ 上。则对一个阿贝尔群 $N$，给出映射 $N\to M^{G}$ 等同于给出映射 $f\colon N\to M$，使得对一切 $g\in G$ 有 $g\cdot f(x)=f(x)$。对偶地，给出映射 $M_{G}\to N$ 等同于给出映射 $f\colon M\to N$，使得 $f(g\cdot x)=f(x)$（这是因为余不变商是 $M$ 的商，而不变子模是子模）。于是，由于对一切 $g\in G$ 有 $N(g\cdot x)=N(x)$ 与 $g\cdot N(x)=N(x)$，这些泛性质给出交换图表
\[
\begin{CD}
M @>{\mathrm{id}}>> M\\
@VVV @VVV\\
M_{G} @>{N}>> M^{G}
\end{CD}
\]
其中 $N:=\sum_{g}g$，且范映射穿过不变子模与余不变商。注意若 $\#G$ 在 $M$ 中可逆，则范映射 $N$ 是同构。模仿 Tate 上同调的定义，我们得到
\[
M_{G}/N(M_{G})=M_{G}/N(M)=\widehat{H}^{0}(G,M),
\]
可见同调代数实际上比我们「通常的」代数表现更好！

现在问：对 $G$-模复形，范映射 $N$ 是什么？我们有典范复合
\[
X_{hG}=P\otimes_{\Z[G]}X\to\underbrace{\Z\otimes_{\Z[G]}X}_{\text{逐项余不变商}}\xrightarrow{\ N\ }\underbrace{\operatorname{Hom}_{G}(\Z,X)}_{\text{逐项不变子模}}\to\operatorname{Hom}_{G}(P,X)=X^{hG},
\]
其中最后一个映射经由拉回，而范映射逐项地由模上的范映射给出——由先前的构造我们知道它按预期作用（尽管它只在同伦意义下良定义，等等）。注意「逐项不变子模」取的是逐次的「朴素」不变子模，它不保持拟同构；「逐项余不变商」类似。总之，这给出一个映射，我们记之为 $N\colon X_{hG}\to X^{hG}$。

取集中在 $0$ 次的复形（一般地，对有下界的复形），同伦不变量把该复形进一步向右推；类似地，余不变商把它向左推。但 Tate 上同调同时做这两件事，所以结果无界，往往非常混乱。不过在若干较简单的情形可以计算，例如下面这个：

\textbf{命题 14.1.} 设 $M$ 是 $G$-模，看作集中在 $0$ 次的复形。则
\begin{enumerate}[label=(\arabic*),leftmargin=2.6em,itemsep=1pt,topsep=2pt]
\item 若 $i\ge 1$，则 $\widehat{H}^{i}(G,M)=H^{i}(G,M)$；
\item $\widehat{H}^{0}(G,M)=M^{G}/N(M)$；
\item $\widehat{H}^{-1}(G,M)=\operatorname{Ker}(N)/(g-1)M=\operatorname{Ker}\bigl(N\colon M_{G}\to M^{G}\bigr)$；
\item 若 $i\ge 2$，则 $\widehat{H}^{-i}(G,M)=H_{i-1}(G,M)$。\footnote{原文此处作 $H^{i-1}(G,M)$；按证明中的长正合列应为同调群 $H_{i-1}(G,M)$，疑为笔误。}
\end{enumerate}

\textit{证明.} 复合
\[
M_{hG}\xrightarrow{\ N\ }M^{hG}\to M_{tG}=\operatorname{hCoker}(N)
\]
给出上同调长正合列
\[
\cdots\to H^{i}(M_{hG})\to H^{i}(M^{hG})\to\widehat{H}^{i}(G,M)\to H^{i+1}(M_{hG})\to\cdots.
\]
若 $i\ge 1$，则 $H^{i}(M_{hG})$ 与 $H^{i+1}(M_{hG})$ 均消没（因为 $M_{hG}$ 位于非正次数），由正合性得同构 $H^{i}(G,M)\cong\widehat{H}^{i}(G,M)$。至于 $i=0,-1$，此时 $H^{-1}(M^{hG})=0$，给出正合列
\[
0\to\widehat{H}^{-1}(G,M)\to M_{G}\xrightarrow{\ N\ }M^{G}\to\widehat{H}^{0}(G,M)\to 0,
\]
这证明了 (2) 与 (3)。最后，若 $i\ge 2$，则 $H^{-i}(M^{hG})$ 与 $H^{-i+1}(M^{hG})$ 均消没，由正合性得同构
\[
\widehat{H}^{-i}(G,M)\cong H^{-i+1}(M_{hG})=H_{i-1}(G,M),
\]
这证明了 (4)。\hfill$\square$

作为推论，上同调作为 Tate 上同调出现在更高次，但不在零次；类似地，同调也出现，唯独（关键地）在 $0$ 次例外。当然，这一切都依赖于 $P$ 有界这一事实。

\textbf{定理 14.2（上同调类域论主定理）.} 设 $L/K$ 是非阿基米德局部域的扩张，Galois 群 $G$ 有限。则
\[
(L^{\times})_{tG}\simeq(\Z[-2])_{tG}.
\]

虽然这个同构怎么构造并非一目了然，但其实并不太复杂！下一讲我们将把它归结为一个（非常典范的！）消没性陈述。

那么，这个定理究竟意味着什么？取零阶上同调，得到
\[
K^{\times}/N(L^{\times})=\widehat{H}^{0}(G,L^{\times})=H^{0}\bigl((L^{\times})_{tG}\bigr)\simeq H^{0}\bigl((\Z[-2])_{tG}\bigr)=\widehat{H}^{-2}(G,\Z)=H_{1}(G,\Z)\cong G^{\mathrm{ab}},
\]
正如上一讲所证。我们已知这对局部域的 $2$ 次扩张成立，所以这为局部域的希尔伯特符号提供了巨大的推广！

我们现已回忆同构 $H_{1}(G,\Z)\cong G^{\mathrm{ab}}$ 的构造。首先，我们证明了 $i>0$ 时 $H_{i}(G,\Z[G])=0$。事实上，
\[
\Z[G]\otimes_{\Z[G]}P\simeq\Z[G]\otimes_{\Z[G]}\Z=\Z
\]
是拟同构，因为 $P$ 拟同构于 $\Z$。特别地，所有低阶同调群消没，因为 $\Z$ 在 $0$ 次。然后我们构造了如下短正合列：
\[
0\to I_{G}\to\Z[G]\xrightarrow{\ \varepsilon\ }\Z\to 0,
\]
其中 $\varepsilon$ 由 $g\mapsto 1$ 定义，且我们令 $I_{G}:=\ker(\varepsilon)$，它是 $\Z[G]$ 中所有系数之和为零的元素构成的理想。取群同调长正合列，得到
\[
\underbrace{H_{1}(G,\Z[G])}_{=0}\to H_{1}(G,\Z)\to I_{G}/I_{G}^{2}\to\Z\xrightarrow{\ \sim\ }\Z\to 0,
\]
因为一般地 $M_{G}=M/I_{G}M$，故有同构 $H_{1}(G,\Z)\cong I_{G}/I_{G}^{2}$。

最后，我们通过双向构造的映射证明 $I_{G}/I_{G}^{2}\cong G^{\mathrm{ab}}$。我们断言映射
\[
G\longrightarrow I_{G}/I_{G}^{2},\qquad g\mapsto g-1
\]
在模 $I_{G}^{2}$ 意义下是同态（因而诱导出 $G^{\mathrm{ab}}\to I_{G}/I_{G}^{2}$）。为看出这一点，对 $g,h\in G$ 展开
\[
(g-1)(h-1)=gh-g-h+1\in I_{G}^{2},
\]
于是模 $I_{G}^{2}$ 有 $(g-1)+(h-1)\equiv gh-1$；由于 $I_{G}/I_{G}^{2}$ 交换，该映射穿过 $G^{\mathrm{ab}}$。至于逆，我们有复合
\[
I_{G}/I_{G}^{2}\hookrightarrow\Z[G]/I_{G}^{2}\xrightarrow{\ \sum_{g}n_{g}g\ \mapsto\ \prod_{g}g^{n_{g}}\ }G^{\mathrm{ab}},
\]
其中我们把 $G$ 写成乘法群；可以验证这是同态，且是前一映射的逆。

\textbf{例 14.3.} 设 $L/K$ 是局部域的非分歧扩张，故 $G:=\Gal(L/K)=\Z/n\Z$ 循环，且该同构是典范的，Frobenius 元素对应 $1$。回忆 $\widehat{H}^{0}(G,\OK_{L}^{\times})=0$，即 $N\colon\OK_{L}^{\times}\to\OK_{K}^{\times}$ 是满射。我们通过滤过证明了这一点：相继的子商给出范映射 $k_{L}^{\times}\xrightarrow{\ N\ }k_{K}^{\times}$，其余给出迹映射 $k_{L}\xrightarrow{\ T\ }k_{K}$，二者都是满射（例如后者是因为扩张可分）。我们还证明了 Herbrand 商 $\chi(\OK_{L}^{\times})=1$，这推出 $\widehat{H}^{1}(G,\OK_{L}^{\times})=0$，即 $(\OK_{L}^{\times})_{tG}\simeq 0$ 是拟同构。构造短正合列
\[
0\to\OK_{L}^{\times}\to L^{\times}\xrightarrow{\ v\ }\Z\to 0,
\]
其中 $v$ 表示 $L^{\times}$ 上的正规化赋值。取 Tate 上同调便给出拟同构
\[
\Z_{tG}[-2]\simeq(L^{\times})_{tG}\simeq\operatorname{hCoker}\bigl(0\to(L^{\times})_{tG}\bigr)\simeq\operatorname{hCoker}\bigl((\OK_{L}^{\times})_{tG}\to(L^{\times})_{tG}\bigr)\simeq\Z_{tG},
\]
由定理 14.2，且因为 Tate 上同调与锥交换并保持拟同构。于是 $(L^{\times})_{tG}$ 是 $2$-周期的。注意这是典范的，尽管需要选取一个生成元——因为我们可以选 Frobenius 元素（即 $x\mapsto x^{q}$）。

我们现在转向一般局部域扩张的讨论。

\textbf{定义 14.4.} $K^{\mathrm{unr}}$ 是 $K\subseteq\overline{K}$ 中极大非分歧扩张的（$p$ 进）完备化。

\textbf{例 14.5.}
\begin{enumerate}[label=(\arabic*),leftmargin=2.6em,itemsep=2pt,topsep=2pt]
\item 设 $K:=\F_{q}((t))$。$K$ 的第 $n$ 个非分歧扩张是 $K_{n}=\F_{q^{n}}((t))$，故 $K$ 的极大非分歧扩张是
\[
\bigcup_{n\ge 1}\F_{q^{n}}((t))\subseteq\overline{\F}_{q}((t))=K^{\mathrm{unr}}.
\]
\item 若 $K:=\Q_{p}$，则 $K^{\mathrm{unr}}=W(\overline{\F}_{p})[1/p]$，其中 $W(-)$ 表示维特向量环。\footnote{原文此处与下文均作 $\F_{q}$、$\F_{p}$；按文意应为代数闭包 $\overline{\F}_{q}$、$\overline{\F}_{p}$。}
\end{enumerate}

因此 $K^{\mathrm{unr}}$ 的基本结构是一个「局部域」（不是在局部紧的意义上，因为剩余域不有限，而是在它是一个完备离散赋值环的分式域的意义上），其剩余域为 $\overline{\F}_{q}$。

现在，设 $\pi$ 是 $K$ 的素元，它在每个 $K_{n}$（即 $K$ 的 $n$ 次非分歧扩张）中仍是素元，则 $\OK_{K^{\mathrm{unr}}}$ 是 $\bigcup_{n}\OK_{K_{n}}$ 的 $\pi$ 进完备化。于是有短正合列
\[
0\to\OK_{K^{\mathrm{unr}}}^{\times}\to K^{\mathrm{unr},\times}\xrightarrow{\ v\ }\Z\to 0,
\]
其中 $v$ 是 $K^{\mathrm{unr}}$ 上的正规化赋值（故 $v(\pi)=1$）。然后定义
\[
\Gal(K^{\mathrm{unr}}/K):=\operatorname{Aut}_{\mathrm{cts}}(K^{\mathrm{unr}}/K),
\]
后者是 $K^{\mathrm{unr}}$ 的连续 $K$-自同构。记 $k$ 为 $K$ 的剩余域、$\sigma$ 为 $\Gal(\bar k/k)$ 的 Frobenius 元素，则
\[
\Gal(K^{\mathrm{unr}}/K)\cong\Gal(\bar k/k)\cong\widehat{\Z},\qquad\sigma\mapsto 1.
\]

由于
\[
K=\{x\in K^{\mathrm{unr}}:x=\sigma x\},
\]
我们有消解
\[
0\to K\to K^{\mathrm{unr}}\xrightarrow{\ 1-\sigma\ }K^{\mathrm{unr}},
\]
进而有序列
\begin{equation}
0\to K^{\times}\to K^{\mathrm{unr},\times}\xrightarrow{\ x\mapsto x/\sigma x\ }K^{\mathrm{unr},\times}\xrightarrow{\ v\ }\Z\to 0. \tag{14.1}
\end{equation}
注意 $\pi$ 不可能落在中间映射的像中，因为对一切 $x$ 有 $v(x)=v(\sigma x)$。这给出用 $K^{\mathrm{unr},\times}$ 表达 $K^{\times}$ 的公式，它将是我们后续几讲的主要工具。

\textbf{断言 14.6.} 序列 (14.1) 正合。

\textit{证明.} 这等价于说每个 $x\in\OK_{K^{\mathrm{unr}}}^{\times}$ 都能写成 $y/\sigma y$，其中 $y\in\OK_{K^{\mathrm{unr}}}^{\times}$。这归结为证明映射
\[
\OK_{K^{\mathrm{unr}}}^{\times}\to\OK_{K^{\mathrm{unr}}}^{\times},\qquad x\mapsto x/\sigma x
\]
是满射。由极大理想的滤过的完备性（因为 $K^{\mathrm{unr}}$ 按定义完备），只需在伴随分次层面证明这一点。% [原文疑误：原印「第二式是 $\F_q\to\F_q$、$x\mapsto(1-q)x$，并称其『就是恒等映射』」。按伴随分次的实际推导，$\sigma$ 在 $\bar\F_q$ 上就是 Frobenius $x\mapsto x^q$，故该映射应为 $x\mapsto x-x^q$；原式 $(1-q)x$ 只是把它在 char $p$ 下形式地写成「标量 $1-q=1$」。此处径改为 $x-x^q$，满射性由 $\bar\F_q$ 代数闭保证]
这给出映射
\[
\overline{\F}_{q}^{\times}\to\overline{\F}_{q}^{\times},\quad x\mapsto x/x^{q}
\qquad\text{与}\qquad
\overline{\F}_{q}\to\overline{\F}_{q},\quad x\mapsto x-x^{q}.
\]
第一个是满射，因为对任意 $y\in\overline{\F}_{q}^{\times}$ 都可解 $x^{q-1}=1/y$，而 $\overline{\F}_{q}$ 代数闭。第二个也是满射，因为对任意 $a\in\overline{\F}_{q}$ 方程 $T^{q}-T=a$ 在代数闭域中总有解。\hfill$\square$

\addcontentsline{toc}{section}{第 15 讲　消没定理蕴含上同调类域论}
% 第 15 讲：消没定理蕴含上同调类域论
% 源：src/lecture15.txt（188 行，含 2 处页眉）
\section*{第 15 讲\quad 消没定理蕴含上同调类域论}

上一讲我们把问题重新表述为：对 Galois 群为 $G$ 的非阿基米德局部域扩张 $L/K$，证明
\begin{equation}
(L^{\times})_{tG}\simeq\Z_{tG}[-2]. \tag{15.1}
\end{equation}
于是我们的新目标是计算 $L^{\times}$ 的 Tate 上同调。回忆我们令 $K^{\mathrm{unr}}$ 表示 $K$ 的极大非分歧扩张的完备化；我们希望用 $K^{\mathrm{unr}}$ 来计算这个 Tate 上同调。

\textbf{断言 15.1.} 若 $x\in K^{\mathrm{unr}}$ 在 $K$ 上代数（由于完备化，未必如此），则 $K':=K(x)$ 在 $K$ 上非分歧。

\textit{证明.} 作为 $K$ 的有限代数扩张，$K'$ 是局部域，且有嵌入
\[
\OK_{K'}/\mathfrak{p}_{K}\OK_{K'}\hookrightarrow\OK_{K^{\mathrm{unr}}}/\mathfrak{p}_{K}\OK_{K^{\mathrm{unr}}}=\bar k,
\]
其中 $k:=\OK_{K}/\mathfrak{p}_{K}$。故 $\OK_{K'}/\mathfrak{p}_{K}\OK_{K'}$ 是域，从而 $K$ 与 $K'$ 的素元相同。\hfill$\square$

\textbf{断言 15.2.} $(K^{\mathrm{unr}})^{\sigma=1}=K$，即被 Frobenius 自同构 $\sigma\in G$ 固定（也就是在其作用下为恒等）的元素恰为 $K$；$\sigma$ 由各个非分歧扩张的 Frobenius 经连续性延拓到完备化而得到。

回忆我们有短正合列
\[
0\to K\to K^{\mathrm{unr}}\xrightarrow{\ 1-\sigma\ }K^{\mathrm{unr}},
\]
它在乘法群上可改写为
\[
1\to K^{\times}\to K^{\mathrm{unr},\times}\xrightarrow{\ x\mapsto x/\sigma x\ }K^{\mathrm{unr},\times}\xrightarrow{\ v\ }\Z\to 0.
\]
我们已证明：$K^{\mathrm{unr},\times}$ 中的元素若能写成 $x/\sigma x$，则它必是整数环中的单位 $x\in\OK_{K^{\mathrm{unr}}}^{\times}$；该映射在每个伴随分次项上是同构，从而在 $\OK_{K^{\mathrm{unr}}}^{\times}$ 上是同构。

现在，我们想显式构造 (15.1) 中的同构。第一次尝试如下：设 $L/K$ 完全分歧（既然上一讲已处理非分歧情形，这个假设相当温和），于是 $L^{\mathrm{unr}}=L\otimes_{K}K^{\mathrm{unr}}$。则我们有如下定理（稍后证明）。

\textbf{定理 15.3（消没定理）.} 若 $L/K$ 完全分歧，则复形 $(L^{\mathrm{unr},\times})_{tG}$ 无环。

\textbf{断言 15.4.} 消没定理蕴含上同调类域论。

\textit{证明.} 设 $L/K$ 完全分歧。我们有四项正合列
\begin{equation}
1\to L^{\times}\to L^{\mathrm{unr},\times}\xrightarrow{\ x\mapsto x/\sigma x\ }L^{\mathrm{unr},\times}\xrightarrow{\ v\ }\Z\to 0. \tag{15.2}
\end{equation}
可将其改写如下：
\[
\begin{array}{c|cccccc}
A & \cdots & 0 & L^{\times} & 0 & 0 & \cdots\\[2pt]
B & \cdots & 0 & L^{\mathrm{unr},\times} & L^{\mathrm{unr},\times} & 0 & \cdots\\[2pt]
\operatorname{Coker}(A\to B) & \cdots & 0 & 0 & \Z & 0 & \cdots
\end{array}
\]
其中 $L^{\times}$ 在 $-1$ 次，$B$ 的微分是 $1-\sigma$、$\operatorname{Coker}(A\to B)$ 的微分是 $v$。最后到由 (15.2) 所得同伦余核的拟同构来自断言 10.12，因为 $A\hookrightarrow B$ 是单射（注意这对任何四项正合列都成立）。逐项余核给出单射
\[
L^{\mathrm{unr},\times}/L^{\times}\xrightarrow{\ x\mapsto x/\sigma x\ }L^{\mathrm{unr},\times},
\]
因为去掉商之后，$L^{\times}$ 恰是该映射的核。

现在，我们有拟同构
\[
B_{tG}=\operatorname{hCoker}\bigl(L^{\mathrm{unr},\times}\xrightarrow{\ 1-\sigma\ }L^{\mathrm{unr},\times}\bigr)_{tG}\simeq\operatorname{hCoker}\bigl((L^{\mathrm{unr},\times})_{tG}\to(L^{\mathrm{unr},\times})_{tG}\bigr),
\]
于是由消没定理，$(L^{\mathrm{unr},\times})_{tG}$ 无环，从而由上同调长正合列，这个同伦余核也无环。因此
\[
(L^{\times}[2])_{tG}=\operatorname{hCoker}\bigl((L^{\times}[1])_{tG}\to 0\bigr)=\operatorname{hCoker}\bigl(A_{tG}\to B_{tG}\bigr)\simeq\Z_{tG},
\]
正是所要。\hfill$\square$

现在设 $L/K$ 是一般的有限 Galois 扩张，$G:=\Gal(L/K)$（虽然我们可以分别处理非分歧与完全分歧情形，因为任何扩张都典范地是这样的扩张的复合）。若 $L/K$ 非分歧，则
\[
L\otimes_{K}K^{\mathrm{unr}}=\prod_{L\hookrightarrow K^{\mathrm{unr}}}K^{\mathrm{unr}}
\]
典范地，按这些嵌入编号。事实上，下述结论成立：

\textbf{定理 15.5（一般消没定理）.} $\bigl[(L\otimes_{K}K^{\mathrm{unr}})^{\times}\bigr]_{tG}$ 无环。

为理解 $L\otimes_{K}K^{\mathrm{unr}}$ 的结构，注意 $\widehat{\Z}\sigma$ 作用在第二个因子、$G$ 作用在第一个因子；这两个作用（即 $x\otimes y\mapsto gx\otimes y$ 与 $x\otimes y\mapsto x\otimes\sigma y$）显然交换。同样，被 $\sigma$ 固定的点构成
\[
L=L\otimes_{K}K\hookrightarrow L\otimes_{K}K^{\mathrm{unr}}.
\]

\textbf{断言 15.6.} 下述序列正合：
\[
1\to L^{\times}\to(L\otimes_{K}K^{\mathrm{unr}})^{\times}\xrightarrow{\ x\mapsto x/\sigma x\ }(L\otimes_{K}K^{\mathrm{unr}})^{\times}\to\Z\to 0.
\]

\textit{证明.} 若 $x\in K^{\mathrm{unr}}$ 是单位，则 $\sigma x$ 也是，故映射 $x\mapsto x/\sigma x$ 良定义；并且 $x$ 在其核中当且仅当 $x$ 在 $\sigma$ 作用下不动，即 $x\in K$；又 $L\otimes_{K}K=L$，故得到 $L$ 的一个单位，这证明了左半部分的正合性。至于到 $\Z$ 的映射，由
\[
\begin{CD}
(L\otimes_{K}K^{\mathrm{unr}})^{\times} @>>> \Z\\
@V{N_{L/K}\otimes\operatorname{id}}VV @A{\ v\ }AA\\
K\otimes_{K}K^{\mathrm{unr},\times} @= K^{\mathrm{unr},\times}
\end{CD}
\]
定义，其中 $N_{L/K}(x):=\prod_{g\in G}gx$。于是其核为 $\OK_{K^{\mathrm{unr}}}^{\times}$，它恰是 $x\mapsto x/\sigma x$ 的像。此外该映射是满射，因为 $1\otimes\pi\mapsto 1$。\hfill$\square$

注意若 $L/K$ 完全分歧，这就回到我们先前的扩张。事实上，若记 $L^{\mathrm{unr}}=L\otimes_{K}K^{\mathrm{unr}}$，则两边的 $\sigma$「吻合」，即 $L^{\mathrm{unr}}$ 与 $K^{\mathrm{unr}}$ 上的诱导 Frobenius 自同构相同，因为 $L$ 与 $K$ 有相同的剩余域。该扩张的范映射 $N_{L/K}\colon L^{\mathrm{unr},\times}\to K^{\mathrm{unr},\times}$ 满足 $v_{K^{\mathrm{unr}}}\circ N=v_{L^{\mathrm{unr}}}$（这样的扩张由 $K$ 的素元的 $n$ 次根生成，于是 $N(\pi^{1/n})=\pi$）。

现在设 $L/K$ 是 $n$ 次非分歧扩张。固定一个嵌入 $L\hookrightarrow K^{\mathrm{unr}}$，并用 $\sigma\in\Gal(L/K)$ 同时表示 $L/K$ 的 Frobenius 元素。则我们有同构
\[
L\otimes_{K}K^{\mathrm{unr}}\xrightarrow{\ \sim\ }\prod_{i=0}^{n-1}K^{\mathrm{unr}},\qquad x\otimes y\mapsto\bigl((\sigma^{i}x)\cdot y\bigr)_{i=0}^{n-1},
\]
其中乘积经由我们固定的嵌入（注意也可以像先前那样对嵌入取乘积，从而更典范地做到这一点）。现在问：$L\otimes_{K}K^{\mathrm{unr}}$ 的自同构 $\mathrm{id}\otimes\sigma$ 在该同构下对应什么？我们有
\[
x\otimes\sigma y\mapsto(x\cdot\sigma y,\ \sigma x\cdot\sigma y,\ \sigma^{2}x\cdot\sigma y,\ \dots)=\sigma\bigl(\sigma^{-1}x\cdot y,\ x\cdot y,\ \sigma x\cdot y,\ \dots\bigr),
\]
所以它是 $\sigma$ 作用在 $x\otimes y$ 的像的「右旋」上（注意 $\sigma$ 在 $K^{\mathrm{unr}}$ 上不是有限阶的，$\mathrm{id}\otimes\sigma$ 也不应是有限阶的，这就排除了把 $\mathrm{id}\otimes\sigma$ 的像写成旋转的可能性）。类似地，范映射 $N_{L/K}\colon\prod K^{\mathrm{unr},\times}\to K^{\mathrm{unr},\times}$ 取所有分量的乘积。

我们希望某个元素 $(x_{0},\dots,x_{n-1})\in\prod K^{\mathrm{unr},\times}$ 落在 $y/\sigma y$ 的像中（即断言 15.6 正合列中间那个映射；这里 $\sigma$ 指自同构 $\mathrm{id}\otimes\sigma$），当且仅当 $\prod x_{i}\in\OK_{K^{\mathrm{unr}}}^{\times}$，即 $\sum v(x_{i})=0$。回忆反向蕴含是平凡的，因为我们已经证明 $\OK_{K^{\mathrm{unr}}}^{\times}\xrightarrow{\ y/\sigma y\ }\OK_{K^{\mathrm{unr}}}^{\times}$ 是满射，因为它在伴随分次层面满射。至于正向，我们有
\[
(y_{0},\dots,y_{n-1})\xrightarrow{\ y/\sigma y\ }\Bigl(\frac{y_{0}}{\sigma y_{n-1}},\frac{y_{1}}{\sigma y_{0}},\dots\Bigr)=:(x_{0},x_{1},\dots).
\]
于是
\[
y_{0}=x_{0}\cdot\sigma y_{n-1},
\]
\[
y_{1}=x_{1}\cdot\sigma y_{0}=x_{1}\cdot\sigma x_{0}\cdots\sigma^{2}y_{n-1},
\]
\[
\cdots=\cdots
\]
\[
y_{n-1}=x_{n-1}\cdot\sigma x_{n-2}\cdots\sigma^{n-1}x_{0}\cdot\sigma^{n}y_{n-1},
\]
即
\[
\frac{y_{n-1}}{\sigma^{n}y_{n-1}}=x_{n-1}\cdot\sigma x_{n-2}\cdots\sigma^{n-1}x_{0}.
\]
注意这里一切都是 $K^{\mathrm{unr}}$ 的元素，所以我们确实没有 $\sigma^{n}=\mathrm{id}$！上一讲我们证明了这当且仅当右端属于 $\OK_{K^{\mathrm{unr}}}^{\times}$，而这等价于说 $\sum v(x_{i})=0$。该正合列的一般情形是这两种情形的某种混合。

我们现在把这些结果与上一讲的结果作比较。假定消没定理成立。对非分歧扩张 $L/K$，我们有两个从 $(L^{\times})_{tG}$ 到 $\Z[-2]_{tG}$ 的拟同构：一个来自我们刚刚所做，另一个来自 $(\OK_{L}^{\times})_{tG}\simeq 0$ 蕴含（由循环性）$(L^{\times})_{tG}\simeq\Z_{tG}\simeq(\Z[-2])_{tG}$。我们断言这两个拟同构重合。证明的思路如下：$G=\Z/n\Z$（生成元为 Frobenius 元素），且有短正合列
\[
0\to\Z\to\Z[G]\xrightarrow{\ 1-\sigma\ }\Z[G]\to\Z\to 0.
\]
如习题集 7 第 1(e) 题所示，$\Z[G]_{tG}\simeq 0$ 是拟同构（这很容易证明，而且我们已经对循环群证明过）。于是得到 $\Z_{tG}[2]\simeq\Z_{tG}$，我们断言这与由复形的 $2$-周期性得到的同构相同。证明靠图表追踪。我们有 $(L\otimes_{K}K^{\mathrm{unr}})^{\times}=\prod K^{\mathrm{unr},\times}$，是有限乘积。于是图表
\[
1\longrightarrow L^{\times}\longrightarrow(L\otimes_{K}K^{\mathrm{unr}})^{\times}\xrightarrow{\ x\mapsto x/\sigma x\ }(L\otimes_{K}K^{\mathrm{unr}})^{\times}\xrightarrow{\ v\ }\Z\longrightarrow 0
\]
与
\[
1\longrightarrow\Z\longrightarrow\Z[G]\xrightarrow{\ 1-\sigma\ }\Z[G]\xrightarrow{\ \varepsilon\ }\Z\longrightarrow 0
\]
竖直映射依次取 $v$、$\prod v$、$\sum v$ 时交换，其中 $\varepsilon$ 表示 $\Z[G]$ 各坐标之和。这恰好说明由两个四项正合列得到的同构相同。

要点是：在 lcft 之下，我们有同构
\[
K^{\times}/N(L^{\times})\cong\Z/n\Z,\qquad\pi\mapsto\mathrm{Frob}.
\]
于是我们把 lcft 归结为消没定理，后者我们将在下一讲证明。

\addcontentsline{toc}{section}{第 16 讲　Tate 上同调群的消没}
% 第 16 讲 Tate 上同调群的消没（译自 src/lecture16.txt）
\section*{第 16 讲\quad Tate 上同调群的消没}

前情回顾：我们把（上同调的）局部类域论归约成了如下命题：设 $L/K$ 是非阿基米德局部域的有限 Galois 扩张，Galois 群为 $G$，则
\[
\big((L\otimes_K K^{\mathrm{unr}})^\times\big)_{tG}\simeq 0,
\]
即该复形无环。为证明这一消没，我们将证明一个（属于 Tate 的）关于 Tate 上同调消没的一般定理，它使上述问题变得更加易于处理。于是我们问：给定 $G$-模的复形 $X$，什么条件能保证 $X_{tG}$ 无环？这类结果的原型如下：

\textbf{定理 16.1.} 若 $G$ 是循环群，则 $X_{tG}$ 无环当且仅当
\[
\widehat H^0(G,X)=0=\widehat H^1(G,X).
\]

\textit{证明.} $X_{tG}$ 可由一个 $2$-周期复形来计算。注意任意两个奇偶性不同的值（例如相邻的两个值）就足够了。\hfill$\square$

本讲的主要结果是：

\textbf{定理 16.2.} 若 $G$ 是 $p$-群（即 $\#G=p^n$，其中 $p$ 为某个素数、$n\ge 0$），定理 16.1 仍然成立。

由此我们将推出下一个结果：

\textbf{定理 16.3.} 设对每个素数 $p$ 及 $G$ 的每个 $p$-Sylow 子群 $G_p\subseteq G$，都有 $\widehat H^0(G_p,X)=0=\widehat H^1(G_p,X)$。则 $X_{tG}$ 无环。

\textbf{注 16.4.} 一般而言，对任意的有限群 $G$，「在相邻两个次数上消没」并不充分。此外，实践中人们往往是对 $G$ 的每个子群（而不只是 $p$-Sylow 子群）验证相邻两个次数上 Tate 上同调的消没。

在接下来的几讲中，我们将由此推出局部类域论。我们先证明定理 16.2。以下始终设 $G$ 为 $p$-群。

\textbf{命题 16.5.} 设 $X$ 是 $\mathbb F_p[G]$-模的复形。若 $\widehat H^0(G,X)=0$，则 $X_{tG}$ 无环。

注意这里我们只需要一个次数上的消没！为此，我们先回忆下述事实。

\textbf{引理 16.6.} 设 $V$ 是非零的 $\mathbb F_p[G]$-模。则 $V^G\neq 0$。

\textit{证明.} 不妨设 $V$ 在 $\mathbb F_p$ 上有限维，因为 $V$ 显然是它的有限维 $G$-子模的并。于是 $\#V=p^r$，其中 $0<r<\infty$。$V$ 的每个 $G$-轨道的大小或者为 $1$、或者被 $p$ 整除（轨道大小必定整除 $\#G$，因为该轨道同构于 $G$ 对稳定化子的商）。由于 $\{0\}$ 是一个大小为 $1$ 的 $G$-轨道，必存在另一个大小为 $1$ 的轨道（因为所有 $G$-轨道的大小之和必须等于 $p^r$），即存在某个 $v\in V^G\setminus\{0\}$。\hfill$\square$

\textit{证明（命题 16.5）.} \textbf{第 1 步.} 先断言：对每个有限维 $\mathbb F_p$-表示 $V$，都有 $\widehat H^0(G,X\otimes V)=0$。这里 $G$ 通过「对角作用」作用，即在 $(X\otimes V)^i=X^i\otimes V$ 上由 $g\cdot(x\otimes v):=gx\otimes gv$ 给出。我们对 $\dim_{\mathbb F_p}V$ 作归纳。由前一引理，有 $\mathbb F_p[G]$-模的短正合列
\[
0\to V^G\to V\to W\to 0,
\]
其中 $\dim W<\dim V$。这给出
\[
\operatorname{hCoker}(V^G\otimes X\to V\otimes X)\simeq W\otimes X.
\]
由于由归纳假设 $\widehat H^0(G,W\otimes X)=0$，而
\[
\widehat H^0(G,V^G\otimes X)=\widehat H^0\Big(G,\bigoplus_{\dim V^G}X\Big)=\bigoplus_{\dim V^G}\widehat H^0(G,X)=0
\]
（由对 $X$ 的假设），Tate 上同调的长正合列便给出 $\widehat H^0(G,V\otimes X)=0$，如所欲证。

\textbf{第 2 步.} 现在证明负次数上的消没。考虑短正合列
\[
0\to V_1\to \mathbb F_p[G]\xrightarrow{\ \varepsilon\ }\mathbb F_p\to 0,
\]
其中 $V_1$ 定义为 $\varepsilon$ 的核，这与我们把 $\mathbb F_p$ 换成 $\mathbb Z$ 时所说的 $I_G$ 类似。设 $\underline X$ 表示 $X$ 配上平凡的 $G$-作用。则
\[
\underline X\otimes\mathbb F_p[G]\to X\otimes\mathbb F_p[G],\qquad x\otimes g\mapsto gx\otimes g
\]
（右端取对角作用）是 $G$-等变映射，并且是双射，从而是同构。事实上，
\[
h(x\otimes g)=x\otimes hg\mapsto hgx\otimes hg=h(gx\otimes g).
\]
习题集 7 第 1(e) 题已经证明 $(\underline X\otimes\mathbb F_p[G])_{tG}$ 无环，故 $(X\otimes\mathbb F_p[G])_{tG}$ 亦然。于是 Tate 上同调的长正合列给出
\[
\widehat H^{i-1}(G,X)\simeq \widehat H^i(G,X\otimes V_1).
\]
第 1 步已表明右端在 $i=0$ 时消没，因此 $\widehat H^{-1}(G,X)=0$。迭代即得 $\widehat H^i(G,X)=0$ 对一切 $i\le 0$ 成立。

\textbf{第 3 步.} 为证明正次数上的消没，注意我们有正合列
\[
0\to \mathbb F_p\xrightarrow{\ 1\mapsto\sum_{g\in G}g\ }\mathbb F_p[G]\to V_2\to 0,
\]
其中 $V_2$ 如前面一样定义为余核。同样的论证给出
\[
\widehat H^i(G,X\otimes V_2)\simeq \widehat H^{i+1}(G,X),
\]
于是第 1 步再次表明 $\widehat H^i(G,X)=0$ 对一切 $i\ge 0$ 成立。\hfill$\square$

\textit{证明（定理 16.2）.} 定义 $X/p:=\operatorname{hCoker}(X\xrightarrow{p}X)$；注意这与把每一项都对 $p$ 取商并不相同。注意，作为 $\mathbb Z[G]$-模的复形，$X/p$ 拟同构于一个 $\mathbb F_p[G]$-模的复形。由于 $X$ 只在拟同构的意义下有定义，不妨设它作为 $\mathbb Z[G]$-模的复形是投射的（特别地，平坦的）。于是我们有拟同构
\[
X/p=X\otimes_{\mathbb Z[G]}\operatorname{hCoker}(\mathbb Z[G]\xrightarrow{p}\mathbb Z[G])\simeq X\otimes_{\mathbb Z[G]}\mathbb F_p[G],
\]
其中右端逐项计算 $X$ 模 $p$ 的商。于是有长正合列
\[
\cdots\to \widehat H^i(G,X)\xrightarrow{p}\widehat H^i(G,X)\to \widehat H^i(G,X/p)\to \widehat H^{i+1}(G,X)\xrightarrow{p}\widehat H^{i+1}(G,X)\to\cdots,
\]
令 $i=0$，由假设即得 $\widehat H^0(G,X/p)=0$。于是由命题 16.5，$(X/p)_{tG}$ 无环。由此 $\widehat H^i(G,X/p)=0$ 对一切 $i$ 成立，因此乘以 $p$ 在 $\widehat H^i(G,X)$ 上是同构，对每个 $i$ 皆然。但正如习题集 7 第 2(c) 题所示，乘以 $\#G$ 在 $\widehat H^i(G,X)$ 上是零映射。由于 $G$ 是 $p$-群，这只有在 $\widehat H^i(G,X)=0$ 对一切 $i$ 成立时才可能。\hfill$\square$

\textit{证明（定理 16.3）.} 如上所述，乘以 $\#G$ 在 $\widehat H^i(G,X)$ 上是零映射，故 $\widehat H^i(G,X)$ 是 $\#G$-挠的。因此只需证明 $\widehat H^i(G,X)[p]=0$ 对一切 $p$ 成立（即 $\widehat H^i(G,X)$ 的 $p$-挠部分消没）。

回忆对 $G$ 的每个子群 $H$，有分别称为限制与余限制的映射 $X_{tG}\to X_{tH}$ 与 $X_{tH}\to X_{tG}$，其复合作为 $X_{tG}$ 的自同态同伦于乘以指数 $[G:H]$。% [原文疑误：原印作 restriction 与 inflation（膨胀），但复合为乘以指数 $[G:H]$ 的正是「限制与余限制」（corestriction），且膨胀映射只对正规子群有定义，故此处径改为余限制]

对 $p$-Sylow 子群 $H=G_p$ 应用上述结论并取上同调，我们得到映射
\[
\widehat H^i(G,X)[p]\subset \widehat H^i(G,X)\to \widehat H^i(G_p,X)\to \widehat H^i(G,X),
\]
其复合是乘以 $[G:G_p]$，而后者按定义与 $p$ 互素。因此它在限制到 $\widehat H^i(G,X)[p]$ 上是同构，特别地 $\widehat H^i(G,X)[p]\to\widehat H^i(G_p,X)$ 是单射。但由定理 16.2，$\widehat H^i(G_p,X)=0$ 对一切 $i$ 成立，于是得到所欲结论。\hfill$\square$

\addcontentsline{toc}{section}{第 17 讲　消没定理的证明}
% 第 17 讲 消没定理的证明（译自 src/lecture17.txt）
\section*{第 17 讲\quad 消没定理的证明}

本讲的目标是证明：对非阿基米德局部域的扩张 $L/K$，设 Galois 群为 $G$，则
\[
\big[(L\otimes_K K^{\mathrm{unr}})^\times\big]_{tG}\simeq 0.
\]
回忆这蕴涵 $(L^\times)_{tG}\simeq \mathbb Z_{tG}[-2]$（即上同调局部类域论的主定理），后者又蕴涵 $G^{\mathrm{ab}}\simeq K^\times/NL^\times$。目前我们先假设 $L/K$ 完全分歧（一般情形到后面再归约），此时有 $L^{\mathrm{unr}}=L\otimes_K K^{\mathrm{unr}}$。上一次我们证明了：只需对每个素数 $\ell$ 与 $G$ 的每个 $\ell$-Sylow 子群 $G_\ell\subseteq G$ 证明
\[
\widehat H^0(G_\ell,L^{\mathrm{unr},\times})=0=\widehat H^1(G_\ell,L^{\mathrm{unr},\times}).
\]
注意，若令 $K':=L^{G_\ell}$，则 $L/K'$ 是 $G_\ell$-Galois 扩张。于是我们可以把 $K$ 换成 $K'$、把 $G$ 换成 $G_\ell$，从而不妨直接假设 $G$ 是 $\ell$-群（即 $\#G=\ell^n$，其中 $n$ 为某个整数）。

现在，上述两个等式中后一个不过是 Hilbert 定理 90（或习题集 7 第 3 题所给的推广）施用于扩张 $L^{\mathrm{unr}}/K^{\mathrm{unr}}$，因此还需证明前一个，即范映射
\[
N:L^{\mathrm{unr},\times}\to K^{\mathrm{unr},\times}
\]
是满射。

我们回忆 $\ell$-群的结构理论：

\textbf{命题 17.1.} 设 $G$ 是 $\ell$-群。则存在正规子群链
\[
1\lneq G_0\lneq\cdots\lneq G_m=G,
\]
使得 $G_{i+1}/G_i$ 对每个 $i$ 都是循环群。

\textit{证明.} 主要一步是证明：若 $G\neq 1$，则 $Z(G)\neq 1$（即 $G$ 的中心非平凡）。让 $G$ 通过伴随作用作用在自身之上，即对 $g,x\in G$ 令 $g\cdot x:=gxg^{-1}$。于是每个 $G$-轨道的大小或者为 $1$、或者被 $\ell$ 整除。由于
\[
\sum_{O\in G\text{-轨道}}\#O=\#G=\ell^n>1,
\]
而 $1$ 的 $G$-轨道大小为 $1$，故 $\ell$ 必须整除那些大小为 $1$ 的 $G$-轨道的个数，于是 $\ell\mid \#Z(G)$（特别地 $\#Z(G)\neq 0$）。% [原文作「hence #Z(G) ≠ 0」，按前文应为 ℓ | #Z(G)；照前文译]
然后取非平凡元素 $x\in Z(G)$，可见 $\langle x\rangle$ 是 $G$ 的正规子群，从而可对商群 $G/\langle x\rangle$ 作归纳即得结论。% [原文作「G/⟨x⟩ is a normal subgroup of G」，商群非子群；已按「⟨x⟩ 正规，对商群归纳」译]\hfill$\square$

于是由 Galois 理论，我们有一列相应的循环扩张
\[
L=L_m/L_{m-1}/\cdots/L_0=K.
\]
由于只需在这些子扩张的每一个上证明范映射是满射（满射的复合仍是满射），我们可以假设 $G$ 是循环群，比如说阶为 $n$。回忆 $N:L^\times\to K^\times$ 并非满射，因为我们已证明 $\#\widehat H^0(G,L^\times)=n$。现在，对每个 $m$，记 $K\subseteq K_m\subseteq K^{\mathrm{unr}}$ 为 $K$ 的 $m$ 次非分歧扩张。主要的一步如下：

\textbf{断言 17.2.} 设 $x\in K^\times$，并记 $L_m:=LK_m$。则对满足 $n\mid m$ 的 $m$，$x$ 属于 $N:L_m^\times\to K_m^\times$ 的像。% [原文只写「对每个 m」且未定义 L_m；按上下文（K_m 为 m 次非分歧扩张、下文 $N_{L_m/L}$）应取 L_m=L K_m。又因 $K^\times/NL^\times$ 的阶为 n，下面复合映射的「取 m 次幂为零」需 n | m（原文作取 n 次幂，见下），故补上此条件；下文只用到 m→∞ 的极限，无碍]

\textit{证明.} 注意 $K^\times/NL^\times=\mathcal O_K^\times/N\mathcal O_L^\times$。事实上，我们有通常的短正合列
\[
0\to \mathcal O_L^\times\to L^\times\xrightarrow{\ v\ }\mathbb Z\to 0,
\]
它给出正合列
\[
0=\widehat H^{-1}(\mathbb Z)\to \widehat H^0(\mathcal O_L^\times)\hookrightarrow \widehat H^0(L^\times)\to \widehat H^0(\mathbb Z)=\mathbb Z/n\mathbb Z,
\]
而最右端的映射为零，因为 $L/K$ 完全分歧（从而对一切 $y\in K^\times$ 有 $n\mid v(y)$）。于是得到同构 $\widehat H^0(L^\times)\simeq \widehat H^0(\mathcal O_L^\times)$，这正是我们的观察。

我们有交换图
\[
\begin{CD}
L^\times @>N>> K^\times @>>> K^\times/NL^\times @>>> 0\\
@VVV @VVV @VVV\\
L_m^\times @>N>> K_m^\times @>>> K_m^\times/NL_m^\times @>>> 0\\
@V{N_{L_m/L}}VV @V{N_{K_m/K}}VV @V{}VV\\
L^\times @>>> K^\times @>N>> K^\times/NL^\times @>>> 0.
\end{CD}
\]

现在，由诱导映射复合而成的映射 $K^\times/NL^\times\to K_m^\times/NL_m^\times\to K^\times/NL^\times$ 是取 $m$ 次幂（$m=[K_m:K]$），因此当 $n\mid m$ 时它是零映射。% [原文作「取 n 次幂」，但按上图复合（包含 $K^\times\hookrightarrow K_m^\times$ 接范映射 $N_{K_m/K}$）实为 $x\mapsto x^{[K_m:K]}=x^m$；又 $K^\times/NL^\times$ 的阶为 n，故须 n | m 才为零，原文此处的 n 与「对每个 m」皆有疏漏，已按数学含义订正] 我们想证明诱导映射
\[
N_{K_m/K}:K_m^\times/NL_m^\times\to K^\times/NL^\times
\]
是同构，而这蕴涵诱导映射 $K^\times/NL^\times\to K_m^\times/NL_m^\times$ 为零，从而证明断言。由断言 7.8(3)，即我们先前对 Herbrand 商的分析，两个群的阶都是 $n$，因此该映射单射当且仅当它满射。此外，由我们的观察，这等价于映射
\[
N_{K_m/K}:\mathcal O_{K_m}^\times/N\mathcal O_{L_m}^\times\to \mathcal O_K^\times/N\mathcal O_L^\times,
\]
而由于 $N:\mathcal O_{K_m}^\times\to \mathcal O_K^\times$ 是满射（其证明与断言 3.4 的证明完全相同），该映射也是满射，证毕。\hfill$\square$

再次设 $G$ 为 $n$ 阶循环群，我们只需证明 $N:L^{\mathrm{unr},\times}\to K^{\mathrm{unr},\times}$ 是满射。把断言施用于 $K_m^\times$，可见 $K_m^\times$ 的每个元素都是 $L_{m+m'}^\times$ 中某个元素的范数，因此
\[
N:\bigcup_m L_m^\times\to \bigcup_m K_m^\times
\]
是满射。接下来只需过渡到完备化。我们知道映射 $N:L^{\mathrm{unr},\times}\to K^{\mathrm{unr},\times}$ 的像包含 $\bigcup_m K_m^\times$，而后者稠密，故只需证明其像包含 $1$ 的一个开邻域。显然
\[
N(L^{\mathrm{unr},\times})\supseteq N(K^{\mathrm{unr},\times})=(K^{\mathrm{unr},\times})^n,
\]
而我们在习题集 1 第 1(a) 题中已经看到，$1+\mathfrak p_{K^{\mathrm{unr}}}^{2v(n)+1}$ 的每个元素都是 $K^{\mathrm{unr}}$ 中的 $n$ 次幂；这就是所要的开邻域。

最后，我们证明消没定理的一般情形，此时非阿基米德局部域的 $G$-扩张 $L/K$ 未必完全分歧。设 $L/L_0/K$ 为 $L$ 中 $K$ 的（唯一）极大非分歧扩张，于是 $L/L_0$ 完全分歧。令 $H:=\Gal(L/L_0)$，则 $L_0/K$ 是 Galois 扩张，群为 $G/H$。

\textbf{引理 17.3.} 设 $X$ 是 $G$-模的复形，并设有正合列
\[
1\to H\to G\to G/H\to 1.
\]
若
\[
X_{tH}\simeq 0\simeq (X^{hH})_{tG/H},
\]
则 $X_{tG}\simeq 0$。

注意，一般而言 $(X_{tH})_{tG/H}=X_{tG}$ 并不成立！例如，若 $H$ 是平凡群，则
\[
X_{tH}=\operatorname{hCoker}(N:X\to X)=0,
\]
其中此处 $N=\mathrm{id}_X$。

\textit{证明.} 由第一个条件，$X_{hH}\xrightarrow{qis}X^{hH}$；由习题集 6 第 3 题，有 $X^{hG}\simeq(X^{hH})^{hG/H}$ 与 $X_{hG}\simeq(X_{hH})_{hG/H}$。于是由第二个条件，
\[
(X^{hH})_{hG/H}\xrightarrow{qis}(X^{hH})^{hG/H}.
\]
容易验证这个拟同构由范映射给出（它是两个范映射的复合），这蕴涵
\[
X_{tG}=\operatorname{hCoker}(X^{hG}\to X^{hG})
\]
无环，如所欲证。\hfill$\square$

现在，我们想证明 $[(L\otimes_K K^{\mathrm{unr}})^\times]_{tG}\simeq 0$。回忆
\[
L\otimes_K K^{\mathrm{unr}}=L\otimes_{L_0}L_0\otimes_K K^{\mathrm{unr}}=L\otimes_{L_0}\prod_{L_0\hookrightarrow K^{\mathrm{unr}}}K^{\mathrm{unr}}=\prod_{L_0\hookrightarrow K^{\mathrm{unr}}}L\otimes_{L_0}K^{\mathrm{unr}}
\]
是典范地成立的（第二个同构由映射 $\alpha\otimes\beta\mapsto (i(\alpha)\beta)_i$ 给出，下标取遍嵌入 $i:L_0\hookrightarrow K^{\mathrm{unr}}$）。我们有
\[
\big[(L\otimes_K K^{\mathrm{unr}})^\times\big]_{tH}\simeq \prod_{L_0\hookrightarrow K^{\mathrm{unr}}}\big[(L\otimes_{L_0}K^{\mathrm{unr}})^\times\big]_{tH}\simeq \prod_{L_0\hookrightarrow K^{\mathrm{unr}}}\big[(L\otimes_{L_0}L_0^{\mathrm{unr}})^\times\big]_{tH}\simeq 0
\]
（用完全分歧情形，因为 $L_0/K$ 非分歧而 $L/L_0$ 完全分歧），这就建立了引理的第一个条件。为验证第二个条件，注意
\[
\prod_{L_0\hookrightarrow K^{\mathrm{unr}}}\big[(L\otimes_{L_0}K^{\mathrm{unr}})^\times\big]^{hH}=\prod_{L_0\hookrightarrow K^{\mathrm{unr}}}K^{\mathrm{unr},\times}\simeq K^{\mathrm{unr},\times}[G/H]
\]
作为 $G/H$-模成立（一旦固定一个嵌入 $L_0\hookrightarrow K^{\mathrm{unr}}$）。但正如习题集 7 第 1(e) 题所示，Tate 上同调对诱导模消没（因此上述等式无关紧要，我们只需要取遍这类嵌入的乘积来构造一个诱导 $G/H$-模）。于是引理 17.3 给出所要结论。

\addcontentsline{toc}{section}{第 18 讲　范数群、Kummer 理论与投射有限上同调}
% 第 18 讲 范数群、Kummer 理论与投射有限上同调（译自 src/lecture18.txt）
\section*{第 18 讲\quad 范数群、Kummer 理论与投射有限上同调}

上一次我们证明了消没定理，并看到它蕴涵：对每个有限 Galois $G$-扩张 $L/K$，有 $(L^\times)_{tG}\simeq \mathbb Z_{tG}[-2]$；取零次上同调即得 $K^\times/NL^\times\simeq G^{\mathrm{ab}}$。注意后者不可能平凡，因为 $G$ 必须是可解群。然而在第一讲中，我们表述了另一个定理：
\[
\Gal(\bar K/K)^{\mathrm{ab}}:=\varprojlim_{L/K}\Gal(L/K)^{\mathrm{ab}}\simeq \widehat{K^\times},
\]
其中逆向极限跑遍有限 Galois 扩张 $L/K$。回忆
\[
\widehat{K^\times}:=\varprojlim_{[K^\times:\Gamma]<\infty}K^\times/\Gamma
\]
是 $K^\times$ 的投射有限完备化，其中 $\Gamma$ 是 $K^\times$ 的有限指数闭子群（这是对拓扑群定义投射有限完备化的唯一合理方式）。于是我们想证明
\[
\varprojlim_{L/K}K^\times/NL^\times\simeq \varprojlim_{[K^\times:\Gamma]<\infty}K^\times/\Gamma,
\]
其中 $L$ 与 $\Gamma$ 如上。

\textbf{定义 18.1.} $K^\times$ 的子群 $\Gamma$ 称为范数群（或范数子群），若对某个有限扩张 $L/K$ 有 $\Gamma=NL^\times$。

\textbf{定理 18.2（存在定理）.} $K^\times$ 的子群 $\Gamma$ 是范数群，当且仅当 $\Gamma$ 是闭的且指数有限。

这显然足以证明上述局部类域论的陈述。

\textbf{注 18.3.} 局部类域论的一个推论是：若 $L/K$ 是 $G$-Galois 扩张，而 $L/L_0/K$ 是 $K$ 在 $L$ 中的极大阿贝尔子扩张，则 $NL^\times=NL_0^\times$。这是因为
\[
K^\times/NL^\times\simeq G^{\mathrm{ab}}\simeq K^\times/NL_0^\times.
\]

我们将在 $\operatorname{char}(K)=0$ 的情形证明存在定理，尽管它在其他情形也成立（只是证明更复杂）。

\textbf{引理 18.4.} 若 $\Gamma\subseteq K^\times$ 是范数群，则 $\Gamma$ 闭且指数有限。

\textit{证明.} 设 $L/K$ 是 $n$ 次扩张，使 $\Gamma=NL^\times$。则 $\Gamma\supseteq N_{L/K}K^\times=(K^\times)^n$，而我们已经看到后者是有限指数的闭子群（因为它包含 $1+\mathfrak p_K^N$，其中 $N$ 充分大），故 $\Gamma$ 亦然。注意若 $\operatorname{char}(K)>0$，则 $(K^\times)^n$ 在 $K^\times$ 中的指数实际上是无限的！\hfill$\square$

于是存在定理的内容是：对一切 $n$，$\pi^{n\mathbb Z}(1+\mathfrak p_K^n)$ 都是范数子群；我们已经证明范数子群「不太小」，现在需要证明它们可以「足够小」。

\textbf{引理 18.5.} 若 $\Gamma'$ 是 $K^\times$ 的子群，且对某个范数子群 $\Gamma$ 有 $K^\times\supseteq\Gamma'\supseteq\Gamma$，则 $\Gamma'$ 也是范数子群。

\textit{证明.} 设 $L/K$ 是有限扩张，使 $\Gamma=NL^\times$。如前，不妨设 $L/K$ 阿贝尔。则由局部类域论，
\[
\Gamma'/\Gamma\subseteq K^\times/NL^\times\simeq\Gal(L/K)
\]
是正规子群，因为按假设 $\Gal(L/K)$ 阿贝尔。于是存在中间扩张 $K'/K$（$K\subseteq K'\subseteq L$），使 $\Gamma'/\Gamma=\Gal(L/K')$，并且
\[
K^\times/N(K')^\times=\Gal(K'/K)=\Gal(L/K)/\Gal(L/K')=(K^\times/NL^\times)/(\Gamma'/\Gamma)=K^\times/\Gamma'
\]
典范地成立。于是 $\Gamma'=N(K')^\times$，即所欲证。

注意我们隐含地使用了如下事实：按我们对局部类域论的明确构造，下述图表交换（对阿贝尔扩张 $L/K$）：
\[
\begin{CD}
\Gal(L/K) @>\sim>> K^\times/NL^\times\\
@V{\alpha}VV @V{}VV\\
\Gal(K'/K) @>\sim>> K^\times/N(K')^\times.
\end{CD}
\]
由于 $\Gamma'/\Gamma=\ker(\alpha)$ 在 $K^\times$ 中的原像既是 $\Gamma'$ 又是 $N(K')^\times$，我们再次得到 $\Gamma'=N(K')^\times$。\hfill$\square$

现在作一段离题讨论：在第二讲中我们说过
\[
K^\times/(K^\times)^2\simeq \Gal^{\mathrm{ab}}(\bar K/K)/2\simeq \operatorname{Hom}\big(K^\times/(K^\times)^2,\mathbb Z/2\mathbb Z\big),
\]
假设 $\operatorname{char}(K)=0$（特别地不为 $2$），其中第一个同构经由局部类域论。也就是说，$K^\times/(K^\times)^2$ 是自对偶的。现在问：如何把它推广到 $n=2$ 之外？答案是使用 Kummer 理论。

回忆：假设 $n\nmid \operatorname{char}(K)$，且 $n$ 次单位根群 $\mu_n\subseteq K^\times$ 的阶为 $n$，则
\[
K^\times/(K^\times)^n\simeq \operatorname{Hom}_{cts}\big(\Gal(\bar K/K),\mu_n\big),
\]
其中右端是群同态。要点是：若 $K$ 还是局部域，我们便期望
\begin{equation}
K^\times/(K^\times)^n\simeq \operatorname{Hom}\big(K^\times/(K^\times)^n,\mu_n\big). \tag{18.1}
\end{equation}

事实上，我们有如下映射：
\begin{align*}
K^\times/(K^\times)^n &= \operatorname{Hom}_{cts}\big(\Gal(\bar K/K),\mu_n\big)\\
&= \operatorname{Hom}_{cts}\big(\Gal^{\mathrm{ab}}(\bar K/K),\mu_n\big)\\
&= \operatorname{Hom}_{cts}\Big(\varprojlim_{L/K}K^\times/NL^\times,\mu_n\Big)\\
&= \varinjlim_{L/K}\operatorname{Hom}\big(K^\times/NL^\times,\mu_n\big)\\
&\hookrightarrow \operatorname{Hom}_{cts}\big(K^\times,\mu_n\big)\\
&= \operatorname{Hom}\big(K^\times/(K^\times)^n,\mu_n\big),
\end{align*}
其中第二个等号是因为所有这类映射都必须穿过 $\Gal(\bar K/K)$ 的阿贝尔化（因为 $\mu_n$ 阿贝尔），第三个由局部类域论，第四个由对偶性。注意这些逆向极限跑遍有限扩张 $L/K$；而「连续」（当定义域有限时并不需要）在这里指映射杀掉某个紧开子群，这就解释了上面的单射。我们想证明这个映射也是同构。注意 $K^\times/(K^\times)^n$ 是有限阿贝尔群且为 $n$-挠；因此只需证明两边阶相同。

\textbf{断言 18.6.} 设 $A$ 是 $n$-挠有限阿贝尔群。则
\[
\#A=\#\operatorname{Hom}(A,\mathbb Z/n\mathbb Z).
\]

\textit{证明.} $A$ 是若干 $\mathbb Z/d\mathbb Z$（$d\mid n$）的直和，故可归约到 $A=\mathbb Z/d\mathbb Z$ 的情形（一般情形下直和与 $\operatorname{Hom}$ 交换）。于是
\[
\operatorname{Hom}(\mathbb Z/d\mathbb Z,\mathbb Z/n\mathbb Z)=(\mathbb Z/n\mathbb Z)[d],
\]
其阶为 $d=\#A$，如所欲证。\hfill$\square$

这表明 (18.1) 是典范同构（不过仅凭断言的一般陈述也足以说明它是同构）。在 $n=2$ 的情形，容易看出这正是希尔伯特符号。

\textbf{推论 18.7.} 若 $\mu_n\subseteq K$，则 $(K^\times)^n$ 是范数子群。

\textit{证明.} 若把我们 Kummer 理论的「图景」对偶化，便得到如下交换图，连同自然的连续映射 $\alpha:K^\times\to\varprojlim_{L/K}K^\times/NL^\times$：
\[
\begin{CD}
\Gal(\bar K) @>{cts}>> \operatorname{Hom}\big(K^\times/(K^\times)^n,\mu_n\big)\\
@V{\text{投影}}VV @A{\beta_{cts}}AA\\
\varprojlim_{L/K}K^\times/NL^\times @>{\sim}>> \Gal^{\mathrm{ab}}(\bar K)
\end{CD}
\]
% [原图在 PDF 中为「Gal(K) --cts--> Hom(K^×/(K^×)^n,μn)」与「K^× --α_cts--> lim K^×/NL^× = Gal^ab(K)」两行、右侧竖直箭头 β_cts（向上）、中间竖直箭头 Gal(K)→Gal^ab(K)，已按原 PDF 第 82 页图式复原]
注意 $\alpha$ 连续，因为逆向极限中的开子群都是范数子群，从而它在 $K^\times$ 中的原像是有限指数开子群。正如我们刚刚看到的，$\ker(\beta\circ\alpha)=(K^\times)^n$，它在逆向极限中是开的（即它是某个 $L/K$ 的 $K^\times/NL^\times$ 的某个子集沿典范投影的完全原像，因为这些映射连续）。于是由引理 18.5，$(K^\times)^n$ 是范数子群。\hfill$\square$

注意上面的映射 $\beta$ 是满射，因为它实现为 $\Gal^{\mathrm{ab}}(\bar K/K)$ 模去 $n$ 次幂。

\textbf{注 18.8.}「先验地」（即如果我们忘掉每个群的阶），这个复合的核可能比 $(K^\times)^n$ 更大。通过论证二者相等，我们便得到了一个「小」的范数子群。

\textit{证明（存在定理）.} 设 $K$ 是一般的特征 $0$ 局部域。令 $L:=K(\zeta_n)$，考虑扩张 $L/K$，其中 $\zeta_n$ 表示本原 $n$ 次单位根的集合。由于由推论 18.7，$(L^\times)^n$ 是 $L^\times$ 中的范数子群，故 $N(L^\times)^n=N((L^\times)^n)\subseteq K^\times$ 是 $K^\times$ 中的范数子群。但 $N(L^\times)^n\subseteq (K^\times)^n\subseteq K^\times$，于是引理 18.5 表明 $(K^\times)^n$ 是 $K^\times$ 中的范数子群。

现在注意，对一切 $N$，存在某个 $n$ 使得
\[
(\mathcal O_K^\times)^n=(K^\times)^n\cap\mathcal O_K^\times\subseteq 1+\mathfrak p_K^N.
\]
事实上，$(\mathcal O_K^\times)^{q-1}\subseteq 1+\mathfrak p_K$，其中 $q=\#\mathcal O_K/\mathfrak p_K$（因为模 $\mathfrak p_K$ 的剩余类取 $q-1$ 次幂必为 $1$）。于是对充分大的 $v(n)$ 有 $(\mathcal O_K^\times)^{(q-1)n}\subseteq 1+\mathfrak p_K^N$，因为一般地 $(1+x)^n=1+nx+\cdots$（省略号表示高阶项），而当 $v(n)\gg 0$ 时除 $1$ 以外的所有项都落在 $\mathfrak p_K^N$ 中。

至于「$\mathbb Z$-方向」上的有限指数子群，即限制到 $\pi^N$ 的倍数上的情形，只需把 $n$ 换成 $nN$，使得只有赋值被 $N$ 整除的元素被实现。于是 $K^\times$ 的每个有限指数开子群都包含某个 $(K^\times)^n$，而后者如上所示是范数子群，故由引理 18.5 它本身也是范数子群。\hfill$\square$

现在让我们快速重访 Kummer 理论；我们将说明，它实际上关于群上同调讲了一件非常一般的事情。设 $G$ 是投射有限群，即 $G=\varprojlim_i G_i$，其中 $G_i$ 是有限群。

\textbf{定义 18.9.} 称 $G$-模 $M$ 是光滑的，若对一切 $x\in M$，存在有限指数开子群 $K\subseteq G$ 使得 $K\cdot x=x$。

\textbf{例 18.10.} 若 $G:=\Gal(\bar K/K)$，则 $G$ 同时作用在 $\bar K$ 与 $\bar K^\times$ 上，二者都是光滑 $G$-模。这是因为这两个 $G$-模中的每个元素都落在某个有限扩张 $L/K$ 里，从而被 $\Gal(\bar K/L)$ 固定，而后者按定义是有限指数开子群。

光滑性使我们能从往往非常复杂的投射有限群归约到有限群的情形。我们现在必须为投射有限群定义一个群上同调的概念，因为我们原先的表述只针对有限群。

\textbf{定义 18.11.} 设 $X$ 是下有界的光滑 $G$-模复形。则
\[
X^{hG}:=\varinjlim_i (X^{K_i})^{hG/K_i},
\]
其中 $K_i:=\ker(G\to G_i)$，而 $X^{K_i}$ 表示被 $K_i$（朴素地）固定的向量。容易看出这构成正向系统。注意 $G_i$ 不作用在 $X$ 上（它只是 $G$ 的商），但它作用在被 $K_i$ 固定的向量上。这些 $K_i$ 是 $G$ 的紧开子群，大小递减。用 $K_i$ 取「朴素不变量」是令人担心的，因为它不保持拟同构；但实际上我们有：

\textbf{断言 18.12.} 若 $X$ 无环，则 $X^{hG}$ 也无环。

证明略去；不过要注意，$X$ 下有界这一点很重要。我们有 Hilbert 定理 90 的如下「无限版本」：

\textbf{命题 18.13.} 若 $L/K$ 是（可能无限的）$G$-Galois 扩张，则
\[
H^1(G,L^\times):=H^1\big((L^\times)^{hG}\big)=0.
\]

\textit{证明.} 写 $L=\bigcup_i L_i$，其中每个 $L_i$ 是 $K$ 的有限 $G_i$-Galois 扩张。则按定义，
\[
H^1(G,L^\times)=\varinjlim_i H^1(G_i,L_i^\times)=0
\]
由 Hilbert 定理 90。\hfill$\square$
% [原文作 $H^1(G_i,K_i^\times)$，但按 $L=\bigcup_i L_i$（各 $L_i$ 为 $K$ 的有限 $G_i$-Galois 扩张）及定义 18.11（其中 $K_i=\ker(G\to G_i)$ 另有含义）应为 $L_i^\times$；已订正]

\textbf{推论 18.14.} 设 $G:=\Gal(\bar K/K)$，$n$ 与 $\operatorname{char}(K)$ 互素。若 $\mu_n\subseteq K$，则
\[
K^\times/(K^\times)^n\simeq \operatorname{Hom}_{cts}(G,\mu_n).
\]

\textit{证明.} 我们有光滑 $G$-模的短正合列
\[
0\to \mu_n\to \bar K^\times\xrightarrow{\ x\mapsto x^n\ }\bar K^\times\to 0.
\]
上同调的长正合列给出
\[
K^\times\xrightarrow{\ x\mapsto x^n\ }K^\times\to H^1(G,\mu_n)\to H^1(G,\bar K^\times),
\]
其中两端的 $K^\times$ 即 $H^0(G,\bar K^\times)$，而最右端由 Hilbert 定理 90（命题 18.13）为零。于是 $K^\times/(K^\times)^n\simeq H^1(G,\mu_n)$。由于 $\mu_n\subseteq K$（如 Kummer 理论的设定），它被 $G$ 固定；正如我们通过上链看到的，对平凡群作用有 $H^1(G,\mu_n)=\operatorname{Hom}_{cts}(G,\mu_n)$，这就给出所要结果。\hfill$\square$

于是我们确实可以从抽象群上同调与 Hilbert 定理 90 非常简单地推出 Kummer 理论。

\addcontentsline{toc}{section}{第 19 讲　布劳尔群}
% 第 19 讲 布劳尔群（源：src/lecture19.txt；已删去抽出时残留的页眉「19. BRAUER GROUPS 83/85」）
\section*{第 19 讲\quad 布劳尔群}

本讲我们概述布劳尔群。叙述将略于证明，但会给出精确的构造与命题陈述；完整证明见 [Mil13, Ser79, Boy07]。我们的动机问题是：「那些关于哈密顿代数的东西究竟是怎么回事？」（见习题集 1 第 5 题、习题集 2 第 3 题与习题集 3 第 2 题。）我们将看到，有两个对象都叫作「布劳尔群」：一个由上同调定义，另一个是更一般的代数定义；我们将证明二者一致。

回忆：若 $L/K$ 是非阿基米德局部域的 $G$-Galois 扩张，则
\[
(\Z[-2])^{tG}\simeq (L^\times)^{tG}.
\]
取 $H^0$（零次上同调）时，我们得到 lcft，即
\[
K^\times/NL^\times\simeq H^1(G,\Z)=G^{\mathrm{ab}}.
\]
这很自然，因为 $L^\times$ 位于 $0$ 次。但 $\Z$ 位于 $2$ 次，那么取 $H^2$ 会怎样？此时我们得到同构
\[
\widehat H^0(G,\Z)\simeq\widehat H^2(G,L^\times),
\]
其左端极简单地同构于
\[
\Z/[L:K]\Z=\Z/\#G\Z=\tfrac{1}{[L:K]}\Z/\Z,
\]
因为 $G$ 平凡地作用在 $\Z$ 上，故不变元正是 $\Z$，而范数对应乘以 $\#G$。然而右端要神秘得多，这促使下述定义。

\textbf{定义 19.1.} $L/K$ 的上同调布劳尔群为
\[
\Br^{\mathrm{coh}}(K/L):=H^2(G,L^\times).
\]

\textbf{注 19.2.} 若变动 $L$ 会怎样？设 $L_2/L_1/K$ 为 Galois 扩张，$\Gal(L_i/K)=G_i$（$i=1,2$）。则有短正合列
\[
0\to \Gal(L_2/L_1)\to G_2\to G_1\to 0,
\]
以及映射
\[
\Br^{\mathrm{coh}}(L_1/K)=H^2(G_1,L_1^\times)\to H^2(G_2,L_1^\times)\to H^2(G_2,L_2^\times)=\Br^{\mathrm{coh}}(L_2/K),
\]
因为关于 $G_2$ 的不变性蕴含关于 $G_1$ 的不变性，并且经由嵌入 $L_1^\times\hookrightarrow L_2^\times$ 过渡。这促使下述定义。
% 原文此处写 $\Br^{\mathrm{coh}}(L_1/K)$，而定义 19.1 写的是 $\Br^{\mathrm{coh}}(K/L)$；此处照原文排，未统一。

\textbf{定义 19.3.} $K$ 的上同调布劳尔群为
\[
\Br^{\mathrm{coh}}(K):=\varinjlim_{L/K}\Br^{\mathrm{coh}}(K/L)=H^2(\Gal(\overline K/K),\overline K^\times),
\]
其中正向极限取遍有限 Galois 扩张 $L/K$。

注意上式最右端的表达用了我们上一讲的记号。

\textbf{断言 19.4.} 在 lcft 之下，下述图表交换：
% [按原 PDF 第 86 页（印刷页 82）复原：上排 $H^2(G_2,L_1^\times)\to H^2(G_1,L_1^\times)$、下排 $\frac{1}{[L_1:K]}\Z/\Z \leftarrow \frac{1}{[L_2:K]}\Z/\Z$（左向包含），两列竖直箭头均为「$\|$」，即 lcft 给出的同构]
\[
\begin{CD}
H^2(G_2,L_1^\times) @>>> H^2(G_1,L_1^\times)\\
@V{\sim}VV @V{\sim}VV\\
\tfrac{1}{[L_1:K]}\Z/\Z @<<< \tfrac{1}{[L_2:K]}\Z/\Z
\end{CD}
\]
其中竖直的同构即 lcft，下方水平箭头为自然的包含 $\tfrac{1}{[L_2:K]}\Z/\Z\hookrightarrow\tfrac{1}{[L_1:K]}\Z/\Z$。

\textbf{推论 19.5.} 对非阿基米德局部域 $K$，有
\[
\Br^{\mathrm{coh}}(K)\simeq\Q/\Z.
\]

\textbf{注 19.6.} 还可以证明最上方的映射是单射。对非阿基米德局部域的扩张 $L/K$，有正合列
\[
0\to \Br^{\mathrm{coh}}(K/L)\to
\underbrace{\Br^{\mathrm{coh}}(K)}_{H^2(\Gal(\overline K/K),\overline K^\times)}\to
\underbrace{\Br^{\mathrm{coh}}(L)}_{H^2(\Gal(\overline K/L),\overline K^\times)},
\]
我们下次课将说明其理由。

现在转向代数视角，它给出布劳尔群的经典定义。

\textbf{命题 19.7.} 设 $K$ 为域，$A$ 是有限维 $K$-代数，其中心为 $K$。则下述条件等价：
\begin{enumerate}[label=(\arabic*)]
\item $A$ 是单的，即它没有非平凡的双边理想。
\item 对某个可分扩张 $L/K$，$A\otimes_K L\simeq M_n(L)$（即 $n\times n$ 矩阵代数）。
\item $A\simeq M_n(D)$，其中 $D$ 是 $K$ 上的中心（即以 $K$ 为中心）可除（即乘法逆元存在，但乘法未必交换）代数。
\end{enumerate}
若这些条件成立，则称 $A$ 为 $K$ 上的\textbf{中心单代数}（central simple algebra，csa；亦称 Azumaya 代数）。

\textbf{推论 19.8.} $K$ 上任何中心单代数的维数都是平方数。

\textit{证明.} $A$ 的维数在 (2) 中的张量运算下保持不变，而 $M_n(L)$ 的维数是平方数。\hfill$\square$

\textbf{例 19.9.} (1) $M_n(K)$。
(2) $K$ 上的中心可除代数。
(3) $\R$ 上的哈密顿（四元数）代数，此时 $K=\R$。
(4) 对一切满足 $\operatorname{char}(K)\neq 2$ 的域 $K$ 与元素 $a,b\in K^\times$，$H_{a,b}$ 是中心单代数。

然而，一般的域扩张 $L/K$（把 $L$ 看作 $K$-代数）并不具有中心性，其维数也未必是平方数。

\textbf{定义 19.10.} $K$ 的中心单代数布劳尔群为
\[
\Br^{\mathrm{csa}}(K):=\{\,K\ \text{上中心单代数的等价类}\,\},
\]
其中的等价关系为：$A\sim B$ 当且仅当 $A$ 与 $B$ 是上述 (3) 中同构的可除代数上的矩阵代数（维数不必相同）。

\textbf{注 19.11.} $A$ 的同构类只依赖于 $D$ 的同构类。
% 原文如此（「同构类」）；按上下文似应为「等价类」，照译并登记于报告。

\textbf{命题 19.12.} 若 $A$ 与 $B$ 是 $K$ 上的中心单代数，则 $A\otimes_K B$ 也是中心单代数（注意张量积使维数相乘，故 $A\otimes_K B$ 的维数也是平方数）。模等价类，这仅依赖于 $[A],[B]\in\Br^{\mathrm{csa}}(K)$，因此 $\Br^{\mathrm{csa}}(K)$ 构成一个阿贝尔群。

证明从略；证明 $A\otimes_K B$ 是单的相当麻烦。但我们指出，$K$ 上的矩阵代数代表 $\Br^{\mathrm{csa}}(K)$ 的单位元（或「$0$ 等价类」）。$A$ 的逆是 $A$ 但取反乘法（即 $x\cdot y:=yx$），我们记作 $A^{\mathrm{op}}$。事实上，有典范的代数同态
\[
A\otimes_K A^{\mathrm{op}}\to \operatorname{End}_K(A),
\]
其中 $A$ 与 $A^{\mathrm{op}}$ 从相反的两侧作用（注意 $\operatorname{End}_K(A)$ 属于 $\Br^{\mathrm{csa}}(K)$ 的 $0$ 等价类）。这个（非零）映射的核必是双边理想，而 $A\otimes_K A^{\mathrm{op}}$ 是单的，故该映射必是单射。又因两端在 $K$ 上的维数都是 $\dim_K(A)^2$，它因而也是满射，从而是同构。

\textbf{定义 19.13.} $L/K$ 的中心单代数布劳尔群为
\[
\Br^{\mathrm{csa}}(K/L):=\{\,\text{满足 } A\otimes_K L\simeq M_n(L)\ \text{的中心单代数 } A\ \text{的等价类}\,\},
\]
并称这样的 $A$ 在 $L$ 上\textbf{分裂}。

这一概念等价于底层的可除代数 $D$ 在 $L$ 上分裂，后者同样意味着 $D\otimes_K L\simeq M_n(L)$ 不再是可除代数。于是显然
\[
\Br^{\mathrm{csa}}(K)=\bigcup_{\substack{L/K\\ \text{可分}}}\Br^{\mathrm{csa}}(K/L).
\]

\textbf{例 19.14.} (1) 若 $K$ 代数闭，则 $\Br^{\mathrm{csa}}(K)=\{0\}$ 显然平凡。
(2) $\Br^{\mathrm{csa}}(\R)=\{\R,\mathbb H\}=\Z/2\Z$，因为 $\R$ 上唯一的另一个可除代数是 $\C$，而它不是中心的。注意哈密顿四元数代数 $\mathbb H$ 在 $\C$ 上分裂，维数为 $4=2^2$。
(3) $H_{a,b}$ 在 $K(\sqrt a)$ 与 $K(\sqrt b)$ 上分裂，它们正是 $H_{a,b}$ 中通常的交换子代数。事实上，对任意 $c\in K$ 都可取 $K(\sqrt c)$，还可以用单位作共轭。因此在下述断言中，$L$ 非常不唯一。注意与这一情形不同，这样的 $L$ 不必是 $K$ 的 Galois 扩张；例子很难找，但 $1970$ 年代已被发现。

\textbf{断言 19.15.} $n^2$ 维中心可除代数 $D/K$ 在 $n$ 次扩张 $L/K$ 上分裂，当且仅当 $K\subseteq L\subseteq D$ 作为子代数。等价地，$L$ 是 $D$ 中的极大交换子代数。

首先有一个（简单的）事实：可除代数的任何交换子代数都是域。这类似于「域上任何有限维整环本身是域」。这是因为乘以任何元素都是单射，且由矩阵给出，故是满射；于是某个元素映到单位元，给出所需的逆元。其次有一个（不显然的）事实：$\dim_K(D)$ 是 $D$ 在 $K$ 上任何极大交换子代数维数的平方。我们来证明断言的一个方向：

\textit{证明.} 设 $K\subseteq L\subseteq D$，$[L:K]=n$，$\dim_K(D)=n^2$。我们断言 $D$ 在 $L$ 上分裂。有映射
\[
D\otimes_K L\to \operatorname{End}_L(D)=M_n(L),
\]
其中 $L$ 经右乘作用在 $D$ 上，$D$ 经左乘作用在自身上；两个作用交换。该映射是单射（若对某个 $d\in D$、$l\in L$ 及所有 $x\in D$ 有 $dxl=x$，则 $dl=1$，从而对所有 $x\in D$ 有 $dx=xl^{-1}=xd$，故 $d\in K$，于是 $d\otimes l=1\otimes 1$），因此是满射，从而是同构，这正是所要。\hfill$\square$

\textbf{定理 19.16.} 对每个 Galois 扩张 $L/K$，有 $\Br^{\mathrm{coh}}(K)\simeq\Br^{\mathrm{csa}}(K)$ 以及 $\Br^{\mathrm{coh}}(K/L)\simeq\Br^{\mathrm{csa}}(K/L)$。

注意上同调布劳尔群只对 Galois 扩张 $L/K$ 定义，而中心单代数布劳尔群对所有扩张 $L/K$ 定义。这使布劳尔群的上同调定义有了意义。我们给出下述（不完整的）证明概要：

\textit{证明.} 对 $G$-Galois 扩张 $L/K$，定义映射
\begin{equation*}
H^2(G,L^\times)\longrightarrow \Br^{\mathrm{csa}}(K/L). \tag{19.1}
\end{equation*}
$H^2$ 中每个元素都由一个 $2$-上闭链代表，即满足
\[
g_1\phi(g_2,g_3)\phi(g_1,g_2g_3)=\phi(g_1g_2,g_3)\phi(g_1,g_2)
\]
的映射 $\phi:G\times G\to L^\times$（对每个 $g_1,g_2,g_3\in G$），其中我们引入了记号 ${}^g x:=g\cdot x$。对每个这样的 $\phi$，如下构造一个中心单代数：作
\[
L[G]=\Bigl\{\sum_{g\in G}x_g[g]:x_g\in L\Bigr\}
\]
并附以关系
\[
[g]x={}^g x[g],\qquad [g][h]=\phi(g,h)[gh]
\]
（对每个 $x\in L$ 与 $g,h\in G$）。下面验证 $2$-上闭链恒等式等价于结合律：
\begin{align*}
[g_1]([g_2][g_3]) &= [g_1]\phi(g_2,g_3)[g_2g_3]\\
&= {}^{g_1}\phi(g_2,g_3)[g_1][g_2g_3]\\
&= {}^{g_1}\phi(g_2,g_3)\phi(g_1,g_2g_3)[g_1g_2g_3]\\
&= \phi(g_1g_2,g_3)\phi(g_1,g_2)[g_1g_2g_3]\\
&= \phi(g_1,g_2)[g_1g_2][g_3]\\
&= ([g_1][g_2])[g_3],
\end{align*}
对一切 $g_1,g_2,g_3\in G$ 成立。我们断言（但不证明）这个中心单代数的等价类只依赖于 $\phi$ 模上边界，并且它在 $L$ 上分裂。此外，映射 (19.1) 是群同构。\hfill$\square$

这并非一个特别深刻的定理，尽管它远非显然；完整证明用到了大量并不特别好记的结构理论。

我们现在问：$K$ 上次数为 $n^2$ 的中心可除代数有多少个同构类？当 $n=1$ 时只有一个；当 $n=2$ 时同样只有一个，如习题集 3 第 2(d) 题所示。
% [原注：原文作 "When n = 1^2, there is only 1; when n = 2^2, there is again only 1"，即按次数 $n^2$（$n=1,2$，次数 1 与 4）表述；译文按 $n=1,2$ 排，语义一致（已按原 PDF 第 84 页核实，非抽出误字）]

\textbf{断言 19.17.} $K$ 上次数为 $n^2$ 的中心可除代数有 $\varphi(n)=\#(\Z/n\Z)^\times$ 个同构类。

\textbf{断言 19.18.} 设 $L/K$ 是非阿基米德局部域的 $n$ 次扩张。则下述图表交换：
\[
\begin{CD}
\Br(K) @>>> \Br(L)\\
@V{\mathrm{lcft}}VV @VV{\mathrm{lcft}}V\\
\Q/\Z @>{\times n}>> \Q/\Z
\end{CD}
\]

注意：鉴于定理 19.16，我们干脆把布劳尔群记作 $\Br$。事实上，我们已经看到两个水平映射都应有非平凡的核；当 $L=\overline K$ 时，$\Br(L)$ 平凡。于是次数为 $n^2$ 的中心可除代数在 $\Br(K)$ 中的阶为 $n$，即它是 $n$-挠的，因为它在自己的 $n$ 次极大交换子代数上分裂。等价地，$H^2(G,-)$ 是 $\#G$-挠的，如习题集 7 第 2(c) 题所证。

\textbf{推论 19.19.} 任何次数为 $n^2$ 的可除代数在 $K$ 的任何 $n$ 次扩张上分裂。

另一方面，若 $A\simeq M_n(D)$ 是 $K$ 上次数为 $n^2$ 的中心单代数，代表 $\Br(K)$ 中某个 $n$-挠类，则 $A$ 在 $D$ 的极大交换子代数上分裂。这意味着：若 $A$ 是 $n$-挠的，但对任何 $m\mid n$ 都不是 $m$-挠的，则 $A$ 本身是可除代数。这给出 $\varphi(n)$：可除代数来自 $\Z/n\Z$ 中那些不来自更小的 $\Z/m\Z$ 的类。另一个推论如下：

\textbf{推论 19.20.} $K$ 上任何次数为 $n^2$ 的中心可除代数都包含 $K$ 的每个 $n$ 次域扩张。
% 原文抽出作 "every degree-n field extension of L"，按上下文应为「$K$ 的」；照上下文译，并登记于报告。

\textit{证明.} 由断言 19.18，$\Br(K)$ 中任何 $n$-挠类在 $K$ 的任何 $n$ 次扩张 $L$ 上的布劳尔群 $\Br(L)$ 中像为零。故它在 $L$ 上分裂，从而由断言 19.15 包含 $L$。\hfill$\square$

对 $n=2$，这就是「每个 $x\in K$ 在 $H_{a,b}$ 中有平方根」这一结果，它已在习题集 3 第 2(c) 题中证明。

\addcontentsline{toc}{section}{第 20 讲　第一不等式的证明}
% 第 20 讲 第一不等式的证明（源：src/lecture20.txt；已删去抽出时残留的页眉「20. PROOF OF THE FIRST INEQUALITY 87/89/91」）
\section*{第 20 讲\quad 第一不等式的证明}

我们首先兑现上一讲的承诺。设 $K^{\mathrm{sep}}=\overline K$，并设 $L/K$ 是非阿基米德局部域的扩张。
% 原文作 "Let $K^{\mathrm{sep}}=\overline K/L/K$ be an extension of nonarchimedean local fields"，即「$K^{\mathrm{sep}}=\overline K$，$L/K$ 为扩张」，抽出时斜杠粘连，译文分句写出

\textbf{断言 20.1.} 有正合列
\[
0\to \Br^{\mathrm{coh}}(K/L)\to \Br^{\mathrm{coh}}(K)\to \Br^{\mathrm{coh}}(L).
\]

回忆 $\Br^{\mathrm{coh}}(K/L)$ 本质上编码了 $K$ 上在 $L$ 上分裂的可除代数，而 $\Br^{\mathrm{coh}}(K)$ 与 $\Br^{\mathrm{coh}}(L)$ 分别编码 $K$ 上与 $L$ 上的可除代数。映射 $\Br^{\mathrm{coh}}(K)\to\Br^{\mathrm{coh}}(L)$ 的核本质上是 $K$ 上在 $L$ 上分裂的可除代数，但我们将在上同调一侧把这一点说清楚。可以用谱序列来做（这有点烦人），但我们着眼于链复形，这会使证明概念上清晰得多。我们需要下述构造：

\textbf{定义 20.2.} 设 $X$ 是 $A$-模的链复形。则 $\tau^{\leq n}X$ 是「截断」链复形
\[
\cdots\to X_{n-2}\xrightarrow{d_{n-2}}X_{n-1}\xrightarrow{d_{n-1}}\operatorname{Ker}(d_n)\to 0\to 0\to\cdots
\]

\textbf{引理 20.3.} 恒等映射 $\tau^{\leq n}X\to X$ 是复形映射，且对一切 $i\leq n$ 给出同构 $H^i(\tau^{\leq n}X)\xrightarrow{\sim}H^i(X)$。

\textit{证明.} 只需注意，由链复形的定义，下图中两个中央方块交换：
\[
\begin{array}{ccccccccccc}
\cdots & \longrightarrow & X_{n-2} & \xrightarrow{d_{n-2}} & X_{n-1} & \xrightarrow{d_{n-1}} & \operatorname{Ker}(d_n) & \longrightarrow & 0 & \longrightarrow & 0\ \cdots\\[8pt]
 & & \big\downarrow & & \big\downarrow & & \big\downarrow & & & & \\[8pt]
\cdots & \longrightarrow & X_{n-2} & \xrightarrow{d_{n-2}} & X_{n-1} & \xrightarrow{d_{n-1}} & X_n & \xrightarrow{d_n} & X_{n+1} & \longrightarrow & \cdots
\end{array}
\]
\hfill$\square$

\textbf{推论 20.4.} 若 $X\xrightarrow{\mathrm{qis}}Y$，则诱导映射 $\tau^{\leq n}X\to\tau^{\leq n}Y$ 也是拟同构。

\textbf{注 20.5.} 截断 $\tau^{\leq n}$ 是唯一不与锥（cone）交换的运算。

\textbf{定义 20.6.} $\tau^{\geq n+1}X:=\operatorname{hCoker}(\tau^{\leq n}X\to X)$（同伦余核）。

\textbf{引理 20.7.} 对一切 $i\geq n+1$，有 $H^iX\xrightarrow{\sim}H^i\tau^{\geq n+1}X$。

\textit{证明.} 简单应用上同调长正合列即可。\hfill$\square$

为应用于布劳尔群，我们先引入一些记号：
\[
G_K:=\Gal(\overline K/K),\qquad G_L:=\Gal(\overline K/L),\qquad H:=\Gal(L/K),
\]
其中注意：为使 $\Br^{\mathrm{coh}}(K/L)$ 有定义，$L/K$ 必须是 Galois 的。回忆
\[
\Br^{\mathrm{coh}}(K):=H^2(G_K,\overline K^\times),\quad
\Br^{\mathrm{coh}}(L):=H^2(G_L,\overline K^\times),\quad
\Br^{\mathrm{coh}}(K/L):=H^2(H,L^\times).
\]
% 原文抽出时 $\overline K$ 的上划线丢失（作 "K ×"），已按 $G_K=\Gal(\overline K/K)$ 的记号补回。
主要观察是
\[
\bigl[(\overline K^\times)^{hG_L}\bigr]^{hH}=(\overline K^\times)^{hG_K}
\]
（由习题集 6 第 3 题），因为我们有短正合列
\[
1\to G_L\to G_K\to H\to 1.
\]

\textit{证明}（断言 20.1）. 注意 $L^\times=\tau^{\leq 0}(\overline K^\times)^{hG_L}$，作为 $H$-模的复形。于是按定义
\[
\operatorname{hCoker}\bigl(L^\times\to(\overline K^\times)^{hG_L}\bigr)
=\tau^{\geq 1}(\overline K^\times)^{hG_L}\simeq\tau^{\geq 2}(\overline K^\times)^{hG_L},
\]
因为这等价于断言 $H^1(G_L,\overline K^\times)=0$，而后者正是希尔伯特定理 90。因此
\[
\operatorname{hCoker}\bigl((L^\times)^{hH}\to(\overline K^\times)^{hG_K}\bigr)
=\bigl(\tau^{\geq 2}(\overline K^\times)^{hG_L}\bigr)^{hH}.
\]
最后，上同调长正合列给出
\[
H^1\bigl(\tau^{\geq 2}(\overline K^\times)^{hG_L}\bigr)^{hH}\to
\underbrace{H^2(H,L^\times)}_{\Br^{\mathrm{coh}}(K/L)}\to
\underbrace{H^2(G_K,\overline K^\times)}_{\Br^{\mathrm{coh}}(K)}\to
\underbrace{H^2(G_L,\overline K^\times)^{H}}_{\Br^{\mathrm{coh}}(L)^{H}}\hookrightarrow\Br^{\mathrm{coh}}(L),
\]
因为对第一项而言，取群上同调只会提高复形的次数，因此不可能产生非平凡的 $1$ 次上同调；而对最后一项而言，群上同调等价于在复形最低非零次处取不变元。\hfill$\square$

现在转向整体类域论，它在许多方面不如局部类域论优美；我们的处理也将相应地不那么详尽。

\textbf{定理 20.8（gcft 的主定理）.} 设 $F$ 为整体域。则
\[
\Gal(\overline F/F)^{\mathrm{ab}}\simeq\widehat{A_F^\times/F^\times},
\]
并且对一切有限 Galois 扩张 $E/F$，有
\[
\Gal(E/F)^{\mathrm{ab}}\simeq N(A_E^\times)\backslash A_F^\times/F^\times.
\]
% [原文疑误：原印作 $\Gal^{\mathrm{ab}}(G)\simeq(\widehat{\mathbb A_F^\times/F^\times})$，其中 $G$ 未在任何地方定义，应为 $\Gal(\overline F/F)$；已按标准主定理订正（原 PDF 第 87 页）]

注意这里 $\widehat{\cdot}$ 照常表示投射有限完备化；在第二式中我们对 $A_F^\times$ 模掉两个不同的对象。

gcft 的「箴言」是：$C_F:=A_F^\times/F^\times$（$F$ 的伊代尔类群）扮演 lcft 中 $K^\times$ 的角色，二者表现出非常相似的行为（有「类构造」理论使这一类比更紧密）。因此，在 gcft 的证明中我们将搬用 lcft 的许多方法；例如，我们预期某种「存在定理」连同 Kummer 理论能使我们从后者推出定理 20.8 的前一结论。课上大概没有时间展示这一蕴含，所以我们转而集中证明第二个结论。

我们的主要「输入」如下：

\textbf{定理 20.9（gcft 的不等式）.} (1) 设 $E/F$ 是整体域的 $n$ 次循环扩张。则 $\chi(C_E)=n$（即 $C_E$ 的 Herbrand 商）。
(2) 对整体域的一切有限 $G$-Galois 扩张 $E/F$，$H^1(G,C_E)=0$。

第二个结论与希尔伯特定理 90 类似，但证明困难得多。这两个结论在 gcft 中分别被称为「第一不等式」与「第二不等式」（各种叫法不尽相同）；我们（任意地）偏好的约定已在上面标明。确实，第一不等式给出
\[
\#\widehat H^0(\Z/n\Z,C_E)=n\cdot\#H^1(\Z/n\Z,C_E)\geq n,
\]
而第二不等式给出 $\#H^0(\Z/n\Z,C_E)\leq n$，故它恰为 $n$。

本讲我们证明第一不等式。以下恒设 $E/F$ 是整体域的 $n$ 次循环扩张。为此，我们将计算 $\chi(E^\times)$ 与 $\chi(A_E^\times)$，或至少它们的商；经由正合列
\[
0\to E^\times\to A_E^\times\to C_E\to 0
\]
这将给出 $\chi(C_E)$。一个小问题是：两个 Herbrand 商都是无穷的。

首先，关于 $A_E^\times$ 的群上同调的结构，有一个一般性评注：

\textbf{断言 20.10.} 对整体域的 $G$-Galois 扩张 $E/F$，对每个 $i$ 有
\[
\widehat H^i(G,A_E^\times)=\bigoplus_{v\in M_F}\widehat H^i\bigl(G,(E\otimes_F F_v)^\times\bigr),
\]
其中 $M_F$ 表示 $F$ 的素位集（即赋值等价类之集）。

需要说明，这个断言完全是局部的：它利用阿代尔的结构把关于整体域的所有局部信息汇聚起来；我们并未使用 $E^\times$ 的对角嵌入（正是它使一个域能整体地出现在阿代尔之中）。注意 $E\otimes_F F_v$ 类似于 $E$ 在 $v$ 处的完备化。

\textit{证明.} 回忆
\begin{align*}
A_E^\times&=\varinjlim_{\substack{S\subset M_F\\ \#S<\infty}}\biggl(\prod_{v\in S}(E\otimes_F F_v)^\times\\
&\qquad\qquad\times\prod_{v\notin S}\bigl(\underbrace{\mathcal O_E\otimes_{\mathcal O_F}\mathcal O_{F_v}}_{\mathcal O_E\otimes_F F_v}\bigr)^\times\biggr),
\end{align*}
其中 $S$ 包含所有分歧素位与无限素位（后者之集我们今后记作 $M_F^\infty$）。注意通常我们取 $S\subset M_E$；这一替代写法改用 $E$ 的素位落在 $F$ 的哪些素位之上来表达它们。由习题集 6 第 6(h) 题，$\widehat H^i(G,-)$ 与一致有下界的正向极限交换，也与直积交换。这意味着
\begin{align*}
\widehat H^i(G,A_E^\times)
&=\varinjlim_{\substack{S\subset M_F\\ \#S<\infty}}\biggl(\bigoplus_{v\in S}\widehat H^i\bigl(G,(E\otimes_F F_v)^\times\bigr)\\
&\qquad\qquad\times\prod_{v\notin S}\widehat H^i\bigl(G,(\mathcal O_E\otimes_{\mathcal O_F}\mathcal O_{F_v})^\times\bigr)\biggr)\\
&=\varinjlim_{\substack{S\subset M_F\\ \#S<\infty}}\bigoplus_{v\in S}\widehat H^i\bigl(G,(E\otimes_F F_v)^\times\bigr).
\end{align*}
确实，由于假定 $S$ 包含所有分歧素位，对任意 $v\notin S$ 与 $w\mid v$，局部域的扩张 $E_w/F_v$ 非分歧，故其 Galois 群就是分解群 $G_w$。于是由习题集 7 第 1(d) 题，对任意 $v\notin S$ 与 $w\mid v$ 有
\[
\widehat H^i\bigl(G,(\mathcal O_E\otimes_{\mathcal O_F}\mathcal O_{F_v})^\times\bigr)
=\widehat H^i(G_w,\mathcal O_{E_w}^\times)=0,
\]
因为
\[
(\mathcal O_E\otimes_{\mathcal O_F}\mathcal O_{F_v})^\times
=\prod_{w\mid v}\mathcal O_{E_w}^\times
=\Z[G]\otimes_{\Z[G_w]}\mathcal O_{E,w}^\times
\]
是诱导 $G$-模，而后一表达按例 14.3 所述消没。\hfill$\square$
% [原文疑误：原印作 $(\mathcal O_E\otimes_{\mathcal O_F}\mathcal O_k)^\times$，下标 $k$ 应为 $v$（与上一行 $(\mathcal O_E\otimes_{\mathcal O_F}\mathcal O_{F_v})^\times$ 及下一行的 $\prod_{w\mid v}$ 分解一致）；已径改]

现在我们设法绕过 Herbrand 商的无穷性。

\textbf{定义 20.11.} 对任意有限 $M_F^\infty\subseteq S\subset M_F$，令
\[
A_{F,S}:=\prod_{v\in S}F_v\times\prod_{v\notin S}\mathcal O_{F_v}
\]
为 $S$-阿代尔环，并类似地记 $S$-伊代尔群为 $A_{F,S}^\times$。

\textbf{引理 20.12.} 存在有限 $S\subset M_F$ 使 $A_{F,S}^\times\cdot F^\times=A_F^\times$。

\textit{证明.} 该等式等价于说映射
\[
A_{F,S}^\times\to A_{F,M_F^\infty}^\times\backslash A_F^\times/F^\times=\mathrm{Cl}(F)
\]
满射到 $F$ 的类群，其中我们也可以模掉 $A_{F,M_F^\infty}^\times$（因为按假设它包含于 $A_{F,S}^\times$），而最后的典范同构来自习题集 2 第 1(b) 题。在该同构下，取 $\mathfrak p\subseteq\mathcal O_K$ 为素理想，$F_{\mathfrak p}$ 的素元映到 $[\mathfrak p]\in\mathrm{Cl}(F)$。由于 $\mathrm{Cl}(F)$ 有限，我们只需取 $S$ 为 $M_F^\infty$ 连同这样一组素位：其中每个素位各对应 $\mathrm{Cl}(F)$ 的一个不同元素。\hfill$\square$
% 原文此处作 "a uniformizer p ⊆ O_K of F_p"，素元记号与素理想记号混同、O_K 疑为 O_F；按上下文译，并登记于报告。

现在回到 $E/F$ 循环的情形，取有限集 $S\subset M_F$，使其包含 $F$ 的无限素位与分歧素位，并满足 $A_{E,S}^\times\cdot E^\times=A_E^\times$；由引理这是可能的（注意 $A_{E,S}^\times:=A_{E,T}^\times$，其中 $T$ 是 $E$ 中位于 $S$ 中 $F$ 的素位之上的素位集；因此我们实际上是对 $E$ 应用引理，再把得到的素位集投影到 $F$）。
% 原文作「containing the infinite and ramified places of S」，按上下文应为「of F」；照上下文译，并登记于报告。

\textbf{断言 20.13.} 有短正合列
\[
0\to E_S^\times\to A_{E,S}^\times\to C_E\to 0,
\]
其中 $E_S^\times:=E^\times\cap A_{E,S}^\times$。

\textit{证明.} 有下述交换图：
% [按原 PDF 第 90 页（印刷页 86）复原：原图上方另有一整行 $0\to0\to0$，各行左端均起于 $0$、右端（前两行）止于 $0$，另有两条虚线斜箭头示意蛇引理；$\varphi$ 标于下排向左的箭头，$\psi$ 标于右列第二处竖直箭头]
\[
\begin{array}{ccccccccc}
 & & 0 & \longrightarrow & 0 & \longrightarrow & 0 & & \\[6pt]
 & & \big\downarrow & & \big\downarrow & & \big\downarrow & & \\[6pt]
0 & \longrightarrow & E_S^\times & \longrightarrow & A_{E,S}^\times & \longrightarrow & A_{E,S}^\times/E_S^\times & \longrightarrow & 0\\[6pt]
 & & \big\downarrow & & \big\downarrow & & \big\downarrow\psi & & \\[6pt]
0 & \longrightarrow & E_S^\times & \longrightarrow & A_E^\times & \longrightarrow & C_E & \longrightarrow & 0\\[6pt]
 & & \big\downarrow & & \big\downarrow & & \big\downarrow & & \\[6pt]
 & & E^\times/E_S^\times & \xleftarrow{\ \varphi\ } & A_E^\times/A_{E,S}^\times & \longrightarrow & 0 & &
\end{array}
\]
其中映射 $\varphi$ 由 $S$ 的选取是满射，并且由 $E_S^\times$ 的定义也是单射（即若 $\varphi(x)\in A_{E,S}^\times$ 则 $x\in E_S^\times$）。蛇引理于是推出映射 $\psi$ 是同构，这正是所要。\hfill$\square$

由断言 20.10，
\[
\widehat H^i(G,A_{E,S}^\times)=\bigoplus_{v\in S}\widehat H^i\bigl(G,(E\otimes_F F_v)^\times\bigr),
\]
% [原文疑误：原印作 $\widehat H^i(E,A_{E,S}^\times)=\bigoplus_{v\in S}\widehat H^\times(G,(E\otimes_F F_v)^\times)$，两处第一变元/上标均有误——左侧第一变元应为群 $G$（原 PDF 第 90 页），右侧上标应为 $i$；已径改]
于是
\[
\chi(A_{E,S}^\times)=\prod_{v\in S}\chi\bigl((E\otimes_F F_v)^\times\bigr)
=\prod_{v\in S}\chi\Bigl(\prod_{w\mid v}E_w^\times\Bigr)
=\prod_{v\in S}\chi(E_w^\times)
=\prod_{v\in S}[E_w:F_v]
\]
（对某个 $w\mid v$ 的选取），由断言 7.8 以及如前
\[
\Bigl(\prod_{w\mid v}E_w^\times\Bigr)^{tG}\simeq(E_w^\times)^{tG_w}.
\]

现在我们要计算 $\chi(E_S^\times)$，并希望它满足
\[
n\cdot\chi(E_S^\times)=\prod_{v\in S}[E_w:F_v],
\]
同样对某个 $w\mid v$。由狄利克雷（Dirichlet）单位定理，$E_S^\times$ 有限生成，因此
\[
E_S^\times\simeq(E_S^\times)_{\mathrm{tors}}\times\Z^r,
\]
其中 $r$ 是 $E_S^\times$ 的秩。于是模去挠部分，$E_S^\times$ 是一个 $r$ 维格。

\textbf{引理 20.14.} 若循环群 $G$ 作用在 $\R$-向量空间 $V$ 上，且 $\Lambda_1,\Lambda_2\subseteq V$ 是两个在 $G$-作用下不变的格，则 $\chi(\Lambda_1)=\chi(\Lambda_2)$（把两个格都看作 $G$-模）。

证明我们留到下一讲，以免时间不够。现在回忆单位定理证明中所用的映射：
\[
E_S^\times\to\prod_{w\in T}\R,\qquad x\mapsto(\log|x|_w)_w,
\]
其中 $T$ 是 $E$ 中位于 $S$ 中素位之上的素位，如前。单位定理的证明表明该映射有有限核，其像则是超平面 $\{(y_w)_w:\sum_w y_w=0\}$ 中的一个格。因此
\[
\operatorname{Im}(E_S^\times)\cup\Z\cdot(1,1,\dots,1)
\]
是 $\prod_{w\in T}\R$ 中的一个格，于是
\[
\chi(E_S^\times)=\chi\bigl(\operatorname{Im}(E_S^\times)\bigr)
=\frac{\chi(\Lambda)}{\chi(\Z)}=\frac{\chi(\Lambda)}{n},
\]
其中 $\Lambda$ 是任一 $G$-不变的格。现在，
\[
\Lambda:=\prod_{w\in T}\Z\subseteq\prod_{w\in T}\R
\]
就是这样一个格。Galois 群 $G$ 作用在 $\prod_{w\mid v}\Z$ 上，因此
\[
\chi(\Lambda)=\chi\Bigl(\prod_{v\in S}\prod_{w\mid v}\Z\Bigr)
=\prod_{v\in S}\chi\Bigl(\prod_{w\mid v}\Z\Bigr)
=\prod_{v\in S}\#G_w=\prod_{v\in S}[E_w:F_v]
\]
（对 $w$ 作一选取），正如所要。这就完成了第一不等式的证明。

\addcontentsline{toc}{section}{第 21 讲　阿廷与布劳尔互反律（上）}
% 第 21 讲 阿廷与布劳尔互反律（上）（源：src/lecture21.txt；已删去抽出时残留的页眉「21. ARTIN AND BRAUER RECIPROCITY, PART I 93/95」）
\section*{第 21 讲\quad 阿廷与布劳尔互反律（上）}

设 $F/\Q$ 是整体域，从而有正合列
\[
1\to F^\times\hookrightarrow A_F^\times\to C_F\to 1,
\]
其中 $C_F:=A_F^\times/F^\times$。上一讲我们几乎证明了：对循环 $n$ 次扩张 $E/F$，有 $\chi(C_E)=n$；剩下的只是证明引理 20.14。

\textbf{引理 21.1.} 设 $L/K$ 是无限域的扩张，$A$ 是 $K$-代数，$M$ 与 $N$ 是 $A$-模，在 $K$ 上有限维，且 $M$ 在 $A$ 上投射。若 $M\otimes_K L\cong N\otimes_K L$ 作为 $A\otimes_K L$-模同构，则 $M\cong N$。

\textbf{证明.} 首先注意有同构
\begin{equation}\tag{21.1}
\operatorname{Hom}_A(M,N)\otimes_K L\xrightarrow{\sim}\operatorname{Hom}_{A\otimes_K L}(M\otimes_K L,\,N\otimes_K L).
\end{equation}
事实上，由于 $M$ 是有限生成投射模，它是某个有限秩自由模的直和项，这就把问题化到 $M=A$ 的情形（因为 Hom 函子与直和交换）。此时 (21.1) 两端都化为 $N\otimes_K L$。（事实上，在更大的一般性下有更基本的恒等式：对任意代数 $A$，只要 $P$ 平坦、$M$ 有限表现，即存在 $n_1,n_2\in\Z$ 使 $A^{n_1}\to A^{n_2}\to M\to 0$ 正合，就有 $\operatorname{Hom}_A(M,N\otimes_A P)=\operatorname{Hom}_A(M,N)\otimes_A P$。）

注意 $M\otimes_K L$ 与 $N\otimes_K L$ 在 $L$ 上有相同的维数，故 $M$ 与 $N$ 在 $K$ 上有相同的维数 $d$。令 $V:=\operatorname{Hom}_A(M,N)$，$W:=\operatorname{Hom}_A(\Lambda^dM,\Lambda^dN)$，其中 $\Lambda^d(-)$ 表示第 $d$ 次外幂。二者都是有限维 $K$-向量空间；特别地，$V$ 是 $d^2$ 维的，$W$ 是 $1$ 维的，即同构于 $K$。$\Lambda^d$ 的函子性给出行列式映射 $\det:V\to W$，它是系数在 $K$ 中的 $d$ 次多项式映射，意思是：取定 $V$ 与 $W$ 的基之后，它由坐标的 $d$ 次多项式给出（因为它是按 $V$ 中 $d\times d$ 矩阵的行列式来计算的）。我们希望证明 $\det$ 非零，因为满足 $\det\varphi\neq0$ 的点 $\varphi\in V$ 恰好就是一个 $A$-模同构 $M\cong N$。由于 $K$ 无限（而 $\det$ 是多项式），只需在将系数扩张到 $L$ 之后检验；由 (21.1)，这给出行列式映射
\[
\operatorname{Hom}_{A\otimes_KL}(M\otimes_KL,\,N\otimes_KL)\to\operatorname{Hom}_{A\otimes_KL}\bigl(\Lambda^d(M\otimes_KL),\,\Lambda^d(N\otimes_KL)\bigr)\cong L.
\]
由于左端含有一个同构 $M\otimes_KL\cong N\otimes_KL$，此映射必非零，证毕。\hfill$\square$

\textbf{证明（引理 20.14）.} 这里的麻烦在于 $\Lambda_1$ 与 $\Lambda_2$ 可能不可公度（即每个都容在另一个的有限指数与倍数之内）。然而我们断言 $\Lambda_1\otimes\Q\cong\Lambda_2\otimes\Q$ 作为 $\Q[G]$-模。事实上，由引理 21.1，只需证明 $\Lambda_1\otimes R\cong\Lambda_2\otimes R$ 作为 $R[G]$-模，而这显然，因为 $\Lambda_i\otimes R\cong V$（$i=1,2$）。于是取 $\Lambda_2$ 在此同构下的像，并与 $R$ 张量以得到 $V$ 中的包含，就存在 $G$-稳定格 $\Lambda_3\subseteq V$，它作为 $\Z[G]$-模同构于 $\Lambda_2$，且与 $\Lambda_1$ 可公度。于是对充分大的 $N$ 有 $N\Lambda_3\subseteq\Lambda_1\subseteq\frac1N\Lambda_3$，故 $\chi(\Lambda_1)=\chi(\Lambda_3)=\chi(\Lambda_2)$，因为所有子商都有限，证毕。\hfill$\square$

现在转向局部 Kronecker--Weber 理论的讨论。

\textbf{定理 21.2.} 对任意素数 $p$，
\[
\Q_p^{\mathrm{ab}}=\Bigl[\bigcup_{n\ge0}\Q_p(\zeta_{p^n})\Bigr]\cdot\Q_p^{\mathrm{unr}},
\]
其中这是与 $\Q_p$ 的未完备化极大非分歧扩张的复合。

\textbf{证明.} 回忆扩张 $\Q_p(\zeta_{p^n})/\Q_p$ 完全分歧，其 Galois 群为 $(\Z/p^n\Z)^\times$。于是 $K\subseteq\Q_p^{\mathrm{ab}}$，其中把右端记为 $K$，且这给出映射
\[
\Z_p^\times\times\widehat\Z=\widehat{\Q}_p^{\times}=\Gal(\Q_p^{\mathrm{ab}}/\Q_p)\longrightarrow\Gal(K/\Q_p)=\Bigl[\varprojlim_n(\Z/p^n\Z)^\times\Bigr]\times\widehat\Z,
\]
其中最左的等号由 lcft 得到，随后是选取 $\Q_p$ 的一个素元；由 Galois 理论该映射满射。由于两端作为抽象群同构，下述引理表明此映射是同构。\hfill$\square$

\textbf{引理 21.3.} 设 $G$ 是投射有限群，使得对一切 $n>0$，$G$ 中指数为 $n$ 的开子群只有有限多个（即 $G$「拓扑有限生成」）。则每个连续同态 $\varphi:G\to G$ 都是同构。

\textbf{证明.} 若 $H\subseteq G$ 是指数至多 $n$ 的子群，则 $\varphi^{-1}(H)\subseteq G$ 也是指数至多 $n$ 的子群。于是有映射
\[
\{H\subseteq G:[G:H]\le n\}\xrightarrow{\ \varphi^{-1}\ }\{H\subseteq G:[G:H]\le n\},
\]
% [原注：原文称此映射「因 $\varphi$ 满射而单射」，但引理假设中 $\varphi$ 只是一般连续同态（如 $\varphi\equiv1$ 则该映射显然不单射）；原文此处有瑕——正确做法是改对 $\varphi(G)$ 施用同一计数论证。译文照译，见报告 §3]
它单射，因为 $\varphi$ 满射。由假设此集合有限，故该映射双射。由于
\[
\operatorname{Im}(\varphi^{-1})=\{H\subseteq G:[G:H]\le n,\ \operatorname{Ker}(\varphi)\subseteq H\},
\]
可知 $\operatorname{Ker}(\varphi)$ 含于 $G$ 的每个有限指数子群，故 $\operatorname{Ker}(\varphi)$ 平凡，$\varphi$ 是同构。\hfill$\square$

现在问：定理 21.2 证明中 $\Z_p^\times\times\widehat\Z$ 的自同构是什么？下述「显式 cft」定理回答此问题，但其证明相当复杂且绝非显然（见 [Dwo58]）。对整体域的类似问题一般未知，只有 $\Q$ 与虚数域的情形除外。

\textbf{定理 21.4（Dwork, Lubin--Tate）.} (1) 元素 $p\in\Q_p^\times$ 在 $\bigcup_n\Q_p(\zeta_{p^n})$ 上平凡作用，而在 $\Q_p^{\mathrm{unr}}$ 上作为 Frobenius 元素作用。(2) 元素 $x\in\Z_p^\times$ 在 $\Q_p^{\mathrm{unr}}$ 上平凡作用，而在 $\bigcup_n\Q_p(\zeta_{p^n})$ 上以 $x^{-1}$ 作用，即 $\theta_p(x)\cdot\zeta_{p^n}=\zeta_{p^n}^{(x^{-1}\bmod p^n)}$，其中
\[
\theta_p:\Q_p^\times\to\Gal\Bigl(\bigcup_n\Q_p(\zeta_{p^n})\cdot\Q_p^{\mathrm{unr}}\Big/\Q_p\Bigr)
\]
是 lcft 给出的同态。

现在我们想介绍两条互反律：阿廷互反律与布劳尔互反律。先从前者开始。设 $E/F$ 是整体域的阿贝尔 $G$-Galois 扩张。回忆我们期待有
\[
F^\times\backslash A_F^\times/N(A_E^\times)=C_F/N(C_E)\xrightarrow{\ \sim\ }G.
\]
我们想构造此映射。

\textbf{断言 21.5.} lcft 给出映射 $\theta:A_F^\times\to G$，满足 $\theta(N(A_E^\times))=1$。

\textbf{证明.} 设 $x\in A_F^\times$。对每个 $v\in M_F$ 有元素 $x_v\in F_v^\times$，lcft 于是给出映射
\[
\theta_v:F_v^\times\to\Gal(E_w/F_v)\subseteq\Gal(E/F)=G
\]
其中 $w\mid v$ 是 $E$ 的一个素位，前者是 $E/F$ 在 $v$ 处的分解群。此嵌入由嵌入 $E\subseteq E_w$ 诱导。回忆当 $E/F$ 阿贝尔时，$\Gal(E_w/F_v)$ 与 $w$ 的选取无关。

我们断言乘积
\[
\theta(x):=\prod_{v\in M_F}\theta_v(x_v)
\]
有意义，即除有限多个 $v$ 外 $\theta_v(x_v)=1$。事实上，对几乎所有的 $v$，$E_w/F_v$ 非分歧且 $x_v\in\mathcal{O}_{F_v}^\times$，从而 $\theta_v(x_v)=1$（因为由 lcft，映射 $\theta_v$ 杀掉 $\mathcal{O}_{F_v}^\times$，并把 $F_v$ 的素元送到 $G$ 的 Frobenius 元素）。由于对一切 $v\in M_F$ 有 $\theta_v(N(E_w^\times))=1$，故 $\theta(N(A_E^\times))=1$，证毕。\hfill$\square$

\textbf{定理 21.6（阿廷互反律）.} 有 $\theta(F^\times)=1$，从而 $\theta$ 给出映射 $C_F/N(C_E)\to G$。

\textbf{例 21.7.} 若 $E/\Q$ 是二次扩张，这归结为二次互反律。事实上，局部阿廷映射恰由希尔伯特符号给出，而我们此前已证明了这一蕴含。

我们将与布劳尔互反律一并证明这一点。
% [原文作 "We will proof this"，疑为 "prove" 之误，见报告 §3]
设 $E/F$ 是整体域的有限 $G$-扩张。有
\[
H^2(G,A_E^\times)=\bigoplus_{v\in M_F}\Br(F_v)=\bigoplus_{v\in M_F\setminus M_F^\infty}\Q/\Z\times\bigoplus_{v\in M_F^\infty}\tfrac12\Z/\Z,
\]
由断言 20.10 以及 $\Br(\R)=\Z/2\Z$。通过
\[
\bigoplus_{v\in M_F/M_F^\infty}\Q/\Z\xrightarrow{\ (x_v)_v\mapsto\sum_vx_v\ }\Q/\Z
\]
定义不变映射 $\iota:H^2(G,A_E^\times)\to\Q/\Z$，即对所有局部因子求和（并忽略一切无限因子）。

\textbf{定理 21.8（布劳尔互反律）.} 复合
\[
\Br(F/E)\to H^2(G,A_E^\times)\xrightarrow{\ \iota\ }\Q/\Z
\]
对一切 $E/\Q$ 为零。
% [原文作 "is zero for all E/Q"；定理设定为 E/F 有限扩张，此处照译，见报告 §3]

注意 $E$ 在这里本质无关紧要；此定理其实关于 $\Br(F)$。第一个映射由对角嵌入 $E\hookrightarrow A_E^\times$ 诱导，而布劳尔群的上同调解释随后表明：$F$ 上每个可除代数除有限多个素位外都是矩阵代数，这本身并不是显然的陈述。

\textbf{例 21.9.} 任意 $a,b\in\Q^\times$ 给出 $\Q$ 上的哈密顿代数 $H_{a,b}$。在有限素位 $p$ 处，不变量为
\[
\begin{cases}
0 & \text{若 } H_{a,b}\otimes\Q_p\cong M_2(\Q_p),\\
1/2 & \text{否则},
\end{cases}
\qquad\Longleftrightarrow\qquad
\begin{cases}
0 & \text{若 } (a,b)_p=1,\\
1/2 & \text{若 } (a,b)_p=-1.
\end{cases}
\]
于是，「所有不变量之和为零」这一断言又是二次互反律。

\textbf{断言 21.10.} 阿廷互反律对 $\Q(\zeta_n)/\Q$ 成立，其中 $\zeta_n$ 是本原 $n$ 次单位根。

\textbf{证明.} 我们借助 Dwork 定理作显式计算。可设 $n=\ell^r$ 是素数幂，因为 $\Q(\zeta_n)$ 是各素数幂因子 $\ell^r\mid n$ 对应的 $\Q(\zeta_{\ell^r})$ 的复合，而 $\Gal(\Q(\zeta_n)/\Q)$ 随后分解为 $\Gal(\Q(\zeta_{\ell^r})/\Q)$ 的积。于是有复合
\[
\Q^\times\hookrightarrow A_\Q^\times\xrightarrow{\ \theta\ }\Gal(\Q(\zeta_{\ell^r})/\Q)=(\Z/\ell^r\Z)^\times,
\]
其中 $\theta=\prod_p\theta_p$ 如前。我们想证明此映射平凡。为此只需证明对一切素数 $p$ 有 $\theta(p)=1$ 且 $\theta(-1)=1$。设 $p\ne\ell$；$p=\ell$ 的情形将在下一讲处理。则
\[
\begin{cases}
\theta_p(p)=p,\\
\theta_\ell(p)=p^{-1},\\
\theta_q(p)=1,\\
\theta_\infty(p)=1.
\end{cases}
\]
事实上，回忆
\[
\theta_p:\Q_p^\times\to\Gal(\Q_p(\zeta_{\ell^r})/\Q_p),\qquad p\mapsto\operatorname{Frob}_p,
\]
且 $\theta_p(\Z_p^\times)=1$。
% [原文作 "θ(Z_p^×) = 1"，疑脱下标 p，见报告 §3]
但 $\operatorname{Frob}_p$ 对应 $p\in(\Z/\ell^r\Z)^\times$。对第二种情形，由 Dwork 定理 $\theta_\ell(p)$ 以 $p^{-1}\in(\Z/\ell^r\Z)^\times$ 作用；对第三种情形，其中 $q\ne p,\ell$，有 $p\in(\Z/q\Z)^\times$ 且扩张 $\Q_q(\zeta_{\ell^r})/\Q_q$ 非分歧。最后，$\theta_\infty$ 对应取符号，因为它是由收缩 $\R^\times$ 的连通分支得到的映射 $\R^\times\to\Gal(\C/\R)$，给出 $\Z/2\Z$ 的、把 $1$ 送到复共轭的映射。\hfill$\square$

\addcontentsline{toc}{section}{第 22 讲　阿廷与布劳尔互反律（下）}
% 第 22 讲 阿廷与布劳尔互反律（下）（源：src/lecture22.txt；已删去抽出时残留的页眉「22. ARTIN AND BRAUER RECIPROCITY, PART II 97/99」）
\section*{第 22 讲\quad 阿廷与布劳尔互反律（下）}

本讲的目标是在模去第二不等式的条件下证明阿廷互反律与布劳尔互反律。回忆阿廷互反律的陈述：对数域 $F$ 与素位 $v$，lcft 给出映射
\[
\theta_v:F_v^\times\to\Gal^{\mathrm{ab}}(F),
\]
而我们断言诱导映射
\[
\theta:A_F^\times\to\Gal^{\mathrm{ab}}(F)
\]
限制到 $F^\times$ 上平凡。这是二次互反律的巨大推广：若 $F=\Q$，则对任意二次扩张 $\Q(\sqrt{\ell})$，有 $\theta_v(\cdot)=(\cdot,\ell)_v$，即 $F_v$ 上的希尔伯特符号，其中 $v$ 是素数或 $\infty$。
% [原文作 "for any quadratic extension Q(`)"，抽出时丢了根号；按上下文还原为 Q(√ℓ)，见报告 §3]
然而，这里我们断言这对整个交换化的 Galois 群成立，而不只是对这个同构于 $\Z/2\Z$ 的特殊商成立。

\textbf{证明（阿廷互反律）.} 首先注意，对任意阿贝尔扩张 $E/F$，我们得到映射 $\theta:A_F^\times\to\Gal(E/F)$；显然只需在每个这样的商中证明消没。

\textbf{情形 1.} 设 $F:=\Q$，$E:=\Q(\zeta_{\ell^r})$，其中 $\ell$ 是素数。必须证明对每个素数 $p$ 有 $\theta(p)=1$，且 $\theta(-1)=1$，因为这些元素生成 $\Q^\times$ 而 $\theta$ 是群同态。我们借助 Dwork 定理作显式计算。先设 $p\ne\ell$。则
\[
\begin{cases}
\theta_q(p)=1,\\
\theta_\infty(p)=1\\
\theta_p(p)=p,\\
\theta_\ell(p)=p^{-1},
\end{cases}
\qquad\text{其中 } q\ne p,\ell.
\]
对第一个等式，注意只要 $q\ne\ell$，分圆扩张 $E/F$ 在 $q$ 处非分歧，而局部阿廷映射在 $q$ 处非分歧的扩张上杀掉 $(\Z/q\Z)^\times$ 的每个元素（见断言 21.5 的证明）。对第二个等式，注意每个 $\theta_v$ 都经过 $v$ 处的分解群分解，而此情形下分解群是 $\Z/2\Z$，且不难显式看出 $\theta_\infty:\R^\times\to\Z/2\Z$ 对应符号函数。对第三个等式，注意 $E/F$ 在 $p$ 处非分歧，因为 $p\ne\ell$。由 Dwork 定理，素元 $p$ 映到 $\Gal(E/F)=(\Z/\ell^r\Z)^\times$ 中的 Frobenius 元素，而它就是 $p$。最后，$E/F$ 在 $\ell$ 处完全分歧，且 $p\in(\Z/\ell\Z)^\times$，故 Dwork 定理给出 $\theta_\ell(p)=p^{-1}$。

现在设 $p=\ell$。则
\[
\begin{cases}
\theta_q(\ell)=1,\\
\theta_\infty(\ell)=1\\
\theta_\ell(\ell)=1,
\end{cases}
\qquad\text{其中 } q\ne\ell.
\]
最后一个等式由 Dwork 定理得到，因为 $\ell$ 在扩张的完全分歧因子上平凡作用（而在非分歧因子上作为其 Frobenius 元素作用）。

最后检验 $\theta(-1)$：
\[
\begin{cases}
\theta_q(-1)=1,\\
\theta_\infty(-1)=-1\\
\theta_\ell(-1)=(-1)^{-1},
\end{cases}
\qquad\text{其中 } q\ne\infty,\ell.
\]
由于 $-1\in(\Z/\ell\Z)^\times$，最后一种情形可如前应用 Dwork 定理。

\textbf{情形 2.} 仍设 $F:=\Q$，$E:=\Q(\zeta_n)$，$n$ 为任意整数。则
\[
n=\prod_{i=1}^m p_i^{r_i}
\]
其中 $p_i$ 为素数，因此
\[
\Q(\zeta_n)=\Q(\zeta_{p_1^{r_1}})\cdots\Q(\zeta_{p_m^{r_m}})
\]
是其各素数幂因子对应的复合，从而
\[
\Gal(\Q(\zeta_n)/\Q)=\prod_{i=1}^m\Gal(\Q(\zeta_{p_i^{r_i}})/\Q).
\]
于是阿廷映射
\[
\theta:A_\Q^\times\to\Gal(\Q(\zeta_n)/\Q)
\]
由 lcft 是各 $\Q(\zeta_{p_i^{r_i}})/\Q$ 的阿廷映射之积给出的，故只需注意下述一般性断言：线性不交扩张的阿廷互反律蕴含它们复合的阿廷互反律。也就是说，$\Q^\times$ 被杀死，因为它的每个坐标都被前一情形杀死。

\textbf{情形 3.} 设 $F$ 是一般数域，$E:=F(\zeta_n)$，$n$ 为某整数。由 lcft（在局部域乘法群的层面上），我们有下述交换图：
\[
\begin{CD}
F^\times @>>> A_F^\times @>{\theta}>> \Gal(F(\zeta_n)/F)\\
@V{N}VV @V{N}VV @VVV\\
\Q^\times @>>> A_\Q^\times @>{\theta}>> \Gal(\Q(\zeta_n)/\Q),
\end{CD}
\]
% [原文作 "at the level of multiplicative groups os local fields"，"os" 疑为 "of" 之误，见报告 §3]
其中 $\theta$ 表示阿廷映射。由于最右的映射是单射，只需证明 $N(F^\times)\subseteq\Q^\times$ 在 $\Gal(\Q(\zeta_n)/\Q)$ 中消没，而这恰是前一情形。

\textbf{情形 4.} 设 $E/F$ 是数域的分圆扩张，即对某个 $n$ 有 $E\subseteq F(\zeta_n)$。则由 lcft 有交换图
\[
\begin{array}{ccc}
A_F & \xrightarrow{\ \theta\ } & \Gal(F(\zeta_n)/F)\\[5pt]
& \stackrel{\theta}{\searrow} & \big\downarrow\\[5pt]
& & \Gal(E/F),
\end{array}
\]
而既然 $\theta$ 在上方的 Galois 群中杀掉 $F^\times$，它在下方那个中也如此。

我们已经证明阿廷互反律对所有数域的分圆扩张成立；对一般情形，我们将使用布劳尔互反律，现在转向它。\hfill$\square$

作为附注，Sam 大概是世上唯一称之为「布劳尔互反律」的人；通常它被称为「从 gcft 计算布劳尔群」之类。

设 $E/F$ 是整体域的 $G$-Galois 扩张。我们有 $G$-模的短正合列
\[
1\to E^\times\to A_E^\times\to C_E\to 1,
\]
它给出复合
\[
\begin{aligned}
\underbrace{\Br(F/E)}_{H^2(G,E^\times)}&\longrightarrow\underbrace{\bigoplus_v\Br(F_v/E_w)}_{H^2(G,A_E^\times)}\hookrightarrow\bigoplus_v\Br(F_v)\\
&\xrightarrow{\ (x_v)_{v\in M_F}\mapsto\sum_{v\in M_F\setminus M_F^\infty}x_v\ }\Q/\Z,
\end{aligned}
\]
其中回忆 $\Br(F/E)$ 表示 $F$ 上那些与 $E$ 张量后成为矩阵代数的可除代数构成的群；$w$ 是 $v$ 上方某个选定的素位；再回忆
\[
\Br(F_v)=
\begin{cases}
\Q/\Z & \text{若 } v \text{ 有限},\\
\tfrac12\Z/\Z & \text{若 } v \text{ 实},\\
0 & \text{若 } v \text{ 复};
\end{cases}
\]
而最右的映射被称为「不变量」映射，因为它是「局部不变量」之和，这些局部不变量分类局部可除代数。此复合对应于把中心单代数与各局部素位上的可除代数作张量、在每一这样的素位取不变量、再把它们相加。直和（自动地）告诉我们：除有限多个素位外，我们在域上得到的是矩阵代数。布劳尔互反律断言此复合为零。于是例如，不存在 $\Q$（或任意数域）上的中心单代数，使得它在除一个素位以外的所有素位上都是（分裂的）矩阵代数。注意当此复合作用于（由两个有理数给出的）哈密顿代数时，它不过是记录了希尔伯特互反律。

\textbf{断言 22.1.} 设 $E/F$ 是整体域的循环扩张。则阿廷互反律等价于布劳尔互反律。

\textbf{证明.} 取定一个生成元，有 $\Gal(E/F)=G\cong\Z/n\Z$。由于 $G$ 循环，它等于自身的交换化，且 Tate 上同调是 $2$-周期的，于是有下述交换图：
\[
\begin{CD}
\widehat H^0(G,E^\times)=F^\times/N(E^\times) @>>> \widehat H^0(G,A_E^\times)=A_F^\times/N(A_E^\times) @>{\theta}>> G=\tfrac1n\Z/\Z\\
@| @| @VVV\\
\Br(F/E)=\widehat H^2(G,E^\times) @>>> \widehat H^2(G,A_E^\times) @>{\iota}>> \Q/\Z,
\end{CD}
\]
其中 $\iota$ 表示不变量映射。现在，左边的方块平凡地交换，而右边的方块由循环扩张的 lcft 交换。断言随即由简单的图追踪得到。\hfill$\square$

\textbf{断言 22.2.} 对任意整体域 $F$ 与每个 $\beta\in\Br(F)$，存在循环分圆扩张 $E/F$ 使 $\beta\in\Br(F/E)$，即 $\beta$ 也落在相对布劳尔群中。

暂且假设此断言。则我们可以一般地推出布劳尔互反律：由于扩张 $E/F$ 循环，我们知道阿廷互反律与布劳尔互反律等价，而且我们还知道阿廷互反律成立，因为该扩张是分圆的。
% [原文作 "we know that Brauer reciprocity is equivalent to Brauer reciprocity"，前一处的 Brauer 应为 Artin，见报告 §3]

此断言对局部域容易验证：给定局部域 $K$ 上的可除代数，若其次数的平方根整除 $[L:K]$，则它在域扩张 $L/K$ 上分裂。事实上，回忆若 $\beta\subseteq\frac1n\Z/\Z\subseteq\Br(K)$，即 $\beta$ 是 $n^2$ 次可除代数，则 $\beta\in\Br(K/L)$ 当且仅当 $n\mid[L:K]$，因为
\[
\Q/\Z=\Br(K)\xrightarrow{\ \times[L:K]\ }\Br(L)=\Q/\Z.
\]

现在回忆下述定理，我们将在下一讲证明：

\textbf{定理 22.3.} 对整体域的一切扩张 $E/F$，有 $H^1(G,C_E)=0$。

这算是希尔伯特定理 90 的类比，而证明它是 gcft 中最困难的部分；我们暂且假设它。短正合列
\[
1\to E^\times\to A_E^\times\to C_E\to 1
\]
于是给出正合列
\[
\underbrace{H^1(A_E^\times)}_{0}\to\underbrace{H^1(C_E)}_{0}\to H^2(E)=\Br(F/E)\to H^2(A_E^\times)=\bigoplus_v\Br(F_v/E_w),
\]
% [原文此二项下有 underbrace，抽出时标注丢失；按紧随其后的「the vanishing is by the previous theorem and Hilbert's Theorem 90」补作 0，见报告 §2]
其中消没由前一定理与希尔伯特定理 90 给出（因为阿代尔的上同调就是各局部上同调的直和）。取正向极限，得到下述推论：

\textbf{推论 22.4.} 对任意整体域 $F$，有典范的单射
\[
\Br(F)\hookrightarrow\bigoplus_{v\in M_F}\Br(F_v).
\]

换言之，数域 $F$ 上的中心单代数是矩阵代数，当且仅当它在 $F$ 的每个完备化处都是矩阵代数。这绝不是显然的陈述！

要点是：为证明断言 22.2，只需证明下述断言：

\textbf{断言 22.5.} 给定整体域 $F$ 的有限素位集 $S$，以及对所有 $v\in S$ 的正整数 $m_v$（要求：$v$ 复时 $m_v=1$；$v$ 实时 $m_v=1,2$；$v$ 有限时 $m_v$ 任意），则存在循环分圆扩张 $E/F$，使得 $E_w/F_v$ 的次数被 $m_v$ 整除（对任意选定的 $w\mid v$）。
% [原文作 "such that m_v = 1 is v is complex"，第二个 "is" 疑为 "if" 之误，见报告 §3]

\textbf{证明（断言 22.5 $\Rightarrow$ 断言 22.2）.} 由推论 22.4，$F$ 上的中心单代数分裂，当且仅当它在每个局部域 $F_v$ 上分裂，而这又当且仅当其次数的平方根整除某个扩张 $E/F$ 与素位 $w\mid v$ 的 $[E_w:F_v]$。于是，对每个素位 $v$ 取 $m_v$ 为其在 $F_v$ 上次数的平方根（等价地，取其在 $v$ 处局部不变量的分母），当 $v$ 实时它必为 $1$ 或 $2$、$v$ 复时为 $1$，此断言蕴含我们的中心单代数在它所提供的扩张 $E$ 上分裂。此外，如前所述，这样的中心单代数在除有限多个素位 $v$ 外的所有 $v$ 上都在 $F_v$ 上分裂，这意味着我们可以取 $S$ 为仅由那些它不分裂的素位组成的集合。\hfill$\square$

在开始证明之前，让我们花点时间注意：确实存在不是循环的分圆扩张！事实上，$\Gal(\Q(\zeta_n)/\Q)=(\Z/n\Z)^\times$，其中 $\zeta_n$ 是本原 $n$ 次单位根，因此例如当 $n=15$ 时，
\[
(\Z/15\Z)^\times\cong(\Z/3\Z)^\times\times(\Z/5\Z)^\times\cong\Z/2\Z\times\Z/4\Z
\]
不是循环的！

\textbf{证明（断言 22.5）.} 首先注意可设 $F=\Q$：我们可把每个局部因子 $m_v$ 换为 $[F_v:\Q_p]\cdot m_v$（其中 $p$ 是 $v$ 下方的 $\Q$ 的素位），在 $\Q$ 上应用本断言，然后把 $E$ 换为复合 $E\cdot F$，它也是 $F$ 的循环分圆扩张（因为循环群的子群循环）。

\textbf{情形 1.} 设对每个局部素位 $v\in S$，有 $m_v=p^{r_v}$，其中 $p$ 是（与 $v$ 无关的）奇素数。对任意 $r$，有分圆扩张 $\Q(\zeta_{p^r})/\Q$，Galois 群为 $(\Z/p^r\Z)^\times$。注意 $\#(\Z/p^r\Z)^\times=(p-1)p^{r-1}$，故可令 $E_r$ 为 $\Q(\zeta_{p^r})/\Q$ 的循环子扩张，$\Gal(E_r/\Q)\cong\Z/p^{r-1}\Z$。则对每个 $v\in S$ 与任意选定的 $w\mid v$，我们断言 $[E_{r,w}:\Q_v]\to\infty$（$r\to\infty$）。事实上，除 $\C$ 外的任何局部域只含有限多个单位根，故这些扩张的次数必递增。现在，
\[
[E_{r,w}:\Q_v]=[\Q_v(\zeta_{p^r}):\Q_v]\mid(p-1)p^\infty,
\]
% [原文抽出作 "[Q_v(ζ_{p^r}/Q_v] | (p−1)p^∞"，缺右方括号与冒号；按图像排作 [Q_v(ζ_{p^r}):Q_v]，见报告 §3]
故此 $p$ 幂因子发散；由于 $S$ 有限，本情形得证。请注意这些扩张是全复的（即每个无限素位都是复的），故不必担心那些 $m_v=2$ 的实素位 $v$。

\textbf{情形 2.} 现设对每个 $v\in S$ 有 $m_v=2^{r_v}$。此情形类似，只差一点：$(\Z/2^r\Z)^\times\cong\Z/2^{r-2}\Z\times\Z/2\Z$，此同构可通过验证下述容易的恒等式得到：
\[
\Z/2^{r-2}\Z\cong\{x\in(\Z/2^r\Z)^\times:x\equiv1\bmod4\}
\]
（于是可通过计算其 $2$-挠部分、再验证它只含 $2$ 个元素等等，来证明它是循环的），并注意 $\Z/2\Z\cong\{1,2^{r-1}-1\}$。于是可令 $\Q(\zeta_{2^r})/E_r/\Q$ 为分圆扩张，其次数为 $2^{r-2}$，Galois 群为 $(\Z/2^{r-2}\Z)^\times$，实现为 $(\Z/2^r\Z)^\times$ 的商。
% [原文作 "a cyclotomic degree-2^{r-2} extension whose Galois group is (Z/2^{r-2}Z)^×"，(Z/2^{r-2}Z)^× 的阶为 2^{r-3}，与次数 2^{r-2} 不符；疑笔误，照译并见报告 §3]

我们断言 $E_r/\Q$ 是全复扩张。由于它是阿贝尔扩张，只需证明复共轭（对应 $-1\in(\Z/2^r\Z)^\times$）在其上非平凡地作用。对充分大的 $r$，$-1\not\equiv1\bmod4$，故由上述显式同构，它在 $(\Z/2^{r-2}\Z)^\times$ 中的投影为 $1-2^{r-1}\ne1$（注意 $2^{r-1}-1\equiv3\bmod4$ 是任一因子中唯一具有此性质的非平凡元素）。于是，由与前一情形类似的论证，扩张 $E_r/\Q$ 即已足够。

\textbf{情形 3.} 最后，对一般情形，取前两种情形所构造的、对应于 $m_v$ 的各素因子的扩张的复合。此外不会发生任何「干扰」（以致复合不再循环），因为每个扩张的 Galois 群都有互不相同的素数幂阶。\hfill$\square$

于是，假设我们关于希尔伯特定理 90 的类比成立，整体域 $F$ 的布劳尔群的每个元素都被某个循环分圆扩张分裂，这蕴含对 $\Br(F)$ 所有元素的布劳尔互反律。这又蕴含循环扩张的阿廷互反律，因为对循环扩张这两条互反律等价。最后，由于任意阿贝尔群都是循环群之积，每个阿贝尔扩张 $E/F$ 都是循环扩张的复合，从而一般地蕴含阿廷互反律（我们已经看到，一组线性不交扩张的阿廷互反律蕴含它们复合的阿廷互反律）。

\addcontentsline{toc}{section}{第 23 讲　第二不等式的证明}
% 第 23 讲 第二不等式的证明（源：src/lecture23.txt；已删去抽出时残留的页眉「23. PROOF OF THE SECOND INEQUALITY 103/105/107」）
\section*{第 23 讲\quad 第二不等式的证明}

本讲的目标是证明「第二不等式」：对整体域的一切扩张 $E/F$，有 $H^1(G,C_E)=0$，其中 $C_E:=A_E^\times/E^\times$ 是 $E$ 的「伊代尔类群」。我们的主要情形是 $E/F$ 为 $p$ 阶循环扩张，且对某个本原 $p$ 次单位根 $\zeta_p$ 有 $\zeta_p\in F$；我们将在本讲末尾归约到此情形（注意 $\operatorname{char}(F)=0$，因为我们假设 $F$ 是数域）。在此情形下，只需证明
\[
\#\widehat H^0(C_E)=\#(C_F/NC_E)=\#\bigl(A_F^\times/F^\times\cdot N(A_E^\times)\bigr)\le p.
\]
事实上，由「第一不等式」我们知道
\[
\frac{\#\widehat H^0(C_E)}{\#\widehat H^1(C_E)}=p,
\]
于是 $p\cdot\#\widehat H^1(C_E)=\#\widehat H^0(C_E)\le p$ 蕴含 $\#H^1(C_E)=1$，正如所愿。我们的方法将是「试错」式的——即先尝试某种做法，它还不够好，然后再加以修正。

取定 $F$ 的一个有限素位集 $S$，使得
\begin{enumerate}[label=(\arabic*),leftmargin=2.4em]
\item 若 $v\mid\infty$，则 $v\in S$；
\item 若 $v\mid p$，则 $v\in S$；
\item $A_F^\times=F^\times\cdot A_{F,S}^\times$，其中回忆
\[
A_{F,S}^\times:=\prod_{v\in S}F_v^\times\times\prod_{v\notin S}\mathcal{O}_{F_v}^\times,
\]
且由引理 20.12 这是可行的；
\item $E=F(\sqrt[p]{u})$，其中 $u\in\mathcal{O}_{F,S}^\times:=F^\times\cap A_{F,S}^\times$ 是 $F$ 的「$S$-单位」。
\end{enumerate}
最后一条由 Kummer 理论可行。

注意最后这个条件蕴含 $E$ 在 $S$ 外非分歧：因为 $u$ 在任何 $v\notin S$ 处是整元素，且由假设所有整除 $p$ 的素位都在 $S$ 中，$p$ 与 $F_v$ 的剩余域的阶互素，故 $F_v(\sqrt[p]{u})/F_v$ 是非分歧扩张。

一个重要的断言（稍后将以略为精细的形式证明）如下：

\textbf{断言 23.1.} $u\in\mathcal{O}_{F,S}^\times$ 是 $p$ 次幂，当且仅当它在 $F_v^\times$ 中的像对每个 $v\in S$ 都是 $p$ 次幂。

令
\[
\Gamma:=\prod_{v\in S}(F_v^\times)^p\times\prod_{v\notin S}\mathcal{O}_{F_v}^\times\subseteq A_{F,S}^\times.
\]
则我们有下述断言：

\textbf{断言 23.2.} $\mathcal{O}_{F,S}^\times\cap\Gamma=(\mathcal{O}_{F,S}^\times)^p$。

\textbf{证明.} 这由前一断言平凡地推出。\hfill$\square$

\textbf{断言 23.3.} $\Gamma\subseteq N(A_E^\times)$。

\textbf{证明.} 扩张 $E/F$ 在每个 $v\notin S$ 处非分歧，故因子 $\prod_{v\notin S}\mathcal{O}_{F_v}^\times\subseteq N(A_E^\times)$。由于 $p$ 杀掉了 $\widehat H^0(E_w^\times)$（对某个选定的 $w\mid v$），可知因子 $\prod_{v\in S}(F_v^\times)^p\subseteq N(A_E^\times)$ 亦然。\hfill$\square$

于是
\[
\#\bigl(A_F^\times/F^\times\cdot N(A_E^\times)\bigr)\le\#\bigl(A_F^\times/F^\times\cdot\Gamma\bigr),
\]
且有短正合列
\[
1\to\mathcal{O}_{F,S}^\times/(\mathcal{O}_{F,S}^\times\cap\Gamma)\to A_{F,S}^\times/\Gamma\to A_F^\times/(F^\times\cdot\Gamma)\to 1.
\]
事实上，第三个映射由上面 $S$ 的性质 (3) 满射，第二个映射单射因为 $\mathcal{O}_{F,S}^\times\subseteq F^\times$，而在 $A_{F,S}^\times/\Gamma$ 处的正合性按定义成立。于是
\[
\#\bigl(A_F^\times/F^\times\cdot\Gamma\bigr)=\frac{\#(A_{F,S}^\times/\Gamma)}{\#(\mathcal{O}_{F,S}^\times/\mathcal{O}_{F,S}^\times\cap\Gamma)},
\]
剩下的就是计算此式的分子与分母。我们有
\[
A_{F,S}^\times/\Gamma=\prod_{v\in S}F_v^\times/(F_v^\times)^p,
\]
% [原文此处作 "∏_{p∈S}"，而通项用 v；按上下文作 v∈S，见报告 §3]
并由 (6.3) 回忆
\[
\#\bigl(F_v^\times/(F_v^\times)^p\bigr)=p\cdot\frac{\#\mu_p(F_v)}{|p|_v}=\frac{p^2}{|p|_v},
\]
因为由假设 $\zeta_p\in F$。于是
\[
\prod_{v\in S}\frac{p^2}{|p|_v}=p^{2\#S}
\]
由乘积公式得到，因为由假设 $v\notin S$ 时 $|p|_v=1$。现在想计算
\[
\#\bigl(\mathcal{O}_{F,S}^\times/\mathcal{O}_{F,S}^\times\cap\Gamma\bigr)=\#\bigl(\mathcal{O}_{F,S}^\times/(\mathcal{O}_{F,S}^\times)^p\bigr).
\]
回忆由 $S$-单位定理，
\[
\mathcal{O}_{F,S}^\times\cong\Z^{\#S-1}\times(\mathcal{O}_{F,S}^\times)_{\mathrm{tors}}.
\]
后者循环，且其阶被 $p$ 整除，故
\[
\#\bigl(\mathcal{O}_{F,S}^\times/(\mathcal{O}_{F,S}^\times)^p\bigr)=p^{\#S-1}\cdot p=p^{\#S}.
\]
合并这两个结果，得到
\[
\#\widehat H^0(C_E)\le\frac{p^{2\#S}}{p^{\#S}}=p^{\#S},
\]
遗憾的是这还不够好。

下面是我们改进此结果的办法：

\textbf{断言 23.4.} 给定这样的集合 $S\subseteq M_F$，存在集合 $T\subseteq M_F$，使得
\begin{enumerate}[label=(\arabic*),leftmargin=2.4em]
\item $\#T=\#S-1$；
\item $S\cap T=\varnothing$；
\item 每个 $v\in T$ 在 $E$ 中分裂，即对一切 $w\mid v$ 有 $E_w=F_v$；
\item 任意 $u\in\mathcal{O}_{F,S\cup T}^\times$ 是 $p$ 次幂，当且仅当 $u\in F_v^\times$ 对一切 $v\in S$ 都是 $p$ 次幂。
\end{enumerate}

注意这里与先前情形的关键差别：在性质 (4) 中，我们不要求 $u\in(F_v^\times)^p$ 对一切 $v\in S\cup T$ 成立，只要求对一切 $v\in S$ 成立。给定这样的 $T$，我们重新定义 $\Gamma$ 为
\[
\Gamma:=\prod_{v\in S}(F_v^\times)^p\times\prod_{v\in T}F_v^\times\times\prod_{v\notin S\cup T}\mathcal{O}_{F_v}^\times.
\]

\textbf{断言 23.5.} $\Gamma\subseteq N(A_E^\times)$。

\textbf{证明.} 性质 (3) 给出关于第二个因子的断言；第一个与第三个因子如前。\hfill$\square$

用以 $A_{F,S\cup T}^\times$ 代替 $A_{F,S}^\times$ 的方式重做我们的计算，由性质 (4) 如前得到
\[
\#(A_{F,S\cup T}^\times/\Gamma)=p^{2\#S},
\]
并且同样如前，
\[
\#\bigl(\mathcal{O}_{F,S\cup T}^\times/(\mathcal{O}_{F,S\cup T}^\times\cap\Gamma)\bigr)=p^{\#(S\cup T)}=p^{2\#S-1},
\]
故二者之商为 $p$，正如所愿！于是只需证明上述断言。

\textbf{断言 23.6.} 对整体域的任意阿贝尔扩张 $F'/F$，$v\notin S$ 的 Frobenius 元素生成 $\Gal(F'/F)$。

我们希望纯粹代数地证明这一点，而不用切博塔廖夫密度定理（后者无论如何给出的是略为不同的陈述）。

\textbf{证明.} 设 $H$ 是由一切 $v\notin S$ 的 Frobenius 生成的子群，$F'':=(F')^H$ 为不动域。我们希望证明 $F''=F$。注意对一切 $v\notin S$，$\operatorname{Frob}_v$ 在 $\Gal(F'/F)/H=\Gal(F''/F)$ 中平凡，故每个 $v\notin S$ 都在 $F''/F$ 中分裂（因为由假设它们非分歧）。于是对一切 $w\mid v$ 与 $v\notin S$ 有 $F''_w=F_v$，而我们断言这不可能。

可设 $F''/F$ 是 $n$ 次循环扩张（必要时以更小的扩张代替之）。由第一不等式，$\chi(C_{F''})=n$，这给出
\[
\#\bigl(A_F^\times/N(A_{F''}^\times)\cdot F^\times\bigr)=\#\widehat H^0(C_{F''})\ge n.
\]
但因此扩张对一切 $v\notin S$ 分裂，平凡地有 $N((F''_v)^\times)=F_v^\times$，从而 $\prod\coprod_{v\notin S}F_v^\times\subseteq N(A_{F''}^\times)$，其中这是限制直积。强逼近随后给出 $F^\times\cdot\prod\coprod_{v\notin S}F_v^\times$ 在 $A_F^\times$ 中稠密，而它又是开的，故除非 $n=1$ 便得到矛盾，正如所愿。\hfill$\square$

我们希望对此断言应用于 $F':=F(\{\sqrt[p]{u}:u\in\mathcal{O}_{F,S}^\times\})$。先有一个断言：

\textbf{断言 23.7.} 对上述 $F'$，有 $\Gal(F'/F)=(\Z/p\Z)^{\#S}$。

\textbf{证明.} 这实质上是 Kummer 理论，因为 $\mathcal{O}_{F,S}^\times/(\mathcal{O}_{F,S}^\times)^p\subseteq F^\times/(F^\times)^p$。我们知道 $F$ 的一切指数 $p$ 的扩张都通过添加 $F^\times$ 中元素的 $p$ 次根得到。其 Galois 群必是 $\Z/p\Z$ 的若干个拷贝之积，但其中某些子群可能重合——迭代应用 Kummer 理论便给出该陈述。\hfill$\square$

现在我们有 $F'/E/F$，因为假设 $E/F$ 是通过添加某个 $S$-单位的 $p$ 次根得到的。选取 $E$ 的素位 $w_1,\dots,w_{\#S-1}$，它们不整除 $S$ 中任何素位，且其 Frobenius 给出 $\Gal(F'/E)\cong(\Z/p\Z)^{\#S-1}$ 的一组基；由断言 23.6 的论证这是可能的。然后令 $T:=\{v_1,\dots,v_{\#S-1}\}$ 为各 $w_i$ 到 $F$ 的限制。

\textbf{断言 23.8.} 每个 $v\in T$ 在 $E$ 中分裂。

\textbf{证明.} 由于 $\operatorname{Frob}_v\in\Gal(F'/E)$，它在 $E$ 上平凡作用，故对任意 $w\mid v$，$\Gal(E_w/F_v)$ 平凡，正如所愿。\hfill$\square$

这建立了 $T$ 的条件 (3)；还需证明条件 (4)，因为条件 (1) 与 (2) 在 $T$ 的构造中已隐含。

\textbf{断言 23.9.} 元素 $x\in\mathcal{O}_{F,S\cup T}^\times$ 是 $p$ 次幂，当且仅当对每个 $v\in S$ 有 $x\in(F_v^\times)^p$。

\textbf{证明.} \emph{第一步.} 我们断言
\[
\mathcal{O}_{F,S}^\times\cap(E^\times)^p=\{x\in\mathcal{O}_{F,S}^\times:x\in(F_v^\times)^p\ \text{对一切 } v\in T\}.
\]
正向包含平凡，因为由前一断言 $(F_v^\times)^p=(E_w^\times)^p$。对反向，注意对任意 $x\in\mathcal{O}_{F,S}^\times$，有扩张 $F'/E(\sqrt[p]{x})/E$。若对每个 $w\mid v$ 与 $v\in T$ 有 $x\in(E_w^\times)^p$，则此扩张在 $w$ 处分裂，故 $\operatorname{Frob}_w$ 在 $E(\sqrt[p]{x})$ 上平凡作用；于是 $\Gal(F'/E)$ 在 $E(\sqrt[p]{x})$ 上平凡作用，因为它由这些 Frobenius 生成；故 $E(\sqrt[p]{x})=E$ 且 $x\in(E^\times)^p$，正如所愿。

\emph{第二步.} 现在我们断言典范映射
\[
\mathcal{O}_{F,S}^\times\xrightarrow{\ \varphi\ }\prod_{v\in T}\mathcal{O}_{F_v}^\times/(\mathcal{O}_{F_v}^\times)^p
\]
满射。正是这一步真正建立了对 $T$ 的大小的限制，第二不等式由此完美地推出。我们将通过计算两端的阶来进行。由第一步，
\[
\operatorname{Ker}(\varphi)=\{x\in\mathcal{O}_{F,S}^\times:x\in(E^\times)^p\}.
\]
则 $\mathcal{O}_{F,S}^\times/\operatorname{Ker}(\varphi)$ 的阶为 $p^{\#S-1}$。事实上，我们先前算出 $\mathcal{O}_{F,S}^\times/(\mathcal{O}_{F,S}^\times)^p$ 的阶为 $p^{\#S}$，且由于
\[
(\mathcal{O}_{F,S}^\times)^p=\{x\in\mathcal{O}_{F,S}^\times:x\in(F^\times)^p\}
\]
而 $E/F$ 是添加某个 $S$-单位的 $p$ 次根得到的 $p$ 次扩张，可知 $[\operatorname{Ker}(\varphi):(\mathcal{O}_{F,S}^\times)^p]=p$。现在，用我们先前的公式关于 $\mathcal{O}_{F_v}^\times$（而非 $F_v^\times$）的版本，右端的阶为
\[
\prod_{v\in T}\frac{\#\mu_p(F_v)}{|p|_v}=p^{\#T}=p^{\#S-1},
\]
故该映射确实满射。

\emph{第三步.} 现在希望建立此断言：若对一切 $v\in S$ 有 $x\in(F_v^\times)^p$，则 $x\in(\mathcal{O}_{F,S\cup T}^\times)^p$（反向平凡）。我们希望证明 $F(\sqrt[p]{x})=F$。令
\[
\Gamma:=\prod_{v\in S}F_v^\times\times\prod_{v\in T}(\mathcal{O}_{F_v}^\times)^p\times\prod_{v\notin S\cup T}\mathcal{O}_{F_v}^\times\subseteq A_{F,S}^\times,
\]
这里又是一个与先前不同的 $\Gamma$。则事实上
\[
\Gamma\subseteq N(A_{F(\sqrt[p]{x})}^\times)\subseteq A_F^\times,
\]
其中第三个因子是因为 $F(\sqrt[p]{x})/F$ 在 $S\cup T$ 外非分歧，第二个因子是因为 $[F(\sqrt[p]{x}):F]\le p$，第一个因子是因为由假设该扩张在 $S$ 的一切素位处分裂。现在我们要证明 $F^\times\cdot\Gamma=A_F^\times$，因为第一不等式随后便如断言 23.6 那样给出结果。由第二步，有
\[
\mathcal{O}_{F,S}^\times\twoheadrightarrow\prod_{v\in T}\mathcal{O}_{F_v}^\times/(\mathcal{O}_{F_v}^\times)^p=A_{F,S}^\times/\Gamma,
\]
故 $\mathcal{O}_{F,S}^\times\cdot\Gamma=A_{F,S}^\times$。这意味着
\[
F^\times\cdot\Gamma=F^\times\cdot A_{F,S}^\times=A_F^\times
\]
由关于 $S$ 的假设。\hfill$\square$

现在希望由已证的特殊情形推出第二不等式的一般情形。第一步如下：

\textbf{断言 23.10.} 若第二不等式对任意基域含 $p$ 次单位根的 $p$ 阶循环扩张成立，则它对任意 $p$ 阶循环扩张成立。

\textbf{证明.} 设 $E/F$ 是整体域的 $p$ 次循环扩张。回忆 $E/F$ 的第二不等式等价于典范单射
\[
\Br(F/E)\hookrightarrow\bigoplus_{v\in M_F}\Br(F_v)
\]
的存在性。事实上，有短正合列
\[
0\to E^\times\to A_E^\times\to C_E\to 0,
\]
而上同调长正合列于是给出
\[
\underbrace{H^1(G,A_E^\times)}_{\oplus H^1(E_w^\times)=0}\to H^1(G,C_E)\to\Br(F/E)\to\bigoplus_v\Br(F_v/E_w)\subseteq\bigoplus_v\Br(F_v)
\]
（对某个选定的 $w\mid v$），其中第一个等号由希尔伯特定理 90 得到。为证明 $H^1(G,C_E)$ 的消没，只需证明最后的映射单射。现在，域扩张
\[
\begin{array}{ccccc}
& & E(\zeta_p) & &\\
& \nearrow & & \nwarrow &\\
E & & & & F(\zeta_p)\\
& \nwarrow & & \nearrow &\\
& & F & &
\end{array}
\]
诱导出交换图
\[
\begin{array}{ccc}
\Br(F/E) & \xrightarrow{\ \alpha\ } & \bigoplus_v\Br(F_v/E_w)\\[4pt]
\big\downarrow{\scriptstyle\gamma} & & \big\downarrow{\scriptstyle\delta}\\[4pt]
\Br(F(\zeta_p)/E(\zeta_p)) & \xrightarrow{\ \beta\ } & \bigoplus_v\Br(F(\zeta_p)_w),\\[4pt]
\big\downarrow & & \\[4pt]
\Br(F/E) & &
\end{array}
\]
% [按原 PDF 第 106 页复原：原图在最左端用花括号把左上角的 $\Br(F/E)$ 与左下角的 $\Br(F/E)$ 括在一起，并标以 $\times[F(\zeta_p):F]$，即左列两个竖直映射（限制与膨胀）的复合等于乘以 $[F(\zeta_p):F]$]
其中最左的两个映射分别是上同调上的限制映射与膨胀映射，其依据是布劳尔群的上同调解释（见习题集 7 第 2 题）。此外，左列复合 $\Br(F/E)\to\Br(F(\zeta_p)/E(\zeta_p))\to\Br(F/E)$ 是乘以 $[F(\zeta_p):F]$，它在 $\Br(F/E)$ 上单射，因为后者是 $p$-挠的（由习题 2(c)），且 $[F(\zeta_p):F]\mid(p-1)$。于是 $\gamma$ 也单射。由于第二不等式对 $E(\zeta_p)/F(\zeta_p)$ 由假设成立，$\beta$ 单射，从而 $\alpha$ 也单射。
% [原文作 "Since the second equality holds for E(ζ_p)/F(ζ_p)"，此处应为 second inequality，见报告 §3]
\hfill$\square$

\textbf{断言 23.11.} 若第二不等式对任意数域的 $p$ 阶循环扩张成立，则它对任意扩张成立。

\textbf{证明.} 我们希望证明 $H^1(G,C_E)=0$。如同有限群的任何 Tate 上同调群，有单射
\[
H^1(G,C_E)\hookrightarrow\bigoplus_p H^1(G_p,C_E),
\]
其中 $G_p$ 是 $G$ 的 $p$-Sylow 子群。于是可设 $G$ 是 $p$-群。由于每个 $p$-群 $G$ 都含有同构于 $\Z/p\Z$ 的正规子群 $H$，可设我们有域扩张 $E_2/E_1/F$，其中 $\Gal(E_2/E_1)\cong H$，$\Gal(E_1/F)\cong G/H$。可设定理对 $H$ 作用在 $E_2$ 上以及 $G/H$ 作用在 $E_1$ 上成立，于是只需重复前一断言证明中关于布劳尔群的单射性的那类论证。

首先我们断言 $C_{E_2}^H=C_{E_1}$。事实上，有短正合列
\[
0\to E_2^\times\to A_{E_2}^\times\to C_{E_2}\to 0,
\]
而上同调长正合列于是给出
\[
0\to\underbrace{H^0(H,E_2^\times)}_{E_1^\times}\to\underbrace{H^0(H,A_{E_2}^\times)}_{A_{E_1}^\times}\to\underbrace{H^0(H,C_{E_2})}_{C_{E_2}^H}\to\underbrace{H^1(H,E_2^\times)}_{0}
\]
由希尔伯特定理 90。注意 $A_{E_2}^{\times,H}=A_{E_1}^\times$，因为取有限群的不变子模与阿代尔定义中的正向极限和直积交换。于是有
\[
\operatorname{hKer}\Bigl(C_{E_2}^{hG}=(C_{E_2}^{hH})^{hG/H}\to(\tau^{\ge2}C_{E_2}^{hH})^{hG/H}\Bigr)\cong(\tau^{\le0}C_{E_2}^{hH})^{hG/H}=(C_{E_1})^{hG/H},
\]
其中第一个等号由习题集 6 第 3 题，映射由截断的定义而来，拟同构是因为 $H^1(H,C_{E_2})$ 由假设消没，而最后这个表示式不过是 $C_{E_2}$ 的朴素 $H$-不变子模，因为截断杀掉了次数大于 $0$ 的一切上同调，故前一断言给出等式。上同调长正合列随后给出
\[
\underbrace{H^1\bigl((C_{E_1})^{hG/H}\bigr)}_{H^1(G/H,C_{E_1})=0}\to\underbrace{H^1\bigl((C_{E_2})^{hG}\bigr)}_{H^1(G,C_{E_2})}\to\underbrace{H^1\bigl((\tau^{\ge2}C_{E_2}^{hH})^{hG/H}\bigr)}_{0},
\]
因为最右的复形处于次数至少 $2$。于是 $H^1(G,C_{E_2})=0$，正如所愿。\hfill$\square$

\section*{参考文献}

原书参考文献条目保留英文原貌；方括号内为原书引用键。

\noindent\textbf{[Boy07]} Dmitriy Boyarchenko, \emph{Topics in Algebra}, 2007.\par\smallskip
\noindent\textbf{[Dwo58]} Bernard Dwork, \emph{Norm residue symbol in local number fields}, Abh. Math. Semin. Univ. Hambg. 22 (1) (December 1958), 180--190.\par\smallskip
\noindent\textbf{[Mil13]} J. S. Milne, \emph{Class Field Theory}, 2013.\par\smallskip
\noindent\textbf{[Ser73]} Jean-Pierre Serre, \emph{A Course in Arithmetic}, Springer Science \& Business Media New York, 1973.\par\smallskip
\noindent\textbf{[Ser79]} Jean-Pierre Serre, \emph{Local Fields}, Springer Science \& Business Media New York, 1979.\par\smallskip
\noindent\textbf{[Wei74]} André Weil, \emph{Basic Number Theory}, Springer-Verlag Berlin Heidelberg, 1974.\par\smallskip

\section*{术语表（冻结版）}
下表为本次试译冻结的术语对照（共 109 条）。术语在译文中统一使用中文列；表格未覆盖的词在正文首次出现处附英文原文。
\begin{longtable}{p{0.46\textwidth}p{0.46\textwidth}}
\toprule 英文 & 中文 \\ \midrule \endhead
class field theory & 类域论 \\
local class field theory & 局部类域论 \\
global class field theory & 整体类域论 \\
cohomological class field theory & 上同调类域论 \\
Hilbert symbol & 希尔伯特符号 \\
norm residue symbol & 范数剩余符号 \\
norm group & 范数群 \\
norm map & 范映射 \\
Hilbert's Theorem 90 & 希尔伯特定理 90 \\
reciprocity law & 互反律 \\
reciprocity map & 互反映射 \\
local reciprocity map & 局部互反映射 \\
global reciprocity map & 整体互反映射 \\
Artin reciprocity & 阿廷互反律 \\
quadratic reciprocity & 二次互反律 \\
Brauer group & 布劳尔群 \\
Brauer–Manin obstruction & 布劳尔–马宁障碍 \\
Hasse principle & Hasse 原理 \\
Hasse–Minkowski theorem & Hasse–Minkowski 定理 \\
Witt cancellation & Witt 消去 \\
quadratic form & 二次型 \\
first inequality & 第一不等式 \\
second inequality & 第二不等式 \\
maximal abelian extension & 极大阿贝尔扩张 \\
abelianization & 交换化 \\
profinite completion & 投射有限完备化 \\
profinite group & 投射有限群 \\
separable closure & 可分闭包 \\
absolute Galois group & 绝对 Galois 群 \\
inertia group & 惯性群 \\
tame inertia group & 驯惯性群 \\
wild inertia group & 野惯性群 \\
decomposition group & 分解群 \\
ramification index & 分歧指数 \\
unramified & 非分歧 \\
tame ramification & 驯分歧 \\
wild ramification & 野分歧 \\
uniformizer & 素元 \\
valuation & 赋值 \\
ring of integers & 整数环 \\
residue field & 剩余域 \\
local field & 局部域 \\
global field & 整体域 \\
number field & 数域 \\
adèle & 阿代尔 \\
idèle & 伊代尔 \\
idèle class group & 伊代尔类群 \\
adèle ring & 阿代尔环 \\
adèle pairing & 阿代尔配对 \\
Kummer theory & Kummer 理论 \\
Kummer extension & Kummer 扩张 \\
cyclotomic field & 分圆域 \\
cyclotomic character & 分圆特征标 \\
primitive root of unity & 本原单位根 \\
Galois cohomology & Galois 上同调 \\
group cohomology & 群上同调 \\
Tate cohomology & Tate 上同调 \\
cochain & 上链 \\
cochain complex & 上链复形 \\
chain complex & 链复形 \\
cohomology group & 上同调群 \\
homology & 同调 \\
exact sequence & 正合列 \\
short exact sequence & 短正合列 \\
inflation & 膨胀 \\
restriction & 限制 \\
corestriction & 余限制 \\
cup product & 杯积 \\
projective resolution & 投射消解 \\
injective resolution & 内射消解 \\
derived functor & 导出函子 \\
quasi-isomorphism & 拟同构 \\
homotopy & 同伦 \\
mapping complex & 映射复形 \\
coinvariant & 余不变商 \\
invariant & 不变子模 \\
Herbrand quotient & Herbrand 商 \\
Euler characteristic & 欧拉示性数 \\
vanishing theorem & 消没定理 \\
fundamental class & 基本类 \\
class formation & 类构造 \\
Lyndon–Hochschild–Serre spectral sequence & Lyndon–Hochschild–Serre 谱序列 \\
spectral sequence & 谱序列 \\
direct limit & 正向极限 \\
inverse limit & 逆向极限 \\
topological group & 拓扑群 \\
open subgroup & 开子群 \\
norm & 范数 \\
trace & 迹 \\
Galois extension & Galois 扩张 \\
separable & 可分的 \\
algebraic closure & 代数闭包 \\
p-adic & p 进 \\
archimedean & 阿基米德的 \\
nonarchimedean & 非阿基米德的 \\
place & 素位 \\
completion & 完备化 \\
splitting field & 分裂域 \\
solvable & 可解的 \\
radical & 根式 \\
Kronecker–Weber theorem & 克罗内克–韦伯定理 \\
Diophantine equation & 丢番图方程 \\
cocycle & 上闭链 \\
coboundary & 上边缘 \\
Claim & 断言 \\
non-degenerate & 非退化 \\
tame symbol & 驯符号 \\
boundary map & 边界映射 \\
filtered complex & 滤复形 \\
\bottomrule
\end{longtable}



\end{document}
